\documentclass[11pt]{article}
\usepackage[letterpaper,margin=0.78in]{geometry}
\usepackage[T1]{fontenc}
\usepackage{lmodern,amsmath,amssymb,graphicx,booktabs,array,xcolor,fancyhdr,longtable,calc}
\usepackage[hidelinks]{hyperref}
\usepackage{xurl}
\hypersetup{pdftitle={Function Families and Their Structure},pdfauthor={Komorebi P.},pdfsubject={Realization Calculus: special functions, growth and finite resolution}}
\setcounter{tocdepth}{1}
\providecommand{\tightlist}{\setlength{\itemsep}{0pt}\setlength{\parskip}{0pt}}
\definecolor{ink}{HTML}{123B50}
\definecolor{accent}{HTML}{087F8C}
\definecolor{warm}{HTML}{B74D36}
\setlength{\parindent}{0pt}
\setlength{\parskip}{6pt}
\setlength{\emergencystretch}{2em}
\pagestyle{fancy}\fancyhf{}
\setlength{\headheight}{14pt}
\newcommand{\pageheading}[2]{\clearpage\phantomsection\addcontentsline{toc}{section}{#2}{\color{accent}\small\textbf{#1}}\par\vspace{3pt}{\color{ink}\LARGE\bfseries #2}\par\vspace{8pt}}
\newcommand{\fig}[2]{\begin{center}\includegraphics[width=\linewidth]{figures/#1.pdf}\end{center}{\small\color{ink}#2}\par}
\newcommand{\claim}[1]{\par\textbf{\color{accent}#1}\par}
\newcommand{\smallfig}[3][0.78]{\begin{center}\includegraphics[width=#1\linewidth]{figures/#2.pdf}\end{center}{\small\color{ink}#3}\par}
\newcommand{\dd}{\,d}
\newcommand{\A}{A_{a,\beta,Q}}
\newcommand{\Z}{Z_Q}

% Operation Komorebi: shared publication profile.
\usepackage{komorebi-publication}
\renewcommand{\KomorebiTitle}{Function Families and Their Structure}
\renewcommand{\KomorebiSubtitle}{Coupled growth, arithmetic depth, and finite resolution}
\renewcommand{\KomorebiSeries}{SPECIAL FUNCTIONS}
\renewcommand{\KomorebiKind}{MATHEMATICAL MONOGRAPH}
\renewcommand{\KomorebiEdition}{}
\renewcommand{\KomorebiShortTitle}{Function Families and Their Structure}
\renewcommand{\KomorebiTitleSize}{33}
\renewcommand{\KomorebiTitleLeading}{36}
\renewcommand{\KomorebiProfile}{classical}
\renewcommand{\KomorebiOrientation}{portrait}
\KomorebiSetup
% Realization Calculus: mathematical monograph visual system.
% This file styles the presentation. The mathematical source files remain separate.
\definecolor{RCPlum}{HTML}{3B283E}
\definecolor{RCMulberry}{HTML}{76526B}
\definecolor{RCCopper}{HTML}{AB8564}
\definecolor{RCInk}{HTML}{2F2B31}
\definecolor{RCMuted}{HTML}{69646B}
\definecolor{RCCream}{HTML}{F8F6F2}
\definecolor{RCLine}{HTML}{DFD9D4}
\definecolor{RCPaper}{HTML}{FFFFFF}
\colorlet{ink}{RCInk}
\colorlet{navy}{RCPlum}
\colorlet{accent}{RCMulberry}
\colorlet{muted}{RCMuted}
\colorlet{KomorebiInk}{RCInk}
\colorlet{KomorebiPine}{RCPlum}
\colorlet{KomorebiMuted}{RCMuted}
\geometry{letterpaper,left=.84in,right=.84in,top=.89in,bottom=.82in,headheight=17pt,headsep=17pt,footskip=24pt}
\setlength{\parindent}{0pt}
\setlength{\parskip}{6.1pt plus 1pt minus 1pt}
\linespread{1.08}
\setlength{\emergencystretch}{3em}
\setlength{\headheight}{17pt}
\pagestyle{fancy}
\fancyhf{}
\fancyhead[L]{\sffamily\fontsize{8}{10}\selectfont\color{RCMuted}\MakeUppercase{Realization Calculus}}
\fancyhead[R]{\sffamily\fontsize{8}{10}\selectfont\color{RCMuted}Special Functions}
\fancyfoot[L]{\sffamily\fontsize{8}{10}\selectfont\color{RCMuted}Function Families and Their Structure}
\fancyfoot[R]{\sffamily\fontsize{8}{10}\selectfont\color{RCPlum}\thepage}
\renewcommand{\headrulewidth}{.24pt}
\renewcommand{\footrulewidth}{0pt}
\fancypagestyle{plain}{%
 \fancyhf{}\fancyfoot[R]{\sffamily\small\color{RCPlum}\thepage}%
 \renewcommand{\headrulewidth}{0pt}%
}
\titleformat{\section}{\rmfamily\Large\color{RCPlum}}{\thesection}{.7em}{}
\titleformat{\subsection}{\rmfamily\large\bfseries\color{RCInk}}{\thesubsection}{.7em}{}
\titleformat{\subsubsection}{\rmfamily\normalsize\bfseries\color{RCInk}}{\thesubsubsection}{.7em}{}
\titlespacing*{\section}{0pt}{2.1ex plus .8ex minus .2ex}{.9ex}
\titlespacing*{\subsection}{0pt}{1.8ex plus .6ex minus .2ex}{.6ex}
\hypersetup{colorlinks=true,linkcolor=RCPlum,citecolor=RCPlum,urlcolor=RCMulberry,
  pdftitle={Function Families and Their Structure},
  pdfauthor={Komorebi P.},
  pdfsubject={Realization Calculus: Special Functions}}

\renewcommand{\KomorebiCover}{%
 \clearpage\thispagestyle{empty}\begingroup
 \begin{tikzpicture}[remember picture,overlay]
 \fill[RCPlum] (current page.south west) rectangle (current page.north east);
 \draw[RCCopper,line width=1.1pt] ([xshift=.82in,yshift=-.87in]current page.north west) -- ([xshift=2.03in,yshift=-.87in]current page.north west);
 \draw[RCCopper!50,line width=.45pt] ([xshift=.82in,yshift=.72in]current page.south west) -- ([xshift=-.82in,yshift=.72in]current page.south east);
 \foreach \k in {0,...,4}{
   \draw[RCCopper!23,line width=.55pt] ([xshift=-.75in,yshift=-1.20in-\k*.24in]current page.north east)
      .. controls ++(-.6in,-.1in) and ++(.3in,.4in) .. ([xshift=-2.4in,yshift=-1.80in-\k*.16in]current page.north east);
 }
 \end{tikzpicture}
 \vspace*{.32in}
 {\sffamily\fontsize{8.5}{13}\selectfont\color{RCCopper}\bfseries REALIZATION CALCULUS\par}
 \vspace*{1.58in}
 {\color{white}\fontsize{35}{41}\selectfont\bfseries Function Families\\[.11em] and Their Structure\par}
 \vspace{19pt}
 {\color{white!78!RCCream}\fontsize{14}{19}\selectfont Coupled growth, arithmetic depth,\\ finite resolution\par}
 \vspace{29pt}{\color{RCCopper}\rule{1.18in}{.75pt}}\par
 \vspace{13pt}{\color{white}\sffamily\fontsize{9}{12}\selectfont MATHEMATICAL MONOGRAPH / SPECIAL FUNCTIONS\par}
 \vfill
 {\color{white}\fontsize{12}{16}\selectfont Komorebi P.\par}
 \vspace{.47in}
 \endgroup\clearpage
}

\newcommand{\MonographTitleLeaf}{%
 \thispagestyle{empty}\begingroup
 \vspace*{1.03in}
 {\sffamily\fontsize{8}{12}\selectfont\color{RCMulberry}\bfseries SPECIAL FUNCTIONS\par}
 \vspace{.50in}
 {\color{RCPlum}\fontsize{31}{36}\selectfont Function Families\\ and Their Structure\par}
 \vspace{.21in}
 {\color{RCMuted}\fontsize{14}{19}\selectfont Coupled growth, arithmetic depth,\\ and finite resolution\par}
 \vspace{.37in}{\color{RCCopper}\rule{.91in}{.7pt}}\par
 \vfill
 {\fontsize{12}{16}\selectfont Komorebi P.\par}
 \vspace{8pt}{\color{RCMuted}\sffamily\fontsize{9}{13}\selectfont A research monograph in five parts\par}
 \vspace{.84in}
 \endgroup\clearpage
}

\newcommand{\MonographPreface}{%
 \phantomsection\addcontentsline{toc}{section}{Preface}
 {\sffamily\fontsize{8}{11}\selectfont\color{RCMulberry}\bfseries PREFACE}\par\vspace{6pt}
 {\fontsize{25}{30}\selectfont\color{RCPlum}Five different families, one problem of structure\par}
 \vspace{14pt}
 The theory of special functions is built from examples whose definitions carry more information than their initial values suggest. A contraction hidden in a differential equation can determine the growth of an entire function. A prime omitted from a sum can alter its parameter signature. An apparently positive finite moment calculation can conceal a negative mass elsewhere in the underlying arithmetic measure.

 This volume studies these phenomena through explicit families. Some results are extensions of familiar analytic constructions. Others are concrete applications of classical machinery in new coordinates. The proofs keep those distinctions visible. A useful representation does not by itself establish historical novelty, and a numerical illustration does not substitute for a universal argument.

 The exposition begins with scalar nonlinear feedback, passes through complex arithmetic depth and signed prime measures, and reaches spectral zero orbits. It then develops coupled functions and precise limits on finite observation. A final part collects the longer proofs, reproducible calculations, and connections to existing literature.

 \vspace{10pt}
 {\color{RCPlum}\bfseries A note on evidence}\par
 Analytic identities and their domains are stated explicitly. Finite computations are reported with their actual ranges, and unresolved regularity, identifiability, and priority questions remain open. These qualifications are part of the mathematics rather than temporary editorial notes.

 \vspace{10pt}
 {\color{RCPlum}\bfseries Reading this book}\par
 The main chapters are designed to be read consecutively. The technical appendices can be consulted at the point where a proof needs more detail. References and verification records are collected at the back so that the development itself reads as a continuous monograph.
 \clearpage
}

\renewcommand{\pageheading}[2]{%
 \clearpage\phantomsection\addcontentsline{toc}{section}{#2}\markright{#2}
 {\sffamily\fontsize{8}{11}\selectfont\color{RCMulberry}\bfseries\MakeUppercase{#1}\par}
 \vspace{6pt}
 {\rmfamily\fontsize{20.5}{26}\selectfont\color{RCPlum}#2\par}
 \vspace{7pt}{\color{RCCopper}\rule{.87in}{.65pt}}\par
 \vspace{9pt}
}
\renewcommand{\claim}[1]{\par\vspace{2.5pt}\noindent{\color{RCMulberry}\bfseries #1}\par\vspace{1pt}}
\newcommand{\MonographPart}[3]{%
 \clearpage\phantomsection\addcontentsline{toc}{part}{Part #1: #2}\thispagestyle{empty}
 \vspace*{1.45in}
 {\sffamily\fontsize{9}{12}\selectfont\color{RCMulberry}\bfseries PART #1\par}
 \vspace{16pt}
 {\rmfamily\fontsize{30}{35}\selectfont\color{RCPlum}#2\par}
 \vspace{18pt}{\color{RCCopper}\rule{1.05in}{.8pt}}\par
 \vspace{17pt}
 {\fontsize{12}{18}\selectfont\color{RCMuted}#3\par}
 \vfill
 {\sffamily\fontsize{8}{12}\selectfont\color{RCMuted}REALIZATION CALCULUS\enspace /\enspace SPECIAL FUNCTIONS\par}
 \clearpage
}
\newcommand{\MonographSubpart}[1]{%
 \clearpage\phantomsection\addcontentsline{toc}{section}{#1}
 {\sffamily\fontsize{9}{12}\selectfont\color{RCMulberry}\bfseries REFERENCE SECTION\par}
 \vspace{8pt}{\fontsize{24}{28}\selectfont\color{RCPlum}#1\par}
 \vspace{12pt}{\color{RCCopper}\rule{1in}{.7pt}}\par\vspace{14pt}
}

\newcommand{\vfig}[3][.92]{\begin{center}\includegraphics[width=#1\linewidth]{figures_v11/#2.pdf}\end{center}{\small\color{ink}#3}\par}
\begin{document}

\KomorebiCover
\MonographTitleLeaf
\MonographPreface
\begingroup\setlength{\parskip}{0pt}\tableofcontents\endgroup
\clearpage
\pageheading{A READING MAP}{Five families, one mathematical question}
This monograph investigates structural invariants in five families of special functions. Feedback scales, arithmetic frequencies, signed prime masses, rotation orbits, and coupled channels provide distinct settings. Their connections are established through specific formulas, invariants, and obstructions, not assumed from a shared vocabulary.

The argument starts with entire functions defined by contracted nonlinear feedback. It then studies arithmetic generating functions and prime-sieve measures, before passing to orbit spectral zeta functions and exact coefficient laboratories. Coupled entire systems connect several of these ideas, while a resolution atlas measures the difference between finite observation and global mathematical structure.

The main chapters give defining formulas, stated hypotheses, mathematical consequences, and computational figures. Extended proofs follow as technical appendices. The final notes distinguish analytic proofs, finite checks, classical inputs, and questions that remain open.

\pageheading{GUIDE / FOUR STARTING FAMILIES}{Four starting constructions}
The first four families lead to distinct results: an entire-growth law, an obstruction to finite differential equations, a positivity boundary, and a rotation-invariant spectral coordinate. A coupled fifth family is developed in Part IV.

\renewcommand{\arraystretch}{1.45}
\begin{tabular}{@{}>{\raggedright\arraybackslash}p{0.21\linewidth}>{\raggedright\arraybackslash}p{0.40\linewidth}>{\raggedright\arraybackslash}p{0.33\linewidth}@{}}
\toprule
Family & Structure worked out & Status of the construction\\
\midrule
Asymmetric dilation & Entire solutions; exact growth balance;
off-diagonal nonrescaling; intrinsic feedback banks & Candidate extension of a
known one-dilation tangent analog\\
Complex arithmetic depth & Exact exponential type; coupled shift calculus;
strict positive kernel; exact parameter recovery; no polynomial-coefficient linear ODE & Candidate
extension of the project's integer arithmetic energy\\
Prime-sieve packets & Exact moment field; real Hankel singularities;
positivity-margin collapse; multiple-prime geometry; arithmetic resolvent & Candidate arithmetic
specialization of classical orthogonal-polynomial machinery\\
Orbit spectral zeta & Global meromorphic continuation; rational integer values;
regularized product; multiplicity-aware orbit quotient & Worked specialization of known
Mittag--Leffler seeds and spectral-zeta methods\\
\bottomrule
\end{tabular}

\claim{An established diagonal antecedent}
The equation $f_q'(z)=1+f_q(qz)^2$, $f_q(0)=0$, is already present in
Miguel A. M\'endez's 2024 path-length paper, Example 2.4. It remains an
important baseline, but its definition is not a discovery of this project.
The asymmetric extension keeps two separate contractions and permits a
proved distinction from every amplitude/argument rescaling of that baseline.

\claim{How to read the evidence}
Theorem statements below concern exact infinite objects under the displayed
hypotheses. Plots and decimal roots illustrate them. Finite assertions check
implementations within recorded ranges. None of these three forms of evidence
establishes historical priority by itself.

The source files preserve extended arguments, independent reviews, negative controls, and open questions. A separate evidence register records the exact files and calculations reviewed.


\MonographPart{I}{Entire Functions and Their Growth}{Contractions, coefficients, scale balance, and nonlinear feedback taps.}
\pageheading{FAMILY 1}{Two scales inside the function}
For real $0<q,r<1$, define the normalized odd solution
\[
\boxed{T_{q,r}'(z)=1+T_{q,r}(qz)T_{q,r}(rz),\qquad T_{q,r}(0)=0.}
\]
The contractions are part of the defining feedback. The function is symmetric
in $q,r$, but unequal scales retain left/right address information.
\fig{01_asymmetric_functions}{Computed Taylor curves. Each entire curve uses
terms through $z^{321}$; positive-coefficient Cauchy estimates bound the omitted
tail on each displayed interval. The black dashed curve is ordinary tangent.}
For $T_{q,r}(z)=\sum_{n\ge0}c_nz^{2n+1}$, $c_0=1$, coefficient matching gives
\[
c_n=\frac{1}{2n+1}\sum_{i=0}^{n-1}
q^{2i+1}r^{2(n-1-i)+1}c_ic_{n-1-i},\qquad n\ge1.
\]
In particular,
\[
T_{q,r}(z)=z+\frac{qr}{3}z^3+
\frac{q^2r^2(q^2+r^2)}{15}z^5+O(z^7).
\]
The coefficients are positive and bounded by tangent coefficients, giving a
convergent germ. If the solution is known on $|z|<R$, its integral equation
extends it to $|z|<R/\max(q,r)$. Repetition proves that it is entire.

The parameters can be recovered, up to swapping, from $c_1,c_2$. Put
$p=3c_1$ and $u=15c_2/p^2$. Then $q,r$ are the roots of
$x^2-\sqrt{u+2p}\,x+p=0$.

\pageheading{FAMILY 1\; /\; GROWTH}{A scale-balance atlas}
\[
\boxed{\rho(T_{q,r})=\rho\quad\text{where}\quad q^\rho+r^\rho=1.}
\]
Here $\rho(f)=\limsup_{R\to\infty}\log\log M_f(R)/\log R$ is entire-function
order. Moreover $\log M_{T}(R)=O(R^\rho)$: the upper type is finite.
\smallfig{02_order_atlas}{The displayed values numerically solve the proved
characteristic equation. Colors saturate at order 7; the upper corner contains
larger orders. The dashed diagonal is the previously known one-scale family.}
\begin{tabular}{@{}lll@{}}\toprule
Scales & Order & Interpretation\\\midrule
$q=r=1/4$ & $1/2$ & Slow entire growth\\
$q+r=1$ & $1$ & Many distinct functions at one order\\
$q=r=1/\sqrt2$ & $2$ & Quadratic logarithmic growth scale\\
$q=r=3/4$ & $\log2/\log(4/3)$ & Approximately $2.4094208397$\\
\bottomrule\end{tabular}

\claim{Why this equation appears}
Nonnegative coefficients give $M(R)=T_{q,r}(R)$. The integral equation implies
\[
M(R)\le R[1+M(qR)M(rR)],
\]
\[
M(R)\ge(1-\alpha)R M(\alpha qR)M(\alpha rR),\quad0<\alpha<1.
\]
Taking logarithms turns product feedback into a sum over contracted scales.
Comparison with $CR^\rho$ proves the upper bound; comparison with
$cR^{\rho_\alpha}$, where $(\alpha q)^{\rho_\alpha}+
(\alpha r)^{\rho_\alpha}=1$, proves the lower bound. Letting $\alpha\uparrow1$
gives the same order. A complete interval-induction proof is in the source route.

\pageheading{DILATION\; /\; GROWTH THEOREM}{The exact order has positive lower type}
The order theorem can now be strengthened. For every finite positive bank,
with $\rho=\max_{h_j>0}\rho_j$,
\[
\boxed{0<\liminf_{R\to\infty}\frac{\log M(R)}{R^\rho}
\le\limsup_{R\to\infty}\frac{\log M(R)}{R^\rho}<\infty.}
\]
Thus logarithmic growth is trapped between two positive multiples of
$R^\rho$. In particular,
\[
\lim_{R\to\infty}\frac{\log\log M(R)}{\log R}=\rho.
\]
The ratio $\log M(R)/R^\rho$ itself need not have a proved limit here.
\fig{08_growth_type}{Finite-radius logarithmic growth and its proved lower
comparison. The normalized curves illustrate scale dependence; the plot is
not a proof of convergence to a type constant. Analytic truncation envelopes
and ordinary high-precision rounding are recorded separately.}
\claim{The step that removes the old loss of exponent}
Choose a maximizing channel $(h,q,r)$ and integrate over a fixed-width
terminal interval. Let
\[
L=1/h,\quad C=\max\left(\frac{qL}{1-q},\frac{rL}{1-r}\right),\quad
B=L+\frac{2}{\min(q,r)},\quad c_-=\frac{\log2}{(B+C)^\rho}.
\]
The result is the explicit all-radius bound
\[
\log M(R)\ge c_-(R+C)^\rho,\qquad R\ge2.
\]
Unlike a fixed fractional window, the terminal window loses only an additive
radius. The shift $C$ absorbs that loss without changing the exponent.

\pageheading{DILATION\; /\; PROOF}{A translated power comparison}
Write $H(R)=\log M(R)$. Positivity gives
\[
M(R)\ge h\int_{R-L}^{R}M(qt)M(rt)\dd t
\ge M(q(R-L))M(r(R-L)),\qquad R\ge B.
\]
The choice of $C$ ensures
\[
q(R-L)+C\ge q(R+C),\qquad r(R-L)+C\ge r(R+C).
\]
On the seed interval $[2,B]$, $H(R)\ge\log R\ge\log2\ge c_-(R+C)^\rho$.
If this comparison holds at the two child radii, then
\[
\begin{aligned}
H(R)&\ge H(q(R-L))+H(r(R-L))\\
&\ge c_-\{[q(R-L)+C]^\rho+[r(R-L)+C]^\rho\}\\
&\ge c_-(q^\rho+r^\rho)(R+C)^\rho=c_-(R+C)^\rho.
\end{aligned}
\]
Let $Q_* =\max(q,r)$ and $d=(1-Q_*)B+Q_*L>0$. Each child radius is at
most $R-d$ and at least two. Induction on successive intervals of length
$d$ propagates the seed comparison to every $R\ge B$. The argument uses
monotonicity of positive powers and works also when $\rho<1$.

\claim{Explicit upper constants are available too}
Set $Q=\max_j\max(q_j,r_j)$, $m=\min_j\min(q_j,r_j)$,
$\Delta=\log(1/m)$, and $B_+=\max_j\log(1/(q_jr_j))+\log2$.
Choose an integer $K$ with $Q^{-K}\ge3/m$ and put
\[
V_K=(2^{K+1}-1)\log2+(2^{K+1}-K-2)\log(1/Q),
\]
\[
c_+=e^{-\rho}(V_K+1+\Delta+B_+).
\]
Then $\log M(R)\le c_+R^\rho-\log R-B_+$ for $R\ge e$.
The tangent majorant $M(1)<2$, followed by the integral recurrence on radii
$Q^{-k}$, bounds the initial seed by $V_K$. The logarithmic-radius comparison
then supplies the displayed upper constant. The full induction is in Appendix A.

These constants are constructive and conservative. Their purpose is to
establish two-sided growth, not to estimate the optimal type.


\pageheading{FAMILY 1\; /\; SEPARATION}{More than a rescaled one-scale seed}
\claim{Nonrescaling theorem}
If $q\ne r$, then $T_{q,r}$ cannot equal $A f_s(Bz)$ for nonzero constants
$A,B$ and any real $0<s<1$, where $f_s'=1+f_s(sz)^2$ is the known family.

Suppose such an equality held. Normalization gives $AB=1$. Write $p=qr$,
$u=q^2+r^2$, and let $\rho$ be the common order. The first two nontrivial
coefficients and the growth law force
\[
B^2s^2=p,\qquad s^2=u/2,\qquad 2s^\rho=q^\rho+r^\rho=1.
\]
Consequently
\[
2\left(\frac{q^2+r^2}{2}\right)^{\rho/2}=q^\rho+r^\rho.
\]
For $\rho\ne2$, strict convexity or concavity of $x^{\rho/2}$ forces $q=r$.
For the exceptional case $\rho=2$, we have $u=1$, $s^2=1/2$, $B^2=2p$.
The two $z^7$ coefficients become
\[
\frac{p^3(3-p^2)}{315}\quad\text{and}\quad\frac{11p^3}{1260}.
\]
Their equality forces $p=1/2$ and again $(q-r)^2=u-2p=0$.
This handles the exceptional order rather than assuming it away.

\claim{A worked order-one member}
\[
T_{1/3,2/3}(z)=z+\frac{2}{27}z^3+
\frac{4}{2187}z^5+\frac{232}{11160261}z^7+O(z^9).
\]
It has order one, and $T_{1/3,2/3}(1)\approx1.0759239897408249$.
The diagonal point $q=r=1/2$ gives
$f_{1/2}(z)=\sqrt2\sinh(z/\sqrt2)$. Every off-diagonal member is separated
from all rescaled diagonal members by the theorem above.

\claim{A combinatorial address formula}
Sum over full ordered binary trees with $n$ internal vertices. A vertex with
$i$ internal vertices in its left subtree and $j$ in its right contributes
\[
\frac{q^{2i+1}r^{2j+1}}{2(i+j+1)+1}.
\]
Multiply these factors over internal vertices, then sum over trees. Root
splitting reproduces the coefficient recurrence exactly. Unequal contractions
therefore record left and right path lengths separately.

This is a precise separation result from one authenticated antecedent.
It does not prove that no earlier author studied the asymmetric equation.

\pageheading{FAMILY 1\; /\; TAPS}{Feedback banks and weak-channel sensitivity}
Choose finitely many channels $j=1,\ldots,J$, with $h_j>0$, $\sum_jh_j=1$,
and $0<q_j,r_j<1$. Define
\[
F_\Theta'(z)=1+\sum_jh_jF_\Theta(q_jz)F_\Theta(r_jz),\qquad F_\Theta(0)=0.
\]
These taps act inside the nonlinear definition. With
$\mu_1=\sum_jh_jq_jr_j$ and
$\mu_2=\sum_jh_jq_jr_j(q_j^2+r_j^2)$,
\[
F_\Theta(z)=z+\frac{\mu_1}{3}z^3+\frac{\mu_1\mu_2}{15}z^5+O(z^7).
\]
If $q_j^{\rho_j}+r_j^{\rho_j}=1$, then
\[
\boxed{\rho(F_\Theta)=\max_{j\text{ active}}\rho_j.}
\]
The positive weights alter local coefficients; the most expansive active
channel determines the asymptotic order. A zero-weight channel is removed
before taking the maximum.
\fig{07_weak_tap}{A two-channel example: scales $(1/4,1/4)$ with weight
$1-\varepsilon$, and $(3/4,3/4)$ with weight $\varepsilon$.
The function varies continuously on every fixed compact set, while its growth
order jumps when the second channel becomes active.}
For $0<\varepsilon\le1$ the order is $\log2/\log(4/3)$; at $\varepsilon=0$
it is $1/2$. Coefficients depend continuously on $\varepsilon$ and are
dominated by those of the entire $f_{3/4}$. Dominated convergence of the Taylor
series therefore proves uniform convergence on compact sets as
$\varepsilon\downarrow0$.

This gives a useful design distinction: finite-window similarity does not
control entire growth order. For a bank with weights $1/4,3/4$ at these scales,
\[
F(z)=z+\frac7{48}z^3+\frac{427}{30720}z^5+
\frac{176681}{188743680}z^7+O(z^9).
\]

\pageheading{DILATION\; /\; TAPS}{A quantitative compact-continuity law}
Fix the contractions and let $h,g$ be two normalized nonnegative tap vectors.
Set $Q=\max_j\max(q_j,r_j)<1$ and $d=\|h-g\|_1$. Then
\[
\boxed{\sup_{|z|\le R}|F_h(z)-F_g(z)|
\le\frac{d}{4}\,[R f_Q'(R)-f_Q(R)].}
\]
The right side is an explicitly defined entire majorant. It measures local
sensitivity without asserting continuity of the asymptotic growth order.
\fig{09_tap_continuity}{Computed weak-tap differences lie below the proved
compact bound. For a slow/faster two-channel bank, $d=2\varepsilon$, so the
envelope is $\varepsilon[R f_{3/4}'(R)-f_{3/4}(R)]/2$.}
Each $c_n(h)$ is a positive polynomial homogeneous of degree $n$ in the
weights. Coefficientwise comparison gives
$c_n(h)\le d_n(Q)(\sum_jh_j)^n$, where $f_Q=\sum_nd_n(Q)z^{2n+1}$.
Every weight partial derivative on the simplex lies in $[0,nd_n(Q)]$.
The positive and negative parts of $h-g$ each have mass $d/2$, giving
$|c_n(h)-c_n(g)|\le(d/2)nd_n(Q)$.
Summing this inequality uses
$\sum_nnd_n(Q)R^{2n+1}=[Rf_Q'(R)-f_Q(R)]/2$.

\pageheading{DILATION\; /\; PARAMETERS}{An analytic family with a discontinuous invariant}
For fixed finite $J$, allow complex $h_j,q_j,r_j$, with
$|q_j|,|r_j|<1$. The coefficient recurrence defines polynomials in all these
parameters. On a compact parameter set choose $Q<1$ and $W>0$ such that
$|q_j|,|r_j|\le Q$ and $\sum_j|h_j|\le W$. Then
\[
|c_n(h,q,r)|\le W^n d_n(Q),
\qquad
\sum_n|c_n|R^{2n+1}\le W^{-1/2}f_Q(\sqrt W R).
\]
Normal convergence proves joint holomorphic dependence on parameters and
argument throughout this domain. Signed and complex banks have entire
solutions, but the positive-bank growth theorem is not asserted for them.

\claim{The weak-tap example now has both sides quantified}
For the bank with scales $(1/4,1/4)$ and $(3/4,3/4)$, weights
$1-\varepsilon,\varepsilon$, the functions converge on every compact set
as $\varepsilon\downarrow0$. The preceding envelope gives a linear bound
in $\varepsilon$. At every $\varepsilon>0$, the order is
$\log2/\log(4/3)$, while at zero it is $1/2$.

The new lower-type theorem does not contradict this: its lower constant
depends on the active faster weight and can tend to zero as that channel
vanishes. For $\varepsilon=10^{-3}$ one constructive lower constant is
about $1.4496\times10^{-9}$; a conservative upper constant is about $90.4392$.

\claim{Identifying a bank has unavoidable equivalences}
Permuting channels, swapping the scales within a channel, or merging
duplicate scale pairs by adding their weights leaves the function unchanged.
The two-scale single-channel inverse is exact up to swapping. A general
minimal-bank identifiability theorem beyond these equivalences is not proved.

The resulting research target is specific: determine optimal growth type,
possible logarithmic oscillations, and minimal-bank identifiability. A finite
plot cannot settle any of those questions.



\MonographPart{II}{Arithmetic Depth and Signed Measures}{Complex arithmetic parameters, positive kernels, and the exact boundaries of prime-sieve positivity.}
\pageheading{FAMILY 2}{Arithmetic energy at complex depth}
For real $a>1/2$, $\beta>0$, and integer $Q\ge1$, define
\[
\boxed{\A(s)=\sum_{\substack{n\ge2\\(n,Q)=1}}
\mu(n)^2n^{-2a}(1-n^{-\beta})^{2s},\qquad s\in\mathbb C.}
\]
The complex power uses the real logarithm of $1-n^{-\beta}>0$.
The atom $n=1$ is excluded because it has no such logarithm. Thus positive
integer depths match the inherited energy, while $A(0)=E(0)-1$.
\fig{03_complex_depth}{Left: $a=\beta=2$ at real depths, using the exact
zeta--Newton representation at 90-digit working precision. Right: initial
squarefree frequencies and weights for $Q=1$. Only finitely many atoms are
shown; the proof uses the infinite support.}
Set $\lambda_n=-2\log(1-n^{-\beta})$. Then
$A(s)=\sum w_ne^{-\lambda_ns}$, $w_n=\mu(n)^2n^{-2a}$.
The weights are summable, and the distinct frequencies decrease to zero.
If $p_*$ is the least prime not dividing $Q$, then
\[
\tau=-2\log(1-p_*^{-\beta}),\qquad M=\Z(2a)-1,
\]
\[
|A(s)|\le M e^{\tau\max(-\Re s,0)}.
\]
This proves local uniform convergence and entire continuation. The first atom
dominates $A(-t)$ as $t\to\infty$, so the entire function has exact order one
and exponential type $\tau$.

\pageheading{FAMILY 2\; /\; CALCULUS}{Closed shifts, infinite differential complexity}
Write
\[
\Z(v)=\frac{\zeta(v)}{\zeta(2v)}\prod_{p\mid Q}(1+p^{-v})^{-1},\qquad v>1.
\]
Only $\operatorname{rad}(Q)$ matters. Expansion of the depth factor gives
\[
\boxed{\A(s+1)=\A(s)-2A_{a+\beta/2,\beta,Q}(s)
+A_{a+\beta,\beta,Q}(s).}
\]
The fractional transfer and Newton formulas are
\[
\A(s+t)=\sum_{j\ge0}(-1)^j\binom{2t}{j}A_{a+j\beta/2,\beta,Q}(s),
\]
\[
\A(t)=\sum_{j\ge0}(-1)^j\binom{2t}{j}[\Z(2a+j\beta)-1].
\]
They converge absolutely and locally uniformly: binomial coefficients have
at most polynomial growth on compact parameter sets, while the shifted sums
decay geometrically. The subtraction of one is essential at noninteger depth.

\claim{No polynomial-coefficient linear ODE}
Suppose $\sum_{j=0}^{r}P_j(s)A^{(j)}(s)=0$ with polynomial $P_j$.
At $s=-t$, divide by the exponential of the first support atom. The remaining
tail is exponentially smaller, so the polynomial
$\sum_jP_j(-t)(-\lambda_{p_*})^j$ tends to zero exponentially and must vanish
identically. Remove that atom and repeat. Every support frequency therefore
satisfies $\sum_jP_j(s)(-\lambda_n)^j=0$ for all $s$.
For fixed $s$, this polynomial in $\lambda$ has infinitely many distinct
roots. All its coefficients vanish. Thus all $P_j$ are zero, a contradiction.

The function is \textbf{not D-finite in depth}. This excludes identification
in the depth variable with fixed finite-parameter hypergeometric or Meijer
G functions and ordinary finite D-finite combinations. It does not exclude
nonalgebraic substitutions or representations by known non-D-finite functions.

\claim{Strict positive complex kernel}
For distinct complex points $s_i$, $K_{ij}=A(s_i+\overline{s_j})$ is strictly
positive definite. A zero quadratic form would force a finite exponential
sum to vanish at every $\lambda_n$. Their accumulation at zero invokes the
identity theorem; derivatives at zero then give a nonsingular Vandermonde
system. On the real axis, $(-1)^jA^{(j)}(s)>0$ for every $j\ge0$.

\pageheading{FAMILY 2\; /\; ASYMPTOTICS}{An arithmetic sector law}
The negative real direction is governed by the first admissible prime. The
positive direction is governed by the density of infinitely many squarefree
integers. These are different pieces of the same entire function.

Let
\[
\alpha=\frac{2a-1}{\beta},\qquad
\delta_Q=\frac1{\zeta(2)}\prod_{p\mid Q}\frac{p}{p+1}.
\]
As $|s|\to\infty$ in any closed sector
$|\arg s|\le\pi/2-\epsilon$ with $\epsilon>0$,
\[
\boxed{\A(s)=\frac{\delta_Q}{\beta}B(\alpha,2s+1)
+O_{a,\beta,Q,\epsilon}\left(|s|^{-(2a-1/2)/\beta}\right).}
\]
In particular,
\[
\A(s)\sim\frac{\delta_Q\Gamma(\alpha)}{\beta}(2s)^{-\alpha}.
\]
The principal power is used in the right half-plane. The beta expression is
an asymptotic continuous-density model, not an identity with the arithmetic sum.

\claim{Proof mechanism}
Squarefree inclusion--exclusion gives
\[
\sum_{\substack{n\le x\\(n,Q)=1}}\mu(n)^2
=\delta_Q x+O_Q(\sqrt{x}).
\]
Integrate against $g_s(x)=x^{-2a}(1-x^{-\beta})^{2s}$ on $[1,\infty)$.
The density term becomes the beta integral after $y=x^{-\beta}$. Partial
integration bounds the error by
$C_Q\int_1^\infty\sqrt{x}|g_s'(x)|\dd x$.
The two derivative terms are beta integrals with decay exponent
$(2a-1/2)/\beta$; $|s|/\Re s$ stays bounded in the stated sector.
This gives the error term and then the leading asymptotic.

\claim{One complex example}
At $a=\beta=2$, $Q=1$,
\[
A(1/2+i/4)\approx
0.06033183289343401791658
-0.007412384053785436251981\,i.
\]
Direct finite sums and a zeta--Newton evaluation agree within the recorded
omitted-tail envelope. This is an ordinary high-precision example, not an
interval certificate.

Mixed Dirichlet-character weights provide matrix-valued positive kernels.
Individual off-diagonal entries may be signed or complex, so complete
monotonicity is asserted only for the positive principal family.

\pageheading{ARITHMETIC DEPTH\; /\; INVERSE}{The function encodes its parameters}
The complete exact analytic germ determines
\[
\boxed{a,\quad\beta,\quad\operatorname{rad}(Q).}
\]
Precisely, two admissible triples define the same germ if and only if these
three quantities agree. Prime powers in $Q$ are invisible because the terms
use only coprimality. A finite Taylor polynomial or noisy measurements are
different inputs from the complete germ in this theorem.

Let the support integers be $n_1<n_2<\cdots$, frequencies
$\lambda_j=-2\log(1-n_j^{-\beta})$, and weights $w_j=n_j^{-2a}$.
Starting at $R_1=A$, peel the ordered atoms:
\[
\lambda_j=\lim_{t\to\infty}\frac{\log R_j(-t)}t,\qquad
w_j=\lim_{t\to\infty}e^{-\lambda_jt}R_j(-t),
\]
\[
R_{j+1}(s)=R_j(s)-w_je^{-\lambda_js}.
\]
Every limit follows from the isolated largest remaining frequency and the
positive summable tail. Thus the exact function determines its spectral list.

The counting law $N_Q(x)=\delta_Qx+O_Q(\sqrt x)$ gives
\[
\lambda_j=2\delta_Q^\beta j^{-\beta}
[1+O_{Q,\beta}(j^{-1/2}+j^{-\beta})],
\quad
\beta=-\lim_{j\to\infty}\frac{\log\lambda_j}{\log j}.
\]
Once $\beta$ is recovered,
\[
n_j=(1-e^{-\lambda_j/2})^{-1/\beta},\qquad
a=-\frac{\log w_j}{2\log n_j}.
\]
Prime membership in this integer support then recovers the excluded prime
set. The reconstruction is exact from infinite data; it does not certify
that a finite search has found every excluded prime.

\pageheading{ARITHMETIC DEPTH\; /\; RESOLUTION}{Separation improves while subtraction worsens}
\fig{10_arithmetic_inverse}{The frequency-rank estimates converge toward
the depth exponent. Later-atom peeling at large negative depth also has a
rapidly growing subtraction condition factor. These are finite examples
from the checked arithmetic run, not certified floating-point intervals.}
If earlier atoms have been removed exactly, set
$\Delta_j=\lambda_j-\lambda_{j+1}>0$, $M_{j+1}=\sum_{k>j}w_k$.
The logarithmic secant
\[
\widehat\lambda_j(t,h)=\frac{\log R_j(-t-h)-\log R_j(-t)}h
\]
obeys, for $t\ge0$, $h>0$,
\[
0\le\lambda_j-\widehat\lambda_j(t,h)
\le\frac1h\log\left(1+\frac{M_{j+1}}{w_j}e^{-\Delta_jt}\right).
\]
This improves with negative depth. But relative error in $A(-t)$ can be
amplified after subtraction by
\[
\kappa_j(t)=\frac{A(-t)}{R_j(-t)}
\sim\frac{w_1}{w_j}e^{(\lambda_1-\lambda_j)t}.
\]
For the displayed $Q=1$ example, the third-atom subtraction at $t=500$
has a condition factor of roughly $10^{108.8}$. Exact uniqueness therefore
does not promise easy recovery from ordinary precision.

\pageheading{ARITHMETIC DEPTH\; /\; STABILITY}{A large prime can hide in every finite window}
For a prime $p\nmid Q$, any $R\ge0$, and integer $k\ge0$,
\[
\boxed{\sup_{|s|\le R}|A_{a,\beta,Q}^{(k)}(s)-A_{a,\beta,Qp}^{(k)}(s)|
\le\lambda_p^k e^{R\lambda_p}\frac{Z_Q(2a)}{p^{2a}+1}.}
\]
The difference consists of the $p$-divisible admissible squarefree atoms.
Their frequencies are at most $\lambda_p$ and their total weight is the
Euler-factor expression on the right. Termwise differentiation proves the
bound for every derivative.
\fig{11_prime_instability}{The proved compact-window envelope shrinks as
the excluded prime increases. At $a=\beta=2$, $Q=1$, $R=5$, excluding
$p=10007$ changes the entire window by at most approximately
$1.0749\times10^{-16}$. The depth-zero difference is known exactly.}
The radical differs for every such prime, while all bounded-window derivatives
converge to those of $A_Q$. Without a bound on excluded primes, there is no
uniformly stable inverse from these windows to the radical parameter.

The same distinction appears in weak taps: exact structural invariants can
change under perturbations that are small at finite resolution.

\pageheading{ARITHMETIC DEPTH\; /\; REALIZATION}{An infinite feature space with arithmetic labels}
For the ordered support, define $\phi_s(j)=\sqrt{w_j}e^{-\lambda_js}$.
The kernel $A(s+\overline t)$ is its weighted $\ell^2$ Gram kernel.
Summability and bounded frequencies make every feature vector well defined.

If a weighted $\ell^2$ combination produces the zero entire function,
Cauchy--Schwarz makes its coefficients absolutely summable. Successive
dominant-frequency peeling then forces every coefficient to vanish. This
gives an injective Hilbert realization, and all finite derivative Hankel
matrices have full rank: a nonzero polynomial in $\lambda$ cannot vanish
on the infinite frequency support.

\claim{Fixed-modulus character resolution}
For characters $\chi,\psi$ modulo a fixed $Q$, let
$\eta=\chi\overline\psi$. When $\eta$ is nonprincipal, the elementary
squarefree character sum is $O_Q(\sqrt x)$, with no density term.
The same sector partial-integration argument therefore gives
\[
A_{\chi,\psi}(s)
=O_{a,\beta,Q,\epsilon}(|s|^{-(2a-1/2)/\beta}),
\qquad |\arg s|\le\pi/2-\epsilon.
\]
Compared with the principal $|s|^{-(2a-1)/\beta}$ scale, this is faster
by $|s|^{-1/(2\beta)}$. The modulus is fixed; the constants are not asserted uniform as $Q$ grows. Individual off-diagonal functions can be complex and are not asserted
completely monotone.

Classical Laplace uniqueness provides a comparison for the inverse mechanism.
The parameter-specific extraction, frequency-rank law, and explicit instability
are proved locally in the complete arithmetic route.


\pageheading{FAMILY 3}{An infinite prime shelf changes the packets}
For a prime $p$ and $\lambda\ge-1$, let
\[
x_n=(1-n^{-2})^2,\qquad
\nu_{p,\lambda}=\sum_{n\ge1}\mu(n)^2n^{-4}
(1+\lambda\mathbf1_{p\mid n})\delta_{x_n}.
\]
Define monic orthogonal polynomials $P_m^{(k;p,\lambda)}$ for the measure
$x^k\dd\nu_{p,\lambda}$, $k\ge0$ integer. The zero atom has mass one at
$k=0$ and is annihilated when $k\ge1$.
\fig{04_prime_polynomials}{Left: degree-two packets for $p=2$, depth zero.
Right: initial atom masses. Strength $-1$ deletes every even squarefree atom;
strength $0$ is the original measure; strength $1$ doubles that shelf.}
The exact moments are finite zeta stencils:
\[
M_j^{p,\lambda}=\sum_{r=0}^{2j}(-1)^r\binom{2j}{r}
\frac{\zeta(4+2r)}{\zeta(8+4r)}
\left(1+\frac{\lambda}{p^{4+2r}+1}\right).
\]
All polynomial coefficients belong to $\mathbb Q(\pi^{-2},\lambda)$.
With $\Delta=M_0M_2-M_1^2$,
\[
P_1=x-\frac{M_1}{M_0},\qquad
P_2=x^2-\frac{M_0M_3-M_1M_2}{\Delta}x+
\frac{M_1M_3-M_2^2}{\Delta}.
\]
The infinitely many positive nodes make every Hankel matrix strictly positive
definite throughout $\lambda\ge-1$. Hence all degrees exist, with positive
norms, real simple zeros in $(0,1)$, and interlacing.

\pageheading{FAMILY 3\; /\; BOUNDARY}{Finite positivity can conceal a signed defect}
Split $\nu_{p,\lambda}=\nu_{\mathrm{non}p}+(1+\lambda)\nu_p$.
Let $A_m,B_m$ be their order-$m$ Hankel matrices. Both are positive definite.
Thus
\[
\tau_m(\lambda)=\det(A_m+(1+\lambda)B_m)
=\det A_m\prod_{i=1}^{m}[1+(1+\lambda)\eta_i],
\]
where $\eta_i>0$ are the eigenvalues of $A_m^{-1/2}B_mA_m^{-1/2}$.
Every determinant root is real and below $-1$; possible coefficient poles
occur there, although numerator cancellation is possible.
\fig{05_sieve_boundary}{The line is a proved bound for $p=2$; dots are
ordinary high-precision pencil roots. A defect $\lambda=-1.001$ eventually
violates positivity even though it can pass low-order checks.}
The nearest root is $-1-\delta_m$, with
\[
\delta_m=\min_{0\ne f,\,\deg f<m}
\frac{\int f^2\dd\nu_{\mathrm{non}p}}{\int f^2\dd\nu_p}>0,
\qquad\delta_m\downarrow0.
\]
To prove the limit, approximate a continuous bump supported at the isolated
node $x_p$ by polynomials on $[0,1]$. The numerator tends to zero while the
denominator retains the $p^{-4}$ atom. Hence every $\lambda<-1$ eventually
produces a negative finite-degree quadratic form.

For $p=2$, test with $f(x)=x(1-x)^{m-2}$, $m\ge2$:
\[
\delta_m\le\frac{4096}{81}\left(\frac{16M_0}{17}-1\right)
\left(\frac{272}{567}\right)^{2(m-2)},\qquad M_0=105\pi^{-4}.
\]
This follows from $x_2=9/16$ and $x_n\ge64/81$ at nonzero odd nodes.
Exactly $\delta_1=p^4$; the limit is an infinite-resolution phenomenon.

\pageheading{FAMILY 3\; /\; TRANSPORT}{Classical operations with arithmetic coordinates}
The depth tilt $x^k$ is a Christoffel transformation. Suppressing $p,\lambda$,
\[
P_m^{(k+1)}(x)=\frac{P_{m+1}^{(k)}(x)
-[P_{m+1}^{(k)}(0)/P_m^{(k)}(0)]P_m^{(k)}(x)}{x}.
\]
The numerator vanishes at zero; all roots are positive, so the displayed
denominator coefficient is nonzero. For shifted Hankel determinants $\tau_m^{(k)}=\det[M_{k+i+j}]_{i,j<m}$,
with $\tau_0^{(k)}=1$,
\[
\tau_m^{(k)}\tau_m^{(k+2)}-(\tau_m^{(k+1)})^2
=\tau_{m+1}^{(k)}\tau_{m-1}^{(k+2)}>0,\qquad m\ge1.
\]
This is the Desnanot--Jacobi identity applied to exact arithmetic moments.
The construction mechanism is classical.

\claim{An endpoint mass is numerically delicate}
For the undeformed depth-zero measure, set $\nu_\theta=\nu+\theta\delta_0$,
$\theta\ge-1$. If $K_m(x,y)=\sum_{j=0}^{m}P_j(x)P_j(y)/h_j$, the Uvarov formula is
\[
P_m^{[\theta]}(x)=P_m(x)-
\frac{\theta P_m(0)}{1+\theta K_{m-1}(0,0)}K_{m-1}(x,0),
\]
with $K_{-1}=0$ for degree zero. The mass at zero is recovered by
\[
1\le\frac1{K_m(0,0)}\le1+(M_0-1)(7/16)^{2m}.
\]
The lower bound comes from the mass-one zero atom. The test polynomial
$(1-x)^m$ gives the upper bound since every other node is at least $9/16$.
At finite degree $K_m(0,0)<1$, because no nonzero polynomial can vanish at
all the other infinitely many nodes. Complete mass removal is valid, while
$1-K_m(0,0)$ becomes very small. This explains the conditioning problem.

\claim{The resolvent is non-D-finite}
The Stieltjes function $S_{p,\lambda}(z)=\int(z-x)^{-1}\dd\nu_{p,\lambda}(x)$
has infinitely many isolated poles at the positive-weight nodes $x_n\to1$.
A solution of a finite rational-coefficient linear ODE has only finitely many
possible finite singular points. The germ at $z=2$ therefore is not D-finite.
This is a general infinite-atomic-resolvent mechanism, specialized here to
arithmetic data; it is not a priority claim.

The full packet sequence also cannot be an affine copy of a continuous Jacobi
sequence: compact moment determinacy would equate a purely atomic measure with
a continuous one. An inherited Laguerre scaling limit remains compatible with
this separation of the complete families.

\pageheading{PRIME SIEVES\; /\; MULTIPLE PARAMETERS}{Overlap is part of the arithmetic}
For a finite set $F$ of distinct primes, define
\[
\nu_{F,\boldsymbol\lambda}
=\sum_{n\ge1}\mu(n)^2n^{-4}
\prod_{p\in F}(1+\lambda_p\mathbf1_{p\mid n})\delta_{x_n},
\quad x_n=(1-n^{-2})^2.
\]
For real strengths, the exact positivity domain is
\[
\boxed{\lambda_p\ge-1\quad\text{for every }p\in F.}
\]
Sufficiency is termwise positivity. If one strength is below $-1$, the
atom $n=p$ isolates its negative factor. The other strengths do not occur
at that atom. At every simultaneous deletion endpoint, primes outside
$F$ still supply infinitely many positive nodes, so all Hankel matrices
remain strictly positive definite.

The exact coordinate formula is
\[
M_j(\boldsymbol\lambda)=\sum_{r=0}^{2j}(-1)^r\binom{2j}{r}
\frac{\zeta(4+2r)}{\zeta(8+4r)}
\prod_{p\in F}\left(1+\frac{\lambda_p}{p^{4+2r}+1}\right).
\]
Coefficients of the monic families lie in
$\mathbb Q(\pi^{-2},\lambda_p:p\in F)$.
For two primes, the overlap term is visible:
\[
M_j=M_j(0,0)+\lambda_pU_{p,j}+\lambda_qU_{q,j}
+\lambda_p\lambda_qU_{pq,j}.
\]
Here $U_{pq,j}>0$ is the $pq$-divisible submeasure moment. The family
does not reduce to adding separate prime responses without their overlap.

\pageheading{PRIME SIEVES\; /\; FINITE TESTS}{Moment positivity is a weaker requirement}
\fig{12_multi_prime_domain}{The exact positive-measure domain and the
order-one moment test have different parameter geometry. Regions passing
a finite test outside the safe orthant still contain negative atoms.}
For primes two and three,
\[
M_0=105\pi^{-4}(1+\lambda_2/17)(1+\lambda_3/82).
\]
Both $(0,0)$ and $(-34,-164)$ have positive total mass. Their midpoint
$(-17,-82)$ has zero mass. Thus the parameter set passing the order-one
positive-definiteness test is nonconvex. This particular midpoint argument
does not establish nonconvexity of the semidefinite test, which permits zero.

\claim{A separate additive family has different geometry}
For weights $1+\sum_{p\in F}\gamma_p\mathbf1_{p\mid n}$, positivity is
equivalent to
\[
1+\sum_{\gamma_p<0}\gamma_p\ge0.
\]
Every prime subset appears at its squarefree product, proving necessity.
Checking each atom proves sufficiency. Its finite Hankel matrices are affine
in $\gamma$, so the semidefinite feasible sets are convex spectrahedra.
The multiplicative family has genuine overlap terms and polynomial parameter
dependence. These are two distinct deformations, with distinct domains.

\pageheading{PRIME SIEVES\; /\; CONSTRUCTIVE WITNESS}{Every negative prime atom has a finite detector}
For an integer $a\ge2$, let
\[
L_a(x)=\prod_{n=1}^{a-1}\frac{x-x_n}{x_a-x_n},\quad
C_a=L_a(1),\quad r_a=\frac{1-x_{a+1}}{1-x_a}<1,
\]
\[
f_{a,m}(x)=L_a(x)\left(\frac{1-x}{1-x_a}\right)^m.
\]
This polynomial kills every lower node, equals one at $x_a$, and has
absolute value at most $C_ar_a^m$ at every higher node.
\fig{13_signed_detector}{A rational detector for $p=3$, $\lambda_3=-1.1$,
$\lambda_2=1$. The first exact bound proving negativity is $m=5$, so a
Hankel matrix of order eight detects the signed defect.}
Set $C_F=\prod_{q\in F}\max(1,|1+\lambda_q|)$. The summable tail gives
\[
\int f_{p,m}^2\dd\nu
\le\frac{1+\lambda_p}{p^4}+
\frac{C_FC_p^2}{3p^3}r_p^{2m}.
\]
Hence the quadratic form is negative whenever
$pC_FC_p^2r_p^{2m}/3<-(1+\lambda_p)$.
For rational strengths this is an exact rational certificate. It needs no
numerical eigenvalue calculation. A depth tilt $x^k$ replaces the negative
atom contribution by $(1+\lambda_p)x_p^k/p^4$.

\pageheading{PRIME SIEVES\; /\; FIBERS}{Real singularities survive multiple deletions}
Fix all strengths except $\lambda_p$, with the others at least $-1$.
The remaining moment matrix has the form $A_m+(1+\lambda_p)B_m$, where
$A_m$ is the $p$-free part and $B_m$ the $p$-divisible part. Both are
positive definite, even at other deletion endpoints: primes outside $F$
supply the first support, and integers $p\ell$ the second.

Thus every one-variable determinant fiber has exact degree $m$ and all
roots real below $-1$. Its nearest margin satisfies
\[
\delta_m=\min_{0\ne f,\deg f<m}
\frac{\int f^2\dd\nu_{p\text{-free}}}
{\int f^2\dd\nu_{p\text{-divisible}}},\qquad\delta_m\downarrow0.
\]
The detector yields, for $C_\kappa=\prod_{q\ne p}\max(1,1+\lambda_q)$,
\[
\delta_{p+m}\le\frac{pC_\kappa C_p^2}{3}r_p^{2m}.
\]
The denominator retains the atom $p$; all lower nodes in the numerator are
killed. This extends the prime-two bound to every prime and to simultaneous
positive or deletion strengths. It does not assert simple roots or global
multivariate real stability.

If at least two positive-domain strengths are nonzero, the complete monic
sequence cannot equal a one-prime sequence on the same variable and nodes.
Compact moment determinacy equates normalized measures; the fixed mass at
zero fixes their scale; atom weights at the individual primes then read off
every strength. An arbitrary affine change of variable is outside this separator.



\MonographPart{III}{Orbit Spectra and Exact Coordinates}{Zero orbits, spectral continuation, and the coefficient language of polyphase transforms.}
\pageheading{FAMILY 4}{One spectral address for each rotation orbit}
For integer $N\ge2$, define the known Mittag--Leffler seed
\[
d_N(w)=\sum_{n\ge0}\frac{(-1)^nw^n}{(Nn)!}=E_{N,1}(-w),
\qquad D_N(z)=d_N(z^N).
\]
Let $w_j$ be the nonzero zeros of $d_N$, counted with multiplicity. Each
corresponds to one rotation orbit of $N$ zeros of $D_N$. Using
$\operatorname{Arg}w\in(-\pi,\pi]$, define
\[
\boxed{\mathcal Z_N(s)=\sum_jw_j^{-s},\qquad\Re s>1/N.}
\]
\fig{06_orbit_zeros}{For $N=4$, the four real/imaginary zeros in each orbit
share one positive coordinate $w=z^4$. The plot shows the first six orbits.
Multiplicity is retained by the definition; simplicity is not assumed in general.}
The bound $|d_N(w)|\le e^{|w|^{1/N}}$ and the finite-order zero-counting theorem
give absolute local uniform convergence in the stated half-plane. Since the
order is below one and $d_N(0)=1$, the Hadamard factorization has genus zero:
\[
d_N(w)=\prod_j(1-w/w_j).
\]
At positive integers, a log derivative gives
\[
-w\frac{d_N'(w)}{d_N(w)}=\sum_{k\ge1}\mathcal Z_N(k)w^k.
\]
If $a_n=(-1)^n/(Nn)!$ and $b_k=\mathcal Z_N(k)$, then
\[
b_k=-ka_k-\sum_{j=1}^{k-1}a_jb_{k-j}.
\]
Thus every positive integer value is rational without numerical zero finding.

\pageheading{FAMILY 4\; /\; REDUCTIONS}{Classical anchors and the open spectral frontier}
The first values are
\[
\begin{aligned}
\mathcal Z_N(1)&=\frac1{N!},\\
\mathcal Z_N(2)&=\frac1{(N!)^2}-\frac2{(2N)!},\\
\mathcal Z_N(3)&=\frac1{(N!)^3}-\frac3{N!(2N)!}+\frac3{(3N)!}.
\end{aligned}
\]
\begin{tabular}{@{}llll@{}}\toprule
$N$ & $\mathcal Z_N(1)$ & $\mathcal Z_N(2)$ & $\mathcal Z_N(3)$\\\midrule
2 & $1/2$ & $1/6$ & $1/15$\\
3 & $1/6$ & $1/40$ & $53/13440$\\
4 & $1/24$ & $17/10080$ & $691/9979200$\\
5 & $1/120$ & $5/72576$ & $997/1743565824$\\
\bottomrule\end{tabular}

\claim{Two exact Hurwitz-zeta reductions}
For $N=2$, $d_2(w)=\cos\sqrt w$, whose zeros are
$w_j=\pi^2(j+1/2)^2$, $j\ge0$. Therefore
\[
\mathcal Z_2(s)=\pi^{-2s}\zeta(2s,1/2).
\]
For $N=4$,
\[
D_4(z)=\cos(z/\sqrt2)\cosh(z/\sqrt2),
\]
and the rotation quotient has one zero
$w_j=(\sqrt2\pi)^4(j+1/2)^4$ per orbit. Hence
\[
\mathcal Z_4(s)=(\sqrt2\pi)^{-4s}\zeta(4s,1/2).
\]
These anchors are classical reductions. The general spectral-zeta mechanism
also has substantial antecedents, including hypergeometric zero zeta functions.

\claim{Global meromorphic continuation}
The original half-plane definition now has a meromorphic extension to the
whole complex plane. The next pages prove its sole pole, special values,
and regularized product using a branch-safe Mellin representation.
These are worked spectral results from established analytic mechanisms.
Residue-weighted sums involving derivatives at zeros remain a separate
question, especially when multiplicities occur.

\pageheading{ORBIT SPECTRA\; /\; CONTINUATION THEOREM}{The continuation gap is closed}
For every integer $N\ge2$, the orbit zeta has a meromorphic continuation to
the whole complex plane with exactly one pole:
\[
\boxed{s=1/N,\qquad
\operatorname{Res}_{s=1/N}\mathcal Z_N(s)=\frac{\sin(\pi/N)}{N\pi}.}
\]
Its exact absolute-convergence abscissa is $1/N$. At nonpositive integers,
\[
\boxed{\mathcal Z_N(-m)=0\quad(m\ge0),\qquad
\mathcal Z_N'(0)=-\log N.}
\]
The zeta-regularized product of the orbit coordinates is consequently
$\exp[-\mathcal Z_N'(0)]=N$.
\fig{15_spectral_poles}{Pole and zero locations follow from the proved
continuation. The second panel shows the shifted extension, where a pole
ladder appears. These are exact structural locations rather than fitted roots.}
The continuation uses classical Mellin regularization. The contribution here is a complete local application with multiplicities,
principal branches, and all convergence hypotheses checked.

\pageheading{ORBIT SPECTRA\; /\; LARGE RADIUS}{A finite exponential sum controls the tail}
Put $E_N(t)=d_N(-t)=\sum_{n\ge0}t^n/(Nn)!$ and
$\omega=e^{2\pi i/N}$. For $x\ge0$,
\[
E_N(x^N)=\frac1N\sum_{j=0}^{N-1}e^{\omega^jx}.
\]
Root filtering proves this identity term by term: only powers divisible by
$N$ remain. The positive real exponential is dominant, so
\[
\log E_N(t)=t^{1/N}-\log N+O(e^{-\kappa_N t^{1/N}}),
\quad \kappa_N=1-\cos(2\pi/N)>0,
\]
with corresponding differentiated estimates on the positive ray.
\fig{14_spectral_tail}{Computed logarithmic derivatives of the finite
exponential sum approach one in the root-radius variable. The decaying
remainder supplies the Mellin continuation's large-radius input.}
\claim{Why complex zero arguments do not break the integral}
In a sufficiently narrow sector around negative $w$, choose the root
$x=(-w)^{1/N}$ near the positive real ray. The first exponential dominates
the other $N-1$ terms by a uniform real-part gap at large radius, so this
sector is eventually zero free. The finitely many remaining zeros avoid
the negative real ray because $d_N(-t)>0$ for $t\ge0$.
Thus all zero arguments have a uniform gap from the branch cut.
In particular $|t+w_j|\ge c_N(t+|w_j|)$ for $t\ge0$.

\pageheading{ORBIT SPECTRA\; /\; MELLIN PROOF}{The two ends have different jobs}
For $1/N<\Re s<1$, the genus-zero log derivative and the angular gap give
absolute interchange in
\[
\mathcal Z_N(s)=\frac{\sin\pi s}{\pi}
\int_0^\infty t^{-s}\frac{d}{dt}\log E_N(t)\dd t.
\]
Each term uses the principal-branch beta integral
$\int_0^\infty t^{-s}(t+w)^{-1}\dd t=\pi w^{-s}/\sin\pi s$.
The denominator estimate bounds the absolute integral by a constant times
$|w|^{-\Re s}$, whose sum converges in this strip.

Near zero write $f(t)=(\log E_N)'(t)=\sum_{k\ge1}\alpha_kt^{k-1}$,
where $\alpha_k=(-1)^{k+1}\mathcal Z_N(k)$.
For any integer $K\ge1$, subtract $K$ terms on $[0,1]$:
\[
\begin{aligned}
I_K(s)={}&\int_0^1t^{-s}\left[f(t)-\sum_{k=1}^K\alpha_kt^{k-1}\right]\dd t
+\sum_{k=1}^K\frac{\alpha_k}{k-s}\\
&+\int_1^\infty t^{-s}\left[f(t)-\frac1N t^{1/N-1}\right]\dd t
+\frac{1}{N(s-1/N)}.
\end{aligned}
\]
The first integral is holomorphic for $\Re s<K+1$; the second is entire
because its remainder decays exponentially in $t^{1/N}$.
The factor $\sin\pi s$ cancels the small-radius poles at positive integers.
Letting $K$ increase covers the plane and leaves only the pole at $1/N$.
At nonpositive integers $I_K$ is regular, giving the stated zero values.

At zero the finite parts combine to
\[
I_K(0)=\lim_{T\to\infty}[\log E_N(T)-T^{1/N}]=-\log N.
\]
Since $(\sin\pi s/\pi)'_{s=0}=1$, this is $\mathcal Z_N'(0)$.
The pole also proves the exact abscissa: absolute convergence at any smaller
real exponent would give a holomorphic series across $1/N$, a contradiction.

\pageheading{ORBIT SPECTRA\; /\; SHIFTS}{A Hurwitz-type deformation and its determinant}
For real $a\ge0$, define initially
\[
\mathcal Z_N(s;a)=\sum_j(w_j+a)^{-s},\qquad\Re s>1/N.
\]
The rightward shift preserves avoidance of the negative branch cut.
Normalize the product with $E_N(a)$ and use
$\log E_N(a+t)$ in the Mellin argument.
For $a>0$ there are genuine simple poles at $s=1/N-k$, $k\ge0$:
\[
\boxed{\operatorname{Res}_{s=1/N-k}\mathcal Z_N(s;a)
=\frac{\sin(\pi/N)}{N\pi}\frac{(1-1/N)_k}{k!}a^k.}
\]
Indeed the large-radius derivative expands as
\[
\frac1N(a+t)^{1/N-1}
=\frac1N\sum_{k\ge0}\frac{(-1)^k(1-1/N)_k}{k!}a^kt^{1/N-1-k}
\]
as an asymptotic series. The alternating sign is canceled by
$\sin(\pi(1/N-k))=(-1)^k\sin(\pi/N)$.
At $a=0$ all secondary residues vanish.

The negative-integer values remain zero, and
\[
\mathcal Z_N'(0;a)=-\log[N E_N(a)],\qquad
\boxed{\prod_j^{\rm reg}(w_j+a)=N E_{N,1}(a).}
\]
This uses the large-radius finite part $-\log N$ and the normalization
term $-\log E_N(a)$. It is a zeta-regularized product, not an ordinary
convergent product of the coordinates.

For $N=2$, the result is $2\cosh\sqrt a$, a classical half-integer
spectral determinant. The four shifted numerical benchmarks reproduce this
anchor. Functional equations and derivative-weighted zero sums remain
separate questions; the continuation proof does not settle them.

\pageheading{EXACT LABORATORY\; /\; COEFFICIENTS}{Integer coordinates for the nonlinear family}
Normalize the Taylor coefficients by
\[
A_n(q,r)=(2n+1)!c_n(q,r),\qquad A_0=1.
\]
The recurrence becomes
\[
\boxed{A_n=\sum_{i=0}^{n-1}\binom{2n}{2i+1}
q^{2i+1}r^{2(n-1-i)+1}A_iA_{n-1-i},\quad n\ge1.}
\]
Therefore $A_n\in\mathbb Z[q,r]$ has nonnegative coefficients and is
symmetric in its two variables. The first coordinates are
\[
\begin{aligned}
A_1&=2qr,\\
A_2&=8q^2r^2(q^2+r^2),\\
A_3&=q^3r^3[48(q^2+r^2)(q^4+r^4)+80q^2r^2].
\end{aligned}
\]
At $(q,r)=(1,1)$ these are the tangent numbers
$1,2,16,272,7936,\ldots$, the known OEIS A000182 anchor.
At $(1,2)$ the exact sequence begins
$1,4,160,35200,38041600,\ldots$.
\smallfig{17_tree_addresses}{Ordered trees retain left/right path information.
The two displayed shapes contribute different monomials to the same
factorial-normalized coefficient. Their labels are exact tree data.}
At rational scales, an integer denominator clearing is available. If $D$
clears every contraction denominator and $H$ clears every bank weight,
\[
H^nD^{n(n+1)}A_n\in\mathbb Z.
\]
The maximum root-path exponent sum over a full binary tree is $n(n+1)$.
The clearing provides exact sequence exports without implying a new OEIS entry.

\pageheading{EXACT LABORATORY\; /\; REGIMES}{A formal sequence is not always an analytic germ}
\smallfig[0.68]{18_parameter_regimes}{The full nonnegative real parameter map separates
finite-order entire functions, infinite-order entire edges, tangent, and
zero-radius formal series. Zero scales give the linear function.}
If $q>0$, $r>1$, the recurrence includes its $i=0$ term and gives
\[
c_n\ge\frac{q^nr^{n^2}}{3\cdot5\cdots(2n+1)}.
\]
The root test proves zero radius. By symmetry this covers either scale above
one, provided the other is positive. The exported positive integer asymmetric
rows belong to this formal regime.

If $0<q<1$, $r=1$, the equation is linear in the uncontracted value:
$T'=1+T(qz)T(z)$. A known coefficient on $|z|<R/q$ gives a linear ODE
there, so disk bootstrapping proves entire extension. On the positive ray,
$(\log T)'=T(qR)+1/T(R)\ge qR$, giving an exponential lower bound for
$T$. Substitution back gives, for $R\ge2/q$,
\[
\log T(R)\ge\frac R2\exp(q^3R^2/8-q/2).
\]
Thus the entire order is infinite. The corner $(1,1)$ is tangent with poles.
These exact boundaries protect the analytic interpretation of the catalog.

\pageheading{EXACT LABORATORY\; /\; PROJECTORS}{Polyphase sections preserve the full sequence}
For an ordinary generating function $G(z)=\sum_{m\ge0}a_mz^m$ and
integer $b\ge1$, write
\[
G(z)=\sum_{j=0}^{b-1}z^jG_j(z^b),\qquad
G_j(t)=\sum_{n\ge0}a_{bn+j}t^n.
\]
The root projector is
\[
P_{b,j}G(z)=\frac1b\sum_{k=0}^{b-1}\omega^{-jk}G(\omega^kz)
=z^jG_j(z^b),\qquad\omega=e^{2\pi i/b}.
\]
The same projector works on an EGF, but the section retains its original
factorials: $\sum_na_{bn+j}z^{bn+j}/(bn+j)!$.
Removing those denominators changes the representation and requires an
explicit transport. Root filtering does not do it automatically.

For causal FIR taps $h_\ell$, $b_m=\sum_\ell h_\ell a_{m-\ell}$ with
zero extension to negative indices, split $A(z)$ and $B(z)$ into $b$ phases.
Then $B_i(t)=\sum_jH_{ij}(t)A_j(t)$, where
\[
\boxed{H_{ij}(t)=\sum_{\substack{\ell\ge0\\j+\ell\equiv i\pmod b}}
h_\ell t^{(j+\ell-i)/b}.}
\]
The exponent is the nonnegative carry to the next phase block. This is an
exact coefficient law, with no approximation or loss of terms.

\pageheading{EXACT LABORATORY\; /\; TAPS}{The carry matrix is part of the filter}
\smallfig{16_polyphase}{A three-phase filter with taps $[1,2,-1]$.
The matrix carries terms across block boundaries. Interleaving its output
reproduces the direct scalar FIR calculation exactly.}
With these taps and $b=3$,
\[
H(t)=\begin{pmatrix}
1&-t&2t\\2&1&-t\\-1&2&1
\end{pmatrix}.
\]
The entries arise from the same congruence condition. For instance an input
in phase two shifted by one lands in phase zero of the next block, contributing
$2t$ to $H_{0,2}$.

\claim{Three operations that must be kept distinct}
For an EGF $E_a(z)=\sum_ma_mz^m/m!$, differentiation transports future
coefficients: $E_a^{(k)}=\sum_ma_{m+k}z^m/m!$.
Causal ordinary taps instead act by repeated zero-constant integration:
\[
E_b=\sum_\ell h_\ell I^\ell E_a.
\]
Multiplication of two EGFs gives Hurwitz convolution,
$b_m=\sum_{k=0}^m\binom mk a_kh_{m-k}$.
It is not ordinary sequence FIR convolution. Intrinsic nonlinear feedback
taps alter the defining family and are a fourth operation with their own
coefficient recurrence.

\pageheading{EXACT LABORATORY\; /\; REUSABLE CODE}{A catalog with explicit provenance}
The standard-library \texttt{family\_lab.py} exposes exact coefficient
polynomials, normalized sequences, positive-bank coefficients, rational
denominator clearing, root sections, Hurwitz products, ordinary FIR taps,
and polyphase filter matrices. The catalog contains 18 rows through index 12.

Every row records parameters, coefficient normalization, arithmetic regime,
and whether its entries are integers or rationals. Eight positive integer
asymmetric specializations are labeled zero-radius formal sequences. Entire
interior examples and an infinite-order edge are labeled separately.

\begin{tabular}{@{}p{0.27\linewidth}p{0.35\linewidth}p{0.30\linewidth}@{}}
\toprule
Example & First normalized coordinates & Interpretation\\\midrule
$(1,1)$ & $1,2,16,272,7936$ & Known tangent anchor\\
$(1,2)$ & $1,4,160,35200,38041600$ & Integer, formal radius zero\\
$(1/2,1/2)$ & $1,1/2,1/4,1/8,1/16$ & Scaled sinh, entire order one\\
$(1/3,2/3)$ & $1,4/9,160/729,18560/177147,\ldots$ & Asymmetric entire order one\\
\bottomrule
\end{tabular}

\claim{Implementation boundaries are enforced}
Phase residues and bases must be integers, and a phase bank cannot be empty.
The independent review found that fractional residues and zero phases were
previously accepted. Both domain-validation defects were fixed and covered
by regression assertions before publication.

The 8,038 exact API checks compare separate constructions: Picard integration,
sine/cosine tangent division, root filters in a cyclotomic quotient,
direct and phase-matrix taps, and coefficient-level antiderivatives.
They verify finite implementation ranges. The algebraic laws above have
separate proofs in Appendix E.

The catalog is ready for comparison with OEIS, with known anchors identified.
It makes no claim that an unmatched prefix establishes a new sequence, and
no sequence is represented here as a submitted or independently assigned OEIS entry.



\MonographPart{IV}{Coupled Systems and Finite Resolution}{How interacting channels alter global growth, and what a finite observer can certify.}
\pageheading{COUPLED FAMILIES}{Coupled growth and the cost of resolution}
The new family consists of two entire functions whose derivatives depend on
contracted values of both components. Its growth exponent is the unique
positive number $\rho$ for which
\[
\operatorname{spr}\begin{pmatrix}q^\rho&r^\rho\\a^\rho&b^\rho\end{pmatrix}=1,
\qquad 0<q,r,a,b<1.
\]
Both functions have positive lower type and finite upper type at this order.
The positive eigenvector fixes the proportions of both type vectors. The
existence of an ordinary type limit remains open.

\claim{Local coefficients reveal the coupling}
Two explicit parameter surfaces describe exactly when one normalized component
is an amplitude and argument rescaling of the other. Outside these surfaces,
the coefficients through degree seven recover all four ordered scales. This
provides an exact local inverse for the labeled pair.

\claim{Three ways finite observations can miss structure}
The resolution atlas compares a weak tap, a large excluded prime, and a negative
sieve atom. Each example links a structural change to a concrete observation
scale and a proved certificate.
\begin{center}
\renewcommand{\arraystretch}{1.35}
\begin{tabular}{@{}p{.22\linewidth}p{.31\linewidth}p{.39\linewidth}@{}}
\toprule
Example & Observation & What the certificate gives\\\midrule
Weak feedback tap & Values on a complex disk & An explicit compact error and a sufficient radius for separation\\
Excluded prime & Values and derivatives & Exact disk bounds and a negative-depth visibility bracket\\
Negative sieve atom & All polynomial tests up to a degree & Positive norms through degree four and a negative degree-seven witness\\
\bottomrule
\end{tabular}
\end{center}

The computed plates are presented in the following chapters.
Complete derivations accompany the full monograph, and the source archive
contains the alternative growth proof, exact laboratory, and review records.
The six-source literature comparison identifies classical antecedents.
Historical priority of the specific coupled results remains unresolved.

\claim{How to read the figures}
An exact rational bound supplies a certificate. A decimal display illustrates
its size. Numerical curves retain their stated truncation and precision, and
the universal growth and recovery statements rest on their proofs.

% Family 5 visual chapter. Included by the monograph and the visual guide.
% Exact proof sources: research_v11/coupled_growth_proof.md,
% coupled_identifiability.md, and coupled_structure.md.

\pageheading{FAMILY 5}{Two functions exchange growth}
For four real contractions $0<q,r,a,b<1$, define the normalized pair
\[
\boxed{\begin{aligned}
U'(z)&=1+U(qz)V(rz),\\
V'(z)&=1+U(az)V(bz),\qquad U(0)=V(0)=0.
\end{aligned}}
\]
Each equation combines one value of $U$ with one value of $V$. The four
scales specify where that exchange occurs. Both functions begin with the
same unit slope, while later coefficients retain the two equations' different
scale addresses.

\vfig[0.96]{coupled_01_matrix}{The worked pair
$(q,r,a,b)=(1/5,2/5,3/5,7/10)$ has exact order one. The two displayed
identities give the positive growth vector $(1/2,1)^T$.}

There is a unique entire solution pair. Both functions are odd, and every
odd Taylor coefficient is positive. The common order is governed by the
positive matrix
\[
P(s)=\begin{pmatrix}q^s&r^s\\a^s&b^s\end{pmatrix},
\qquad \operatorname{spr}P(\rho)=1.
\]
The eigenvector supplies the relative weights that a row sum misses. In
the example, the row sums at $s=1$ are $3/5$ and $13/10$. The coupled
system nevertheless has order exactly one in both components.

\pageheading{FAMILY 5 / CONSTRUCTION}{The coefficients preserve both addresses}
Write $U(z)=\sum_{n\ge0}u_nz^{2n+1}$ and
$V(z)=\sum_{n\ge0}v_nz^{2n+1}$, with $u_0=v_0=1$.
The first equation gives the exact recursion
\[
u_n=\frac1{2n+1}\sum_{i+j=n-1}
q^{2i+1}r^{2j+1}u_i v_j,\qquad n\ge1.
\]
The second replaces $(q,r)$ by $(a,b)$ while retaining the same products
$u_i v_j$. Thus each new coefficient uses the shared history of the pair.

\vfig[0.97]{coupled_02_functions}{For
$(q,r,a,b)=(1/3,2/3,1/2,3/4)$, the components separate despite their
common initial slope. These numerical curves use odd indices through 240.
The two truncations at indices 120 and 240 agree at $x=8$ to displayed
double precision.}

The first distinction already appears at degree three:
\[
U(z)=z+\frac{qr}{3}z^3+O(z^5),\qquad
V(z)=z+\frac{ab}{3}z^3+O(z^5).
\]
Induction gives $0<u_n,v_n\le2^{-n}$, which supplies a convergent germ.
If both functions are known on a disk of radius $R$, their integral
equations extend them to radius $R/Q$, where $Q=\max(q,r,a,b)<1$.
Repeating this step reaches every finite disk.

The numerical agreement above checks the plotted approximation. The later
certificate page supplies an exact bound on the omitted tail.

\pageheading{FAMILY 5 / GROWTH}{One order, two exact type proportions}
The entire functions have the same order $\rho$, with positive lower type
and finite upper type. Positivity places the maximum modulus on the positive
real axis, so $M_U(R)=U(R)$ and $M_V(R)=V(R)$.

\vfig[0.96]{coupled_03_growth}{The order-one example has different component
amplitudes. The right panel shows finite-radius slopes near one. The curves
use odd indices through 800. Agreement with index 400 is a numerical
diagnostic, and the order itself follows from the analytic proof.}

For $F_1=U$ and $F_2=V$, define
\[
\underline\tau_i=\liminf_{R\to\infty}\frac{\log M_{F_i}(R)}{R^\rho},
\qquad
\overline\tau_i=\limsup_{R\to\infty}\frac{\log M_{F_i}(R)}{R^\rho}.
\]
If $P(\rho)w=w$ with $w>0$, then both type vectors have precisely these
Perron proportions:
\[
\boxed{\underline{\boldsymbol\tau}=\underline\tau\,w,
\qquad\overline{\boldsymbol\tau}=\overline\tau\,w,
\qquad0<\underline\tau\le\overline\tau<\infty.}
\]
For the worked pair, both the lower type and the upper type of $V$ are
twice the corresponding type of $U$. Either both ordinary type limits
exist or neither does. Equality of the lower and upper types remains open.
The finite-radius slope picture does not decide that question.

\pageheading{FAMILY 5 / PARAMETER SPACE}{Growth follows a spectral surface}
The order is the unique positive solution of
\[
\boxed{(1-q^\rho)(1-b^\rho)=(ar)^\rho.}
\]
The left side increases from zero to one as $\rho$ increases. The right
side decreases from one to zero. Their unique crossing determines the
growth exponent without fitting a truncated Taylor curve.

\vfig[0.80]{coupled_04_landscape}{Numerical solutions of the proved
characteristic equation, with $r=2/5$ and $a=3/5$. The highlighted pair
$(q,b)=(1/5,7/10)$ has order exactly one by rational arithmetic.}

The order depends on the self scales $q,b$ and on the product $ar$ of
the cross scales. Holding $q,b,ar$ fixed preserves the order, while
changing the cross scales separately generally changes the functions and
their type proportions:
\[
\frac{w_1}{w_2}=\frac{r^\rho}{1-q^\rho}
=\frac{1-b^\rho}{a^\rho}.
\]
Consequently, equal order does not determine the family member. It fixes
one growth exponent within a larger parameter space. The coefficients
carry enough additional information to recover all four ordered scales
away from the reduction surfaces described below.

\pageheading{FAMILY 5 / PRECISION}{A coefficient list with a certified tail}
A finite Taylor list becomes an exact approximation certificate when it
is paired with a bound on a larger disk. Let $F$ be either component,
and let $F_N$ retain the odd indices $0,\ldots,N$.

\vfig[0.95]{coupled_05_certification}{Each curve is computed from an exact
rational upper bound for the same worked pair. Increasing the coefficient
count certifies a smaller uniform error on the stated complex disk.
The logarithmic plotting coordinates are rounded.}

Put $Q=\max(q,r,a,b)$. One elementary majorant starts with
$B(x)=x/(1-x^2)$ for $0\le x<1$ and works outward through the finite
contraction ladder using
\[
B(x)=x\bigl(1+B(Qx)^2\bigr).
\]
The integral equations show that this bounds both components on the disk
$|z|\le x$. For $S>R\ge0$, positivity then gives
\[
\boxed{\max_{|z|\le R}|F(z)-F_N(z)|
\le B(S)\left(\frac RS\right)^{2N+3}.}
\]
The exponent is the first omitted odd degree. Every later term has an
even smaller ratio between its values at radii $R$ and $S$.
For the worked pair, $N=30$, $R=1$ and $S=2$ give an error below
$1.600\times10^{-13}$ for both functions. The laboratory returns
rational interval endpoints, so the certificate survives decimal rounding.

\pageheading{FAMILY 5 / SEQUENCES}{The phase calculus transports the pair}
Use the derivative normalization
$A_n=(2n+1)!u_n$, $B_n=(2n+1)!v_n$. Its recurrence has integer
binomial weights:
\[
A_n=\sum_{i+j=n-1}\binom{2n}{2i+1}
q^{2i+1}r^{2j+1}A_iB_j,\qquad n\ge1.
\]
The second component again replaces $(q,r)$ by $(a,b)$.

\vfig[0.91]{coupled_06_polyphase}{The full derivative array of $V$ is split
into three residue lanes. Applying taps $(1,-2,1)$ with the correct block
carries, then interleaving, agrees exactly with direct convolution.}

These derivative arrays are polynomials with nonnegative integer
coefficients in the four scale variables. Their total degree is at most
$n(n+1)$. If every parameter denominator divides $D_0$, then
\[
D_0^{n(n+1)}A_n,\qquad D_0^{n(n+1)}B_n
\quad\hbox{are integers}.
\]
This uniform clearing factor makes exact rational families available to
integer sequence experiments. It need not be the smallest possible factor.

Ordinary taps on a derivative array correspond to repeated integration
shifts of its exponential generating function. Multiplication of two
exponential generating functions instead requires binomial convolution.
The catalog preserves that distinction and marks generic integer and
signed scale examples as formal series outside the proved analytic domain.

\pageheading{FAMILY 5 / REDUCTION}{Two surfaces where one component rescales to the other}
Set $p=qr$, $k=ab$, and use the ordinary Taylor coefficients $u_2,v_2$
at degree five. The exact separation quantity factors as
\[
\boxed{D=k^2u_2-p^2v_2
=\frac{pk ar}{15}(q^2-b^2)(qb-ar).}
\]

\vfig[0.88]{coupled_07_reduction_geometry}{The two reduction surfaces are
$q=b$ and $qb=ar$. Their union is exactly the zero set of $D$ in the
positive contraction domain. Outside this union, the components cannot
be obtained from each other by amplitude and argument rescaling.}

There exist nonzero complex constants $A,B$ with $V(z)=A\,U(Bz)$
if and only if $D=0$. On either surface, put $s=\sqrt{k/p}>0$.
The scalar family is defined by
$T_{x,y}'(z)=1+T_{x,y}(xz)T_{x,y}(yz)$, $T_{x,y}(0)=0$.
It gives the explicit reduction
\[
U(z)=\sqrt{s}\,T_{q,sr}(z/\sqrt{s}),\qquad
V(z)=s^{-1/2}T_{q,sr}(\sqrt{s}\,z).
\]
Both effective scales lie in $(0,1)$. Equality $U=V$ is more restrictive:
it occurs exactly when $\{q,r\}=\{a,b\}$.

This classifies relations between the normalized components. Separately,
no invertible constant linear coordinate change decouples the defining
operator into two autonomous scalar equations on every analytic input.
That operator statement does not exclude nonlinear changes of variables
or other special-function identities.

\pageheading{FAMILY 5 / RECOVERY}{Degree seven identifies the four scales}
Away from $D=0$, the six coefficients
$(u_1,v_1,u_2,v_2,u_3,v_3)$ uniquely recover the ordered parameters
$(q,r,a,b)$ in $(0,1)^4$. These are the coefficients at degrees three,
five and seven of the two labeled functions.

\vfig[0.92]{coupled_08_inverse}{Cubic terms recover the two scale products.
Fifth-degree terms leave at most two candidates for each remaining scale.
The seventh-degree terms select the unique admissible candidates when
$D\ne0$. Repeated quadratic roots are already unique.}

First $p=3u_1$ and $k=3v_1$. The possible values $x=q^2$ and $y=a^2$
satisfy
\[
p x^2-\frac{15u_2}{p}x+kp^2=0,\qquad
p y^2-\frac{15v_2}{k}y+k^3=0.
\]
For two distinct positive candidates $x,x'$, the predicted seventh
coefficients differ by
\[
u_3(x)-u_3(x')=\frac{pD}{7k^2}(x^2-x'^2)\ne0.
\]
The second component supplies the analogous separation for $y$.
Finally take the positive square roots and recover $r=p/q$, $b=k/a$.

This is an exact-data inverse theorem. Near a reduction surface, small
$D$ can make candidate predictions close, so a uniform stable inverse
does not follow. The degree-seven recovery and the reduction classification
are proved statements. Historical priority of these particular results
is addressed separately by the bounded literature comparison.

\pageheading{RESOLUTION ATLAS\; /\; WEAK TAPS}{A small change can alter the global order}
A feedback channel may be almost invisible on a fixed disk and still determine
the eventual growth of the function. Consider the normalized entire solution
\[
F_\varepsilon'(z)=1+(1-\varepsilon)F_\varepsilon(z/4)^2
 +\varepsilon F_\varepsilon(3z/4)^2,
\qquad F_\varepsilon(0)=0.
\]
Here $0\le\varepsilon\le1$. At $\varepsilon=0$, its entire order is $1/2$.
At every positive weight, the order becomes
$\rho=\log2/\log(4/3)\approx2.40942084$.

\vfig{resolution_01_weak_tap}{Left: divide the compact difference by the tap
weight to compare its shape. Right: the proved order jumps at zero. The dashed
curve is an analytic envelope. The function curves use 100-digit arithmetic
and have a normalized omitted-series bound below $5.26\times10^{-23}$ on
$0\le R\le1$. Floating rounding is separate from that tail bound.}

\claim{The unit disk comes with an exact error certificate}
Let $f_{3/4}=F_1$, with $f_{3/4}(z)=\sum_{n\ge0}d_nz^{2n+1}$.
Positive coefficients give
\[
\sup_{|z|\le R}|F_\varepsilon(z)-F_0(z)|
\le\varepsilon\sum_{n\ge1}nd_nR^{2n+1}.
\]
A rational coefficient sum and a positive-tail estimate prove that, on the
unit disk, the multiplier is less than $0.241890267$.
For $\varepsilon=10^{-12}$, the entire disk therefore changes by less than
$2.41890267\times10^{-13}$.

The observation window and the global invariant answer different questions.
Exact measurements can distinguish the functions. With an absolute error
allowance at least half the displayed bound, midpoint data can fit both.
The order concerns the limit through arbitrarily large radii, where the
active faster channel eventually matters.

\pageheading{RESOLUTION ATLAS\; /\; OBSERVATION RADIUS}{How far is far enough to guarantee separation?}
The growth theorem supplies positive lower type for the weak-tap function
and finite upper type for the slow function. For $0<\varepsilon\le1$ and $R\ge e$,
\[
\log F_\varepsilon(R)\ge c_\varepsilon R^\rho,
\qquad \log F_0(R)\le C_0\sqrt R,
\]
\[
c_\varepsilon=\frac{\log2}{(4/\varepsilon+8/3)^\rho},
\qquad C_0=e^{-1/2}(1+22\log2).
\]

\vfig{resolution_02_growth_resolution}{The left curve is a sufficient
radius obtained from conservative comparison constants. The right panel
shows why the lower bound for the perturbed function eventually exceeds the
upper bound for the slow one. These are bounds on logarithmic growth, rather
than numerical evaluations at the large plotted radii.}

\claim{An explicit sufficient radius}
Put
\[
R_2(\varepsilon)=\max\!\left\{
 e,\left(\frac{2C_0}{c_\varepsilon}\right)^{1/(\rho-1/2)},
 \left(\frac{2\log2}{c_\varepsilon}\right)^{1/\rho}\right\}.
\]
Then $F_\varepsilon(R)\ge2F_0(R)$ for every $R\ge R_2(\varepsilon)$.
The formula makes $c_\varepsilon R^\rho$ at least both
$2C_0\sqrt R$ and $2\log2$, and therefore at least their half-sum.
Exponentiating gives the factor of two.

For $\varepsilon=10^{-6}$, the formula is approximately
$1.237\times10^9$. This large radius is a guaranteed sufficient choice.
It is not the first radius at which a useful difference appears, and it does
not establish failure of every observation at smaller radii.
The exact formula carries the guarantee. Rounded figures illustrate its scale.

\pageheading{RESOLUTION ATLAS\; /\; ARITHMETIC PRECISION}{A large excluded prime hides in a fixed window}
Use the squarefree arithmetic family
\[
A_Q(s)=\sum_{\substack{n\ge2\\\mu(n)^2=1\\(n,Q)=1}}
 n^{-4}(1-n^{-2})^{2s},
\qquad D_p(s)=A_1(s)-A_p(s),
\]
where $p$ is prime. The difference contains exactly the squarefree atoms
divisible by $p$. Its largest frequency is
$\lambda_p=-2\log(1-p^{-2})$, and its total mass is
$D(4)/(p^4+1)$, with $D(4)=105/\pi^4$.

\vfig{resolution_03_prime_precision}{Left: analytic bounds for the
function and three derivatives on $|s|\le5$. The connecting curves show scale,
while the arithmetic statement applies at primes. Right: rational bounds
for three particular excluded primes. The bar height is the negative
base-ten logarithm of the bound, not a count of recovered digits.}

\claim{The full complex disk has a rational certificate}
For integer $k\ge0$,
\[
\sup_{|s|\le5}|D_p^{(k)}(s)|
<\frac{4(1-p^{-2})^{-10}}{3(p^4+1)}
 \left(\frac{2}{p^2-1}\right)^k.
\]
The inequality uses $D(4)<4/3$ and
$\lambda_p\le2/(p^2-1)$. Every quantity on the right is rational.
For $k=0$, upward decimal bounds are $1.283\times10^{-8}$ at $p=101$,
$1.287\times10^{-12}$ at $p=1009$, and $1.330\times10^{-16}$ at $p=10007$.

The excluded prime changes the exact inverse of the function.
Nevertheless, its effect becomes arbitrarily small on every fixed disk,
including in every fixed derivative order. A bounded error model therefore
needs an arithmetic parameter bound if it is to promise uniform recovery.

\pageheading{RESOLUTION ATLAS\; /\; NEGATIVE DEPTH}{Seeing the prime can make subtraction harder}
At negative real depth, the difference grows. The prime atom alone supplies
a lower bound, while the entire removed shelf supplies an upper bound:
\[
p^{-4}e^{\lambda_pT}\le D_p(-T)
\le\frac{D(4)}{p^4+1}e^{\lambda_pT},\qquad T\ge0.
\]
Choose a tolerance $\eta>D(4)/(p^4+1)$. Its unique crossing occurs between
\[
T_- =\frac{\log\!\bigl(\eta(p^4+1)/D(4)\bigr)}{\lambda_p},
\qquad
T_+ =\frac{\log(\eta p^4)}{\lambda_p}.
\]

\vfig{resolution_04_prime_depth}{At tolerance $10^{-6}$, the left panel
shows the analytic crossing brackets after rescaling each row by $T_+$.
The right panel shows a lower bound on the decimal cancellation scale when
subtracting full functions. Displayed endpoints and labels are ordinary
floating evaluations of exact formulas.}

For $p=101$, the difference is no larger than $10^{-6}$ through
$T_-=23307.789\ldots$, and at least that large by
$T_+=23690.516\ldots$. The decimals illustrate the formulas and are not
rounded interval certificates.

\claim{The direct positive difference and subtraction have different costs}
The $n=2$ atom in the full function proves
\[
\frac{A_1(-T)}{D_p(-T)}
\ge\frac{p^4+1}{16D(4)}e^{(\lambda_2-\lambda_p)T}.
\]
At $T_+$ for $p=101$, the base-ten logarithm of this lower bound is about
5924. Subtracting independently approximated full functions can therefore
amplify relative errors enormously. Evaluating the known positive difference
series is a different procedure. These conditioning figures impose no
universal precision requirement on every possible recovery algorithm.

\pageheading{RESOLUTION ATLAS\; /\; NEGATIVE ATOMS}{A degree-seven witness exposes a signed defect}
Let $x_n=(1-n^{-2})^2$, and change the prime-two sieve strength to
$\lambda_2=-1-\delta$. For a real polynomial $f$, its squared norm is
\[
Q_\delta(f)=\sum_{\substack{n\ge1\\\mu(n)^2=1\\n\text{ odd}}}
 \frac{f(x_n)^2}{n^4}
-\delta\sum_{\substack{n\ge1\\\mu(n)^2=1\\n\text{ even}}}
 \frac{f(x_n)^2}{n^4}.
\]
Every even squarefree atom has negative mass. At $\delta=10^{-12}$,
the polynomial below certifies that this is detectable at finite degree.

\vfig{resolution_05_signed_detector}{The detector vanishes on the first
four positive shelves and equals one at the negative atom $n=2$.
The positive bar is an upper bound for the entire remaining positive tail.
The total upper bound is negative, so no omitted atom can reverse the
conclusion.}

\claim{The certificate uses rational arithmetic throughout}
Take $S=\{1,3,5,7\}$ and define
\[
P(x)=\prod_{n\in S}\frac{x-x_n}{x_2-x_n}
 \left(\frac{1-x}{1-x_2}\right)^3.
\]
Its degree is seven. All remaining positive nodes begin at $n=11$.
With
$C=756482048/20777538825$ and $r=3856/102487$, a complete tail bound gives
\[
Q_{10^{-12}}(P)\le-\frac{10^{-12}}{16}
 +\frac{C^2r^6}{3\cdot7^3}
 <-5.8845708899\times10^{-14}.
\]
The inequality is certified as a comparison of exact rational numbers.
The first four zeros remove the dominant positive contributions.
The factor $(1-x)^3$ suppresses the remaining cluster near one.
The construction exhibits a detecting degree, without asserting that degree
seven is the smallest possible choice.

\pageheading{RESOLUTION ATLAS\; /\; ALL LOW DEGREES}{Passing every coarse test can still miss the defect}
The previous witness would be less informative without a converse guarantee
at lower resolution. For the same defect $\delta=10^{-12}$, every nonzero
polynomial of degree at most four has strictly positive norm. This quantifies
the information that the coarse test actually misses.

\vfig{resolution_06_degree_gap}{The blue boundary certifies positivity
for every polynomial through that degree. The orange boundary gives an
explicit detecting polynomial. The shaded gap remains unresolved by these
certificates. At $\delta=10^{-12}$, degrees through four pass, degree seven
fails, and degrees five and six are undecided here.}

\claim{Interpolation controls the whole polynomial space}
Choose the odd squarefree labels $n_i=1,3,5,7,11$, write $y_i=x_{n_i}$,
and put
\[
b_i=\prod_{j\ne i}\frac{\max(y_j,1-y_j)}{|y_i-y_j|},
\qquad K_4=\sum_{i=0}^{4}n_i^4b_i^2.
\]
Lagrange interpolation and weighted Cauchy--Schwarz show that, for every
polynomial $f$ of degree at most four,
\[
Q_\delta(f)\ge
\left(1-\frac{\delta K_4}{12}\right)
\sum_{i=0}^{4}n_i^{-4}f(y_i)^2.
\]
Exact rational comparison gives $12/K_4>7.9790\times10^{-11}>10^{-12}$.
A nonzero polynomial cannot vanish at all five selected nodes, so the right
side is strictly positive. Thus the order-five Hankel matrix is positive
definite, while the degree-seven witness makes the order-eight matrix fail
positive semidefiniteness.

The complete atlas contains the general interpolation and detector bounds.
Its reproducible check records 985 exact assertions and 16 floating algebra
checks. The proofs establish the universal statements. The tests verify the
stated finite certificates and implementations.

\MonographPart{V}{Proofs, Evidence, and Sources}{Detailed arguments, reproducible checks, and a bounded comparison with preceding literature.}
\pageheading{TECHNICAL APPENDICES}{Complete derivations and their dependencies}
The illustrated chapters develop the definitions, decisive proof steps, and worked examples. The appendices provide the complete derivations, explicit constants, computational contracts, source comparisons, and unresolved questions. Editable source files accompany the monograph.

\begin{tabular}{@{}p{0.13\linewidth}p{0.35\linewidth}p{0.45\linewidth}@{}}
\toprule
Appendix & Proof chapter & Main dependencies\\\midrule
A & Dilation type and continuity & Positive coefficient recurrence, radius bootstrap, translated-radius comparison\\
B & Arithmetic inverse and character resolution & Summable positive atoms, squarefree counting, Euler factors, sector partial integration\\
C & Multiple-prime geometry and detectors & Exact moments, isolated-node polynomials, definite matrix pencils\\
D & Orbit continuation and shifts & Genus-zero product, branch separation, exponential sums, two-end Mellin subtraction\\
E & Integer sequences and polyphase laws & Factorial recurrence, ordered trees, root filtering, carry arithmetic\\
F & Coupled entire growth and type & Positive coupled recurrence, characteristic matrix, translated powers, left eigenvector\\
G & Component reduction and exact inverse & Cubic and fifth coefficients, factorization, seventh-coefficient separation\\
H & Constant linear operator structure & Projection fibers, polynomial interpolation, monomial tests\\
I & Quantitative resolution atlas & Tap type bounds, removed-prime mass, exact interpolation and tail certificates\\
J & Coupled exact laboratory & Derivative recurrence, denominator clearing, positive disk majorants\\
\bottomrule\end{tabular}

Inherited entire existence, nonrescaling, arithmetic non-D-finiteness,
one-prime transport, and base orbit convergence are developed in the visual
chapters and in the preserved baseline routes. The preceding proofs were independently reviewed. The foundational audit pins Appendices
A--E and their exact API. The coupled-family audit pins the coupled proofs,
resolution certificates, and exact APIs underlying Appendices F--J.

Classical external leaves are named explicitly. In particular, ordinary
orthogonal-polynomial theory, Hadamard factorization, Laplace uniqueness,
and spectral Mellin regularization are established machinery. The local
proofs verify their application and do not presume historical novelty.

\clearpage\appendix
\hypertarget{positive-lower-type-and-quantitative-tap-continuity}{%
\section{Positive lower type and quantitative tap
continuity}\label{positive-lower-type-and-quantitative-tap-continuity}}

This chapter sharpens the entire-growth theorem for nonlinear contraction families. The symmetric definition has an established antecedent. Historical priority of the asymmetric extension is not claimed.

\hypertarget{the-exact-growth-exponent-has-positive-lower-type}{%
\subsection{The exact growth exponent has positive lower
type}\label{the-exact-growth-exponent-has-positive-lower-type}}

Let

\[
F'(z)=1+\sum_{j=1}^{J}h_jF(q_jz)F(r_jz),\qquad F(0)=0,
\]

where \(h_j>0\), \(\sum_jh_j=1\), and \(0<q_j,r_j<1\). Zero-weight
channels may instead be allowed if they are removed from the active bank
before defining its growth exponent. Put

\[
\rho=\max_{j:h_j>0}\rho_j,\qquad q_j^{\rho_j}+r_j^{\rho_j}=1.
\]

The inherited coefficient recurrence gives nonnegative coefficients,
entire extension, and \(M(R)=\max_{|z|=R}|F(z)|=F(R)\ge R\) for \(R>0\).

\textbf{Theorem 1.} There are explicit constants \(0<c_-\le c_+<\infty\)
such that

\[
c_-\le\liminf_{R\to\infty}\frac{\log M(R)}{R^\rho}
\le\limsup_{R\to\infty}\frac{\log M(R)}{R^\rho}\le c_+.
\]

In particular the bank has positive lower type and finite upper type at
its exact order, including when \(0<\rho<1\). Furthermore

\[
\lim_{R\to\infty}\frac{\log\log M(R)}{\log R}=\rho.
\]

This last radial growth-exponent limit does not assert a type limit or
uniform angular growth.

\hypertarget{explicit-lower-constants-and-proof}{%
\subsubsection{Explicit lower constants and
proof}\label{explicit-lower-constants-and-proof}}

Select any active maximizing channel and abbreviate it to \((h,q,r)\),
so \(q^\rho+r^\rho=1\). Define

\[
L=\frac1h,\qquad
C=\max\left(\frac{qL}{1-q},\frac{rL}{1-r}\right),\qquad
B=L+\frac2{\min(q,r)},\qquad
c_-=\frac{\log2}{(B+C)^\rho}.
\]

We prove the stronger bound

\[
\boxed{\log M(R)\ge c_-(R+C)^\rho\quad\text{for every }R\ge2.}
\]

Write \(H(R)=\log M(R)\). On the seed interval \([2,B]\),
\(H(R)\ge\log R\ge\log2\ge c_-(R+C)^\rho\). For \(R\ge B\), integrate
the positive feedback over the fixed-width interval \([R-L,R]\):

\[
\begin{aligned}
M(R)&\ge h\int_{R-L}^{R}M(qt)M(rt)\,dt\\
&\ge hL M(q(R-L))M(r(R-L))\\
&=M(q(R-L))M(r(R-L)).
\end{aligned}
\]

The two child radii are at least two because \(R\ge B\). Moreover

\[
q(R-L)+C\ge q(R+C),\qquad
r(R-L)+C\ge r(R+C),
\]

because respectively \((1-q)C\ge qL\) and \((1-r)C\ge rL\). If the
comparison bound is known at both child radii, it follows that

\[
\begin{aligned}
H(R)&\ge H(q(R-L))+H(r(R-L))\\
&\ge c_-\{[q(R-L)+C]^\rho+[r(R-L)+C]^\rho\}\\
&\ge c_-(q^\rho+r^\rho)(R+C)^\rho
=c_-(R+C)^\rho.
\end{aligned}
\]

For a rigorous propagation, put \(Q_* =\max(q,r)<1\) and
\(d=(1-Q_*)B+Q_*L>0\). Every child of every \(R\ge B\) is at most
\(R-d\). Induction on the intervals \([B+kd,B+(k+1)d]\), beginning with
the already controlled seed, proves the comparison for all larger radii.
Continuity deals with common endpoints. All power comparisons use only
monotonicity of \(x^\rho\) for positive \(x\); no convexity assumption
is made, so the argument also covers \(\rho<1\).

\hypertarget{explicit-finite-upper-constants}{%
\subsubsection{Explicit finite upper
constants}\label{explicit-finite-upper-constants}}

For all channels set

\[
Q=\max_j\max(q_j,r_j),\quad
m=\min_j\min(q_j,r_j),\quad
\Delta=\log(1/m),\quad
B_+=\max_j\log\frac1{q_jr_j}+\log2.
\]

Choose an integer \(N\ge0\) with \(Q^{-N}\ge3/m\), and put

\[
V_N=(2^{N+1}-1)\log2+
       (2^{N+1}-N-2)\log(1/Q),
\]

\[
c_+=e^{-\rho}(V_N+1+\Delta+B_+).
\]

These constants are intentionally conservative, and require no numerical
evaluation of an unknown solution. The tangent majorant gives
\(M(1)<2\). At \(R_k=Q^{-k}\), the integral equation gives

\[
M(R_k)\le R_k[1+M(R_{k-1})^2].
\]

Define \(U_0=2\), \(U_k=Q^{-k}(1+U_{k-1}^2)\). Then \(M(R_k)\le U_k\),
\(U_k\ge1\), and \(\log U_k\le k\log(1/Q)+\log2+2\log U_{k-1}\). Summing
this elementary recurrence gives \(\log U_N\le V_N\). Since \(e<3\), the
logarithmic-radius seed \(1\le t\le1+\Delta\) has
\(e^t\le e/m<3/m\le R_N\), and therefore \(\log M(e^t)\le V_N\).

Let \(a_j=\log(1/q_j)\), \(b_j=\log(1/r_j)\), and
\(G(t)=\log M(e^t)+t+B_+\). For \(t\ge1+\Delta\), all child moduli are
at least one, so the positive integral and weight normalization give

\[
\log M(e^t)\le t+\log2+
\max_j\{\log M(e^{t-a_j})+\log M(e^{t-b_j})\}.
\]

The definition of \(B_+\) implies

\[
G(t)\le\max_j\{G(t-a_j)+G(t-b_j)\}.
\]

On the seed, \(G(t)\le V_N+1+\Delta+B_+\le c_+e^{\rho t}\). Each channel
satisfies \(q_j^\rho+r_j^\rho\le1\). Induction in intervals of length
\(\min_j\min(a_j,b_j)>0\) therefore propagates the seed comparison,
giving the explicit upper bound

\[
\boxed{\log M(R)\le c_+R^\rho-\log R-B_+\quad(R\ge e).}
\]

Dividing the two bounds by \(R^\rho\), then taking logarithms of their
positive logarithmic-growth bounds, proves all assertions of Theorem 1.
The theorem leaves open whether \(\log M(R)/R^\rho\) has a limit.
Lattice resonances or logarithmic oscillations are possible research
directions, not results here.

\hypertarget{joint-holomorphic-dependence-on-parameters}{%
\subsection{Joint holomorphic dependence on
parameters}\label{joint-holomorphic-dependence-on-parameters}}

Fix the number \(J\) of channels. Allow complex weights \(h_j\), without
normalization, and complex contractions satisfying \(|q_j|,|r_j|<1\).
The same formal recurrence defines coefficients

\[
c_n(h,q,r)\in\mathbb Q[h_1,\ldots,h_J,q_1,r_1,\ldots,q_J,r_J].
\]

Each is homogeneous of degree \(n\) in the weights. If a parameter set
has \(|q_j|,|r_j|\le Q<1\) and \(\sum_j|h_j|\le W\), with \(W>0\),
coefficient induction gives

\[
|c_n(h,q,r)|\le W^n d_n(Q),
\]

where \(f_Q(z)=\sum d_n(Q)z^{2n+1}\) is the established diagonal
function. Consequently, on every disk \(|z|\le R\),

\[
\sum_{n\ge0}|c_n(h,q,r)|R^{2n+1}
\le W^{-1/2}f_Q(\sqrt W R)<\infty.
\]

Every compact subset of the complex parameter domain has such \(Q,W\).
The series therefore converges normally on compact sets of the joint
domain, and its sum is jointly holomorphic in the weights, contractions,
and \(z\). This proves parameter continuity at zero weights, although
the growth order can jump there. Complex or signed parameters do not
inherit the positive-bank growth theorem merely from holomorphic
dependence.

\hypertarget{a-quantitative-compact-bound-for-changing-tap-weights}{%
\subsection{A quantitative compact bound for changing tap
weights}\label{a-quantitative-compact-bound-for-changing-tap-weights}}

Keep fixed real contractions \(0<q_j,r_j\le Q<1\), and let \(h,g\) be
nonnegative normalized weight vectors. Zero weights are permitted. Let
\(d=\|h-g\|_1\). Then for every \(R\ge0\),

\[
\boxed{\sup_{|z|\le R}|F_h(z)-F_g(z)|
\le\frac d4\{R f_Q'(R)-f_Q(R)\}.}
\]

\textbf{Proof.} Regard \(c_n\) as a polynomial in the real weight
variables with the contractions fixed. Its coefficients are nonnegative.
Induction in the recurrence gives the coefficientwise polynomial bound

\[
0\preceq c_n(h)\preceq d_n(Q)\Bigl(\sum_jh_j\Bigr)^n.
\]

Here \(\preceq\) compares every monomial coefficient. Differentiation
preserves this comparison, so on the normalized nonnegative simplex

\[
0\le\partial_{h_j}c_n\le n d_n(Q).
\]

Put \(v=h-g\). Its positive and negative parts both have sum \(d/2\),
since \(\sum_jv_j=0\). At every point of the segment between \(g\) and
\(h\), all partial derivatives lie in the interval \([0,n d_n(Q)]\).
Hence

\[
|v\cdot\nabla c_n|\le\frac d2 n d_n(Q),\qquad
|c_n(h)-c_n(g)|\le\frac d2 n d_n(Q).
\]

Multiply by \(R^{2n+1}\) and sum. Entire convergence justifies the sum,
and
\[R f_Q'(R)-f_Q(R)=2\sum_{n\ge0}n d_n(Q)R^{2n+1},\]
proving the
claim. The same argument applies to the closed simplex and needs no
positive minimum weight.

\hypertarget{a-vanishing-tap-and-discontinuous-order}{%
\subsubsection{A vanishing tap and discontinuous
order}\label{a-vanishing-tap-and-discontinuous-order}}

Take the two diagonal channels \((q,r)=(1/4,1/4)\) and \((3/4,3/4)\),
with weights \(1-\varepsilon\) and \(\varepsilon\), respectively. For
\(0\le\varepsilon\le1\), write the resulting function as
\(F_\varepsilon\). The quantitative bound gives

\[
\sup_{|z|\le R}|F_\varepsilon(z)-f_{1/4}(z)|
\le\frac{\varepsilon}{2}\{R f_{3/4}'(R)-f_{3/4}(R)\}.
\]

Thus convergence to \(f_{1/4}\) is uniform on every fixed disk, at an
explicit linear rate in \(\varepsilon\). Yet

\[
\rho(F_0)=\frac12,\qquad
\rho(F_\varepsilon)=\frac{\log2}{\log(4/3)}
\quad(0<\varepsilon\le1).
\]

For the faster active channel the lower constants become
\(L=1/\varepsilon\), \(C=3/\varepsilon\), \(B=1/\varepsilon+8/3\). Its
certified lower-type constant is

\[
c_-(\varepsilon)=\frac{\log2}{(4/\varepsilon+8/3)^{\log2/\log(4/3)}}>0.
\]

It tends to zero as the faster tap vanishes. These are proved growth and
compact-convergence statements. They do not determine an actual
crossover radius or leading asymptotic constant.

\hypertarget{computational-contract-and-limitations}{%
\subsection{Computational contract and
limitations}\label{computational-contract-and-limitations}}

\path{checks/dilation_type.py} uses exact \texttt{fractions.Fraction}
arithmetic for its finite assertions. It checks the algebraic
fixed-width and interval-propagation conditions for rational scales and
weights, the positive homogeneous weight-polynomial majorant, derivative
bounds on a normalized simplex, the improved total-variation coefficient
bound, and a distinct ordered-tree evaluation against the recurrence.
Scalar examples and the explicit real type constants are displayed at
high precision separately from exact assertions.

The finite checks do not prove a universal analytic theorem, normal
convergence, or historical novelty. Those depend on the proofs above. No
numerical type-limit fit is presented as evidence that a limit exists.

The earlier moving-width lower argument used a fixed \(\alpha<1\); it
proved orders below \(\rho\) and then recovered the order by
\(\alpha\uparrow1\), but did not establish positive lower type at
\(\rho\). This is a limitation of that argument, not a counterexample.
The fixed-width comparison closes precisely that missing implication.
Any execution failure of the new finite checks is recorded with its
cause and retained run metadata rather than erased.

\clearpage\hypertarget{arithmetic-depth-an-exact-inverse-its-conditioning-and-character-resolution}{%
\section{Arithmetic depth: an exact inverse, its conditioning, and
character
resolution}\label{arithmetic-depth-an-exact-inverse-its-conditioning-and-character-resolution}}



The entire squarefree depth function determines its arithmetic
parameters exactly, with one precise equivalence: powers of the primes
in the modulus are invisible. This is a local mathematical theorem. It
does not assert historical priority for Laplace inversion or for this
particular arithmetic application.

\hypertarget{claim-card-and-conventions}{%
\subsection{Claim card and
conventions}\label{claim-card-and-conventions}}

For real \(a>1/2\), real \(\beta>0\), and integer \(Q\ge1\), write

\[
 A_{a,\beta,Q}(s)=\sum_{\substack{n\ge2\\ \mu(n)^2=1\\ (n,Q)=1}}
 n^{-2a}\exp(-\lambda_n s),\qquad
 \lambda_n=-2\log(1-n^{-\beta}).
\]

The logarithm is real. The omitted \(n=1\) atom is not restored in any
inverse formula. Let \(n_1<n_2<\cdots\) be the support integers, and
abbreviate \(\lambda_j=\lambda_{n_j}\), \(w_j=n_j^{-2a}\). Then
\(\lambda_1>\lambda_2>\cdots>0\), \(\lambda_j\to0\), and
\(\sum_jw_j=M=Z_Q(2a)-1<\infty\), where

\[
 Z_Q(v)=\frac{\zeta(v)}{\zeta(2v)}
          \prod_{p\mid Q}(1+p^{-v})^{-1},\qquad v>1.
\]

Here a germ means the complete exact analytic germ, equivalently all
exact derivatives at a point, rather than a finite Taylor polynomial or
noisy measurements. The base route proves that this germ has a unique
entire continuation, with
\(|A(s)|\le M\exp(\lambda_1\max\{-\Re s,0\})\).

\hypertarget{exact-inverse-theorem}{%
\subsection{Exact inverse theorem}\label{exact-inverse-theorem}}

\textbf{Theorem AI-1.} Two admissible parameter triples define the same
analytic germ if and only if

\[
 a=a',\qquad \beta=\beta',\qquad
 \operatorname{rad}(Q)=\operatorname{rad}(Q').
\]

The function itself therefore encodes the weight exponent, the depth
exponent, and the finite excluded prime set. It does not encode
valuations \(v_p(Q)\).

\hypertarget{step-1-recover-the-ordered-spectral-atoms}{%
\subsubsection{Step 1: recover the ordered spectral
atoms}\label{step-1-recover-the-ordered-spectral-atoms}}

Starting with \(R_1=A\), define successively

\[
 \lambda_j=\lim_{t\to\infty}\frac{\log R_j(-t)}t,
 \qquad
 w_j=\lim_{t\to\infty}e^{-\lambda_jt}R_j(-t),
 \qquad
 R_{j+1}(s)=R_j(s)-w_je^{-\lambda_js}.
\]

Every limit exists and recovers exactly the \(j\)-th support frequency
and weight. Indeed, if \(M_{j+1}=\sum_{k>j}w_k\), then for every
\(t\ge0\),

\[
 R_j(-t)=w_je^{\lambda_jt}(1+r_j(t)),\qquad
 0\le r_j(t)\le\frac{M_{j+1}}{w_j}
                    e^{-(\lambda_j-\lambda_{j+1})t}.
\]

The bound follows by replacing every tail frequency by its largest value
\(\lambda_{j+1}\). Taking logarithms and dividing by \(t\) proves the
first limit. Multiplying by \(e^{-\lambda_jt}\) proves the second.
Removing this exact entire atom leaves the same hypotheses for the next
step. Thus the complete ordered atom list is intrinsic to the entire
function.

No approximation to a large negative-depth value is used in this
uniqueness proof. It is an exact limiting argument.

\hypertarget{step-2-recover-beta-without-knowing-the-integer-support}{%
\subsubsection{\texorpdfstring{Step 2: recover \(\beta\) without knowing
the integer
support}{Step 2: recover \textbackslash beta without knowing the integer support}}\label{step-2-recover-beta-without-knowing-the-integer-support}}

Put

\[
 \delta_Q=\frac1{\zeta(2)}\prod_{p\mid Q}\frac{p}{p+1}>0.
\]

The base route establishes the elementary counting estimate

\[
 N_Q(x):=\#\{2\le n\le x:\mu(n)^2=1,(n,Q)=1\}
       =\delta_Qx+O_Q(\sqrt{x}).
\]

Since \(N_Q(n_j)=j\), it follows that

\[
 n_j=\frac{j}{\delta_Q}\left(1+O_Q(j^{-1/2})\right),
 \quad
 \lambda_j=2\delta_Q^{\beta}j^{-\beta}
     \left(1+O_{Q,\beta}(j^{-1/2}+j^{-\beta})\right).
\]

For completeness, the first estimate follows initially from
\(j=\delta_Qn_j+O_Q(\sqrt{n_j})\), which gives \(n_j\asymp_Qj\);
substituting this back improves the error to \(O_Q(\sqrt j)\). The
second uses \(-\log(1-u)=u(1+O(u))\) as \(u\downarrow0\).

Consequently the already recovered frequencies determine

\[
 \boxed{\displaystyle
 \beta=-\lim_{j\to\infty}\frac{\log\lambda_j}{\log j}.}
\]

A rank ratio avoids the slowly decaying constant divided by \(\log j\):
for every fixed real \(c>1\),

\[
 -\frac{\log(\lambda_{\lfloor cj\rfloor}/\lambda_j)}{\log c}
 =\beta+O_{Q,\beta,c}(j^{-1/2}+j^{-\beta}).
\]

The floor contributes an additional \(O(j^{-1})\), which is absorbed by
the displayed error. These are fixed-parameter asymptotics. They do not
give a constant uniform in arbitrary moduli.

\hypertarget{step-3-recover-the-support-integers-and-the-weight-exponent}{%
\subsubsection{Step 3: recover the support integers and the weight
exponent}\label{step-3-recover-the-support-integers-and-the-weight-exponent}}

Once \(\beta\) is known, every spectral atom gives its exact integer
label:

\[
 \boxed{\displaystyle
 n_j=(1-e^{-\lambda_j/2})^{-1/\beta},\qquad
 a=-\frac{\log w_j}{2\log n_j}.}
\]

The recovered real number in the first formula is an integer by the
family definition; rounding is not part of the exact theorem. Since
\(n_j\ge2\), the denominator in the second formula is nonzero. Every
recovered atom yields the same value of \(a\).

\hypertarget{step-4-recover-the-excluded-primes}{%
\subsubsection{Step 4: recover the excluded
primes}\label{step-4-recover-the-excluded-primes}}

A prime \(p\) lies in the recovered integer support exactly when
\(p\nmid Q\): primes are squarefree, and no other restriction is
imposed. Hence

\[
 \boxed{\displaystyle
 \operatorname{rad}(Q)=\prod_{\substack{p\ \mathrm{prime}\\
                                      p\notin\{n_j:j\ge1\}}}p.}
\]

The product is finite because the parameters are assumed to come from
this family. This formula is an exact reconstruction from infinite data,
not an algorithm that can certify completion of the prime list after
finitely many observations. Conversely, equal \(a,\beta\) and equal
radicals produce exactly the same terms. This proves both directions of
AI-1.

\hypertarget{a-constructive-counting-envelope}{%
\subsubsection{A constructive counting
envelope}\label{a-constructive-counting-envelope}}

The rank asymptotic can be made effective without an RH input. Write
\(\omega(Q)=\#\{p:p\mid Q\}\). For real \(x\ge4\), inclusion and
exclusion gives

\[
 |N_Q(x)-\delta_Qx|\le C_Q\sqrt x+1,
 \qquad C_Q=2^{\omega(Q)}+2\frac{\varphi(Q)}Q.
\]

To verify the constants, expand \(\mu(n)^2=\sum_{d^2\mid n}\mu(d)\). For
\((d,Q)=1\), the number of \(m\le x/d^2\) coprime to \(Q\) differs from
\((\varphi(Q)/Q)x/d^2\) by at most \(2^{\omega(Q)}\). Summing over at
most \(\sqrt x\) values of \(d\) gives the first part of \(C_Q\). The
omitted main-term tail is at most
\((\varphi(Q)/Q)x/\lfloor\sqrt x\rfloor\le2(\varphi(Q)/Q)\sqrt x\).
Removing \(n=1\) adds one. The main-term density is
\((\varphi(Q)/Q)\sum_{(d,Q)=1}\mu(d)d^{-2}=\delta_Q\) by its absolutely
convergent Euler product.

For \(j\ge3\), therefore,

\[
 \left(\frac{\sqrt{C_Q^2+4\delta_Q(j-1)}-C_Q}{2\delta_Q}\right)^2
 \le n_j\le
 \left(\frac{\sqrt{C_Q^2+4\delta_Q(j+1)}+C_Q}{2\delta_Q}\right)^2.
\]

The lower and upper bounds are the positive solutions of the
corresponding quadratic inequalities in \(\sqrt{n_j}\). At \(j\ge3\),
all possible supports have \(n_j\ge5\), so the \(x\ge4\) counting
envelope applies.

\hypertarget{what-finite-negative-depth-measurements-can-recover}{%
\subsection{What finite negative-depth measurements can
recover}\label{what-finite-negative-depth-measurements-can-recover}}

\textbf{Theorem AI-2, gap bound.} Suppose the first \(j-1\) atoms have
already been removed exactly, put
\(\Delta_j=\lambda_j-\lambda_{j+1}>0\), and fix \(h>0\). Define

\[
 \widehat\lambda_j(t,h)
 =\frac{\log R_j(-t-h)-\log R_j(-t)}h.
\]

Then for every \(t\ge0\),

\[
 0\le\lambda_j-\widehat\lambda_j(t,h)
 \le\frac1h\log\!\left(1+\frac{M_{j+1}}{w_j}e^{-\Delta_jt}\right),
\]

and

\[
 0\le e^{-\lambda_jt}R_j(-t)-w_j
     \le M_{j+1}e^{-\Delta_jt}.
\]

\textbf{Proof.} The preceding positive remainder formula gives
\(\widehat\lambda_j=\lambda_j+h^{-1}[\log(1+r_j(t+h))-\log(1+r_j(t))]\).
Each tail factor decreases, so \(0\le r_j(t+h)\le r_j(t)\). The two
inequalities now follow from the bound on \(r_j(t)\). For the amplitude,
multiply the positive tail bound by \(e^{-\lambda_jt}\).

In particular, to guarantee frequency error at most \(\varepsilon>0\)
from exact residual values, it suffices to take

\[
 t\ge\max\left\{0,\frac1{\Delta_j}
      \log\frac{M_{j+1}}{w_j(e^{h\varepsilon}-1)}\right\}.
\]

The parameters in this certificate are properties of the true tail. If
they are only estimated from data, this is a conditional recovery bound
rather than an independent numerical certificate.

\hypertarget{positive-tail-envelopes-for-reproducible-examples}{%
\subsubsection{Positive tail envelopes for reproducible
examples}\label{positive-tail-envelopes-for-reproducible-examples}}

For integer \(L\ge n_j\), a direct sum through \(n\le L\) at depth
\(-t\), \(t\ge0\), omits at most

\[
 E_L(t)=e^{t\lambda_{L+1}}\frac{L^{1-2a}}{2a-1}.
\]

Indeed all omitted frequencies are at most \(\lambda_{L+1}\), and
\(\sum_{n>L}n^{-2a}\le\int_L^\infty x^{-2a}\,dx\). Ignoring
squarefreeness and coprimality makes this a valid upper bound. If
\(r_0,r_1\) are the two positive truncated residual sums at \(t,t+h\),
the true secant belongs to

\[
 \left[
 \frac{\log r_1-\log(r_0+E_L(t))}h,
 \frac{\log(r_1+E_L(t+h))-\log r_0}h
 \right].
\]

This interval treats truncation exactly at the mathematical level. Its
displayed floating-point endpoints are not directed-rounding
certificates.

\hypertarget{cancellation-is-a-real-limitation}{%
\subsubsection{Cancellation is a real
limitation}\label{cancellation-is-a-real-limitation}}

Evaluating a residual by subtracting previous atoms is different from
evaluating its known positive series. If an observation of \(A(-t)\) has
relative error at most \(\eta\) and the previous atoms are exact, then
the residual error can have relative size

\[
 \eta\,\kappa_j(t),\qquad
 \kappa_j(t)=\frac{A(-t)}{R_j(-t)}
 \sim\frac{w_1}{w_j}e^{(\lambda_1-\lambda_j)t}.
\]

Thus the large negative depth that separates nearby atoms also worsens
cancellation for later atoms. Uncertain previously subtracted
frequencies introduce further error. The computation evaluates residuals
as positive sums and reports \(\log_{10}\kappa_j\) so this conditioning
cost remains visible.

\hypertarget{exact-uniqueness-does-not-imply-uniform-prime-recovery}{%
\subsection{Exact uniqueness does not imply uniform prime
recovery}\label{exact-uniqueness-does-not-imply-uniform-prime-recovery}}

\textbf{Theorem AI-3, bounded-window instability.} Let \(p\) be a prime
not dividing \(Q\), and let \(R\ge0\). For every integer \(k\ge0\),

\[
 \boxed{\displaystyle
 \sup_{|s|\le R}|A_{a,\beta,Q}^{(k)}(s)-A_{a,\beta,Qp}^{(k)}(s)|
 \le \lambda_p^k e^{R\lambda_p}
      \frac{Z_Q(2a)}{p^{2a}+1}.}
\]

Consequently adding larger and larger excluded primes gives convergence
in every derivative on every bounded complex-depth window, even though
the recovered radical changes each time.

\textbf{Proof.} The difference consists exactly of the squarefree
integers divisible by \(p\) and coprime to \(Q\), all of which satisfy
\(n\ge p\). Their frequencies are at most \(\lambda_p\). The Euler
factor at \(p\) shows that their total weight is

\[
 \sum_{\substack{n\ge1\\\mu(n)^2=1,(n,Q)=1\\p\mid n}}n^{-2a}
 =\frac{Z_Q(2a)}{p^{2a}+1}.
\]

There is no \(n=1\) term in this difference. Bound the \(k\)-th
derivative of every exponential by \(\lambda_p^ke^{R\lambda_p}\) and
sum. Since \(\lambda_p\to0\) and \(p^{-2a}\to0\), the bound tends to
zero. This proves the derivative convergence and the claimed
instability.

More explicitly, for any bounded window and positive measurement
tolerance, two distinct radicals in this family can produce functions
closer than that tolerance throughout the window. This defeats a
recovery guarantee uniform over unbounded excluded primes. It does not
refute recovery under an a priori modulus bound or from complete exact
analytic data.

For \(a=\beta=2,Q=1\), the mass removed by \(p\) is exactly
\(105/[\pi^4(p^4+1)]\). The examples \(p=101,1009,10007\) quantify how
quickly large-prime distinctions disappear on \(|s|\le5\).

\hypertarget{the-function-as-an-arithmetic-hilbert-feature-space}{%
\subsection{The function as an arithmetic Hilbert feature
space}\label{the-function-as-an-arithmetic-hilbert-feature-space}}

This section spells out the classical feature-space construction for
this arithmetic spectrum. It supplies a precise space in which the
complex-depth packets live.

For \(c=(c_j)\in\ell^2\), define

\[
 (Uc)(s)=\sum_{j\ge1}c_j\sqrt{w_j}\,e^{-\lambda_js}.
\]

Cauchy and Schwarz give uniform convergence on bounded complex sets and

\[
 |(Uc)(s)|\le\|c\|_{\ell^2}\sqrt{A(2\Re s)}.
\]

The map is injective: \(\sum|c_j|\sqrt{w_j}\le\|c\|\sqrt M<\infty\), so
the same spectral peeling argument, now for complex coefficients, forces
the first coefficient and then every coefficient to vanish whenever
\(Uc=0\). Positivity was needed to recover unknown frequencies in AI-1;
injectivity here only needs absolute summability and the known ordered
support. Specifically, after dividing by the largest remaining
exponential, the tail is bounded by
\(e^{-\Delta_jt}\sum_{k>j}|c_k|\sqrt{w_k}\), which tends to zero.

The image \(\mathcal H_{a,\beta,Q}=U\ell^2\), with
\(\|Uc\|_{\mathcal H}=\|c\|_{\ell^2}\), is therefore a complete Hilbert
space of entire functions. Its reproducing kernel is

\[
 K(s,t)=A(s+\overline t).
\]

Using an inner product linear in the first argument, the coefficient
vector of \(K(\cdot,t)\) is \((\sqrt{w_j}e^{-\lambda_j\overline t})_j\),
so \(\langle Uc,K(\cdot,t)\rangle=(Uc)(t)\). This also verifies the sign
and conjugation convention directly. Arbitrarily large finite Gram
matrices at distinct complex depths are strictly positive, as proved in
the base route. Thus the feature space has infinite rank, not a
finite-dimensional Prony model.

The monomial moment matrices also have infinite rank:
\(m_k=(-1)^kA^{(k)}(0)=\sum_jw_j\lambda_j^k\), and for every nonzero
polynomial \(P\), \(\sum_jw_j|P(\lambda_j)|^2>0\), because a nonzero
polynomial cannot vanish at infinitely many distinct frequencies. All
finite derivative Hankel matrices are strictly positive definite.

\hypertarget{complete-fixed-modulus-character-resolution}{%
\subsection{Complete fixed-modulus character
resolution}\label{complete-fixed-modulus-character-resolution}}

Let \(\chi,\psi\) be Dirichlet characters modulo the same fixed integer
\(Q\), extended by zero outside the units, and put
\(\eta=\chi\overline\psi\). The base extension is

\[
 A_{\chi,\psi}(s)=\sum_{n\ge2}\mu(n)^2
            \eta(n)n^{-2a}(1-n^{-\beta})^{2s}.
\]

The endpoint subtraction in its Newton expansion is still exactly one:

\[
 A_{\chi,\psi}(s)
 =\sum_{j\ge0}(-1)^j\binom{2s}{j}
 \left[\frac{L(2a+j\beta,\eta)}
              {L(4a+2j\beta,\eta^2)}-1\right].
\]

Every character takes the value one at \(n=1\); at primes dividing
\(Q\), it is zero. The Euler product identity is valid even when
\(\eta\) is imprimitive, because the actual Dirichlet character and its
actual modulus are retained. All arguments of \(L\) lie in \(\Re v>1\),
where its Euler product converges absolutely and is nonzero. No unproved
zero-free-region hypothesis is used.

\textbf{Theorem AI-4, character sector decay.} Fix \(a>1/2,\beta>0,Q\),
and a closed sector \(|\arg s|\le\pi/2-\varepsilon\),
\(0<\varepsilon<\pi/2\). If \(\chi\ne\psi\), then uniformly in that
sector,

\[
 A_{\chi,\psi}(s)
 =O_{a,\beta,Q,\varepsilon}\bigl(|s|^{-(2a-1/2)/\beta}\bigr),
 \qquad
 \frac{A_{\chi,\psi}(s)}{A_{a,\beta,Q}(s)}
 =O_{a,\beta,Q,\varepsilon}\bigl(|s|^{-1/(2\beta)}\bigr).
\]

These statements are uniform over the finite character set of this fixed
modulus. They are not uniform as \(Q\) grows.

\textbf{Proof.} Distinct characters give a nonprincipal \(\eta\). Its
sum over a complete period is zero, so \(|\sum_{m\le y}\eta(m)|\le Q\).
Hence

\[
 \sum_{n\le x}\mu(n)^2\eta(n)
 =\sum_{d\le\sqrt x}\mu(d)\eta(d)^2
     \sum_{m\le x/d^2}\eta(m)=O_Q(\sqrt x).
\]

Use partial integration against \(g_s(x)=x^{-2a}(1-x^{-\beta})^{2s}\),
whose endpoint values vanish when \(\Re s>0\). The counting bound
reduces the magnitude to a constant times
\(\int_1^\infty\sqrt x\,|g_s'(x)|\,dx\). Writing \(u=\Re s\) and
\(\gamma=(2a-1/2)/\beta>0\), direct substitution \(y=x^{-\beta}\) bounds
its two derivative terms by

\[
 \frac{2a}{\beta}B(\gamma,2u+1)
 +2|s|B(\gamma+1,2u),
\]

which is \(O_{a,\beta,\varepsilon}(|s|^{-\gamma})\) because
\(u\ge |s|\sin\varepsilon\). The principal diagonal has the nonzero
leading term \(\delta_Q\Gamma(\alpha)\beta^{-1}(2s)^{-\alpha}\),
\(\alpha=(2a-1)/\beta\), by the base sector theorem. Its modulus is
bounded below by a positive fixed-sector multiple of \(|s|^{-\alpha}\)
for sufficiently large \(|s|\). Division gives the second bound, since
\(\gamma-\alpha=1/(2\beta)\).

At real positive depth, the matrix \((A_{\chi,\psi}(s))\) on any \(d\)
distinct characters is positive semidefinite and has identical diagonal
\(A_{a,\beta,Q}(s)\). Its normalized matrix tends to \(I_d\), with
operator norm error \(O_{a,\beta,Q,d}(s^{-1/(2\beta)})\). Indeed, the
off-diagonal entry bound controls the maximum row and column sums, and
thus the operator norm. This gives quantitative separation of a fixed
character bank. It is consistent with a complex or signed off-diagonal
entry and makes no complete-monotonicity assertion for that entry.

\hypertarget{classical-comparison-and-limits-of-novelty}{%
\subsection{Classical comparison and limits of
novelty}\label{classical-comparison-and-limits-of-novelty}}

The general fact that a positive measure is determined by its Laplace
transform is classical. The primary publisher-distributed preview of
René L. Schilling, Renming Song, and Zoran Vondraček, \emph{Bernstein
Functions: Theory and Applications}, first edition, de Gruyter, 2010,
ISBN 978-3-11-021530-4 / electronic ISBN 978-3-11-021531-1, was checked
in the reviewed source:

\url{https://api.pageplace.de/preview/DT0400.9783110215311_A15706968/preview-9783110215311_A15706968.pdf}

Proposition 1.2, printed page 2 / PDF page 14 under one-based indexing,
states uniqueness of a measure on \([0,\infty)\) from its Laplace
transform and proves recovery from all derivatives. Theorem 1.4, printed
page 3 / PDF page 15, is the Bernstein representation theorem.
Translation here: the source's measure is
\(\sum_jw_j\delta_{\lambda_j}\), its transform variable is the present
real depth \(s>0\), and its finiteness condition holds because
\(\sum_jw_j<\infty\). Classification: general uniqueness is
\textbf{known after translation of notation}. The exact extraction of
\(a,\beta,\operatorname{rad}(Q)\), the rank asymptotic, and the
displayed large-prime instability are proved locally rather than claimed
as results of that book.

The Hilbert feature-space construction and finite Prony mechanism are
also classical constructions. They are explicit realizations for this
arithmetic spectrum, with no priority claim. The parameter-specific
inverse theorem and character extension are not asserted historically
new. The limited searches do not settle whether these results have
appeared together for this exact positive squarefree family.

Bounded discovery queries were
\texttt{\small squarefree Laplace transform inverse identification parameters spectral measure arithmetic}
and
\texttt{\small Bernstein functions theorem 1.4 completely monotone uniquely determines measure Schilling Song Vondracek pdf}.
The direct de Gruyter chapter URL returned an internal retrieval error.
The publisher-distributed preview supplied the needed theorem and full
proof instead. No conclusion depends on the failed URL, and the source package does not contain the publisher PDF. The elementary inverse
and character proofs do not require this external source as an unchecked
leaf.

\hypertarget{computation-contract-and-handoff}{%
\subsection{Computation contract and
handoff}\label{computation-contract-and-handoff}}

\path{checks/arithmetic_inverse.py} is deterministic. Its exact
checks use rational finite atomic surrogates with \(2a=4\), integer
\(\beta=1,2,3\), \(Q=1,2,6,30\), and one through five atoms. They
independently construct the annihilating polynomial from moments and
from its known rational nodes, verify the positive Vandermonde formula
for Hankel determinants, reconstruct weights, recover integer labels at
known \(\beta\), and test further exact recurrence moments. These
surrogates do not turn the infinite family's Hankel matrices into
finite-rank matrices.

The high-precision examples use 180 decimal digits. They show two-rank
\(\beta\) estimates at ranks 100, 1000, and 10000, conditional
negative-depth peeling for the first three atoms at \(t=20,100,500\),
and large-prime bounded-window envelopes. Positive direct sums use
\(n\le5000\), with the analytic omitted-tail envelope above. Rank
enumeration uses \(n\le100000\). Actual cancellation factors are
reported. Numerical values and floating-point endpoints are
illustrations, not ball-arithmetic certificates.

Run from the source directory with
\texttt{\small python checks/arithmetic\_inverse.py}; the project dependency
fallback is discovered by the script in this session, while a portable
environment needs \texttt{mpmath}. Results are written to
\path{evidence/arithmetic_inverse_checks.json}, including
input/source hashes, arithmetic conventions, tested ranges, software
versions, elapsed time, and the source-retrieval failure. Universal
conclusions AI-1 through AI-4 rest on the local proofs above. Historical
priority remains unresolved.

\clearpage\hypertarget{finite-prime-sieve-families-and-explicit-detection-of-signed-atoms}{%
\section{Finite-prime sieve families and explicit detection of signed
atoms}\label{finite-prime-sieve-families-and-explicit-detection-of-signed-atoms}}

The construction extends the one-prime setting to a finite-prime parameter family. The results below are local
mathematical derivations. They do not assert historical priority.
Classical orthogonal-polynomial mechanisms are identified as such.

\hypertarget{definition-and-exact-coordinates}{%
\subsection{Definition and exact
coordinates}\label{definition-and-exact-coordinates}}

Let \(F\) be a finite set of distinct primes, let
\(\boldsymbol\lambda=(\lambda_p)_{p\in F}\), and put

\[
x_n=(1-n^{-2})^2,\qquad
\kappa_F(n;\boldsymbol\lambda)=\prod_{p\in F}(1+\lambda_p\mathbf1_{p\mid n}),
\] \[
\nu_{F,\boldsymbol\lambda}
=\sum_{n\ge1}\mu(n)^2n^{-4}\kappa_F(n;\boldsymbol\lambda)\delta_{x_n}.
\]

The mass at \(x_1=0\) is exactly 1. A strength \(\lambda_p=-1\) deletes
every squarefree atom divisible by \(p\). Several deletions commute. The
depth-\(k\) family uses \(x^k\,d\nu\), for integers \(k\ge0\).

Write \(D(s)=\zeta(s)/\zeta(2s)\). For arbitrary finite complex
strengths and \(\Re s>1\), absolute convergence permits expansion of the
finite prime product. The local squarefree Euler factors give

\[
Z_F(s;\boldsymbol\lambda)
=\sum_{n\ge1}\mu(n)^2n^{-s}\kappa_F(n;\boldsymbol\lambda)
=D(s)\prod_{p\in F}\left(1+\frac{\lambda_p}{p^s+1}\right).
\]

Equivalently, expansion in subsets \(S\subseteq F\) gives the
coefficient \(D(s)/\prod_{p\in S}(p^s+1)\) for
\(\prod_{p\in S}\lambda_p\). This is an independently useful coordinate
formula for prime overlap.

Consequently

\[
M_j(\boldsymbol\lambda)=\int x^j\,d\nu_{F,\boldsymbol\lambda}
=\sum_{r=0}^{2j}(-1)^r\binom{2j}{r}
 D(4+2r)\prod_{p\in F}
 \left(1+\frac{\lambda_p}{p^{4+2r}+1}\right).
\tag{1}
\]

Depth-\(k\) moments are \(M_{k+j}\). With \(v=\pi^{-2}\), the even-zeta
formula makes every \(D(4+2r)\) a rational multiple of \(v^{r+2}\). Thus
the moments belong to \(\mathbb Q[v,\lambda_p:p\in F]\), and monic
orthogonal-polynomial coefficients belong to
\(\mathbb Q(v,\lambda_p:p\in F)\). This is an exact field statement,
with no decimal recognition step.

For two primes the interaction term is explicit:

\[
M_j(\lambda_p,\lambda_q)=M_j(0,0)+\lambda_p U_{p,j}
+\lambda_q U_{q,j}+\lambda_p\lambda_q U_{pq,j},
\]

where \(U_{pq,j}=\sum_{pq\mid n}\mu(n)^2n^{-4}x_n^j>0\). More generally,
with all other strengths at least \(-1\), every mixed moment derivative
in a set of distinct primes is strictly positive. It is the integral
over atoms divisible by those primes, with the remaining nonnegative
prime factors. Infinitely many such atoms remain after any finite
collection of deletions. Pure second derivatives in any one strength
vanish because the moments are separately affine.

\hypertarget{the-exact-positivity-domain}{%
\subsection{The exact positivity
domain}\label{the-exact-positivity-domain}}

\textbf{Theorem 1.} For real strengths, \(\nu_{F,\boldsymbol\lambda}\)
is a positive measure if and only if \(\lambda_p\ge-1\) for every
\(p\in F\). On this closed orthant, every finite Hankel matrix of every
nonnegative integer depth is strictly positive definite.

\textbf{Proof.} The stated inequalities make each atom weight
nonnegative. The primes outside \(F\) supply infinitely many positive
atoms, even when all strengths are deletion endpoints. A nonzero
polynomial cannot vanish on infinitely many distinct nodes; therefore
its squared norm is strictly positive. Conversely, if \(\lambda_p<-1\),
the squarefree atom \(n=p\) has mass \((1+\lambda_p)p^{-4}<0\): no other
strength appears at that atom. The same argument works after every depth
tilt because \(x_p>0\). \(\square\)

The positive monic families \(P_m^{(k;F,\boldsymbol\lambda)}\)
consequently exist at all degrees throughout the closed orthant. Their
Hankel determinants never vanish there, including at simultaneous
deletion endpoints. Each finite-degree coefficient is analytic in a
neighborhood of every point of that orthant. This neighborhood is a
finite-degree statement; the next theorem explains why no common
enlargement of the positivity domain works at all degrees.

Together, Theorems 1 and 2 also prove the exact equivalence: every
finite Hankel matrix is positive semidefinite if and only if every
strength is at least \(-1\). Thus this is a complete positivity-domain
result, not only a sufficient condition.

\hypertarget{an-explicit-polynomial-detector-for-every-bad-strength}{%
\subsection{An explicit polynomial detector for every bad
strength}\label{an-explicit-polynomial-detector-for-every-bad-strength}}

For an integer \(a\ge2\), define

\[
L_a(x)=\prod_{n=1}^{a-1}\frac{x-x_n}{x_a-x_n},\qquad
C_a=L_a(1)=\prod_{n=1}^{a-1}\frac{1-x_n}{x_a-x_n}>0,
\] \[
r_a=\frac{1-x_{a+1}}{1-x_a}
=\frac{a^4\bigl(2(a+1)^2-1\bigr)}{(a+1)^4(2a^2-1)}\in(0,1),
\] \[
f_{a,m}(x)=L_a(x)\left(\frac{1-x}{1-x_a}\right)^m,
\qquad m=0,1,2,\ldots.
\tag{2}
\]

The degree is \(a-1+m\); all lower nodes are killed exactly;
\(f_{a,m}(x_a)=1\). At every higher node,
\(|f_{a,m}(x_n)|\le C_a r_a^m\), because each factor of \(L_a\)
increases on \([x_a,1]\) and \(1-x_n\le1-x_{a+1}\).

Set \(C_F=\prod_{q\in F}\max(1,|1+\lambda_q|)\). Then
\(|\kappa_F(n)|\le C_F\), and the elementary decreasing-series bound
gives

\[
\sum_{n>a}|\mu(n)^2n^{-4}\kappa_F(n)|
\le\frac{C_F}{3a^3}.
\tag{3}
\]

\textbf{Theorem 2.} If \(\lambda_p<-1\), then

\[
\int f_{p,m}(x)^2\,d\nu_{F,\boldsymbol\lambda}(x)
\le\frac{1+\lambda_p}{p^4}
+\frac{C_F C_p^2}{3p^3}r_p^{2m}.
\tag{4}
\]

Hence the Hankel matrix of order \(p+m\) is not positive semidefinite
whenever

\[
\boxed{\frac{p C_F C_p^2}{3}r_p^{2m}<-(1+\lambda_p).}
\tag{5}
\]

All quantities in (5) are rational when the strengths are rational. It
provides an exact finite certificate without calculating Hankel
eigenvalues or relying on floating-point root finding. A suitable finite
\(m\) exists because \(r_p<1\).

\textbf{Proof.} The polynomial kills all atoms below \(p\), equals 1 at
the negative atom \(p\), and has the stated uniform bound above \(p\).
Replacing all higher signed contributions by their absolute values and
applying (3) proves (4). The coefficients of \(f_{p,m}\) then supply a
negative quadratic form for \([M_{i+j}]_{i,j<p+m}\) when (5) holds.
\(\square\)

For depth \(k\), the same polynomial satisfies

\[
\int x^k f_{p,m}(x)^2\,d\nu
\le\frac{(1+\lambda_p)x_p^k}{p^4}
+\frac{C_F C_p^2}{3p^3}r_p^{2m}.
\]

Replace the right side of (5) by \(-(1+\lambda_p)x_p^k\) to obtain an
explicit depth-\(k\) detector. This proof extends to any finite signed
atomic measure with these nodes and the same summable weight envelope.
Its mechanism is an isolating polynomial followed by exponential
suppression of the higher nodes.

\hypertarget{conditional-prime-fibers-and-all-prime-margin-bounds}{%
\subsection{Conditional prime fibers and all-prime margin
bounds}\label{conditional-prime-fibers-and-all-prime-margin-bounds}}

Fix \(p\in F\), with all other strengths at least \(-1\), and write
\(\kappa(n)=\prod_{q\in F\setminus\{p\}}(1+\lambda_q\mathbf1_{q\mid n})\).
Separate the positive measure into its \(p\)-free and \(p\)-divisible
parts:

\[
A_m=[\sum_{p\nmid n}\mu(n)^2n^{-4}\kappa(n)x_n^{i+j}]_{i,j<m},
\quad
B_m=[\sum_{p\mid n}\mu(n)^2n^{-4}\kappa(n)x_n^{i+j}]_{i,j<m}.
\]

Both matrices are positive definite. For \(A_m\), use infinitely many
primes outside \(F\); for \(B_m\), use infinitely many squarefree
integers \(p\ell\), with primes \(\ell\) outside \(F\). This remains
true at every deletion endpoint of the other strengths.

\textbf{Theorem 3.} The one-variable fiber

\[
\tau_m(\lambda_p)=\det\bigl(A_m+(1+\lambda_p)B_m\bigr)
\]

has exact degree \(m\), and all its roots, with multiplicity, are real
and strictly less than \(-1\). The root nearest the positivity boundary
is \(-1-\delta_m\), with

\[
\delta_m=\min_{0\ne f\in\mathbb R[x],\ \deg f<m}
\frac{\int f^2\,d\nu_{p\text{-free}}}{\int f^2\,d\nu_{p\text{-divisible}}}>0.
\tag{6}
\]

The sequence \(\delta_m\) is nonincreasing, and, with
\(C_\kappa=\prod_{q\ne p}\max(1,1+\lambda_q)\),

\[
\boxed{\delta_{p+m}\le\frac{p C_\kappa C_p^2}{3}r_p^{2m}.}
\tag{7}
\]

\textbf{Proof.} The definite pencil is congruent to
\(I+(1+\lambda_p)A_m^{-1/2}B_mA_m^{-1/2}\), whose eigenvalues have
positive slopes. Its roots are \(-1-1/\eta_i\), with \(\eta_i>0\). This
proves the degree and root statements, and the reciprocal largest
eigenvalue is exactly (6). Polynomial-space inclusion proves
monotonicity. Test (6) at order \(p+m\) with (2). The denominator
includes the atom \(p\), of mass \(p^{-4}\), because \(\kappa(p)=1\).
The numerator kills every lower atom, has no atom at \(p\), and its
higher tail is bounded by \(C_\kappa C_p^2r_p^{2m}/(3p^3)\). This proves
(7) and therefore \(\delta_m\to0\). \(\square\)

A sharper arithmetic bound uses the exact remaining tail mass

\[
A_{\rm tail}=D(4)\frac{p^4}{p^4+1}
\prod_{q\ne p}\left(1+\frac{\lambda_q}{q^4+1}\right)
-\sum_{1\le n<p}\mu(n)^2n^{-4}\kappa(n)>0,
\]

and replaces (7) by \(\delta_{p+m}\le p^4 C_p^2 A_{\rm tail}r_p^{2m}\).
Every finite-degree coefficient can have poles only at the fiber's
Hankel roots; numerator cancellations can remove a possible pole. The
theorem does not assert simple roots or global multivariate real
stability.

At depth zero, the initial margin is exactly \(\delta_1=p^4\),
independently of all other safe strengths: the total \(p\)-free and
\(p\)-divisible masses have ratio \(p^4\), since all the other Euler
factors cancel. The higher-order margins contain the arithmetic overlap
information.

The same argument applies to every nonnegative depth, replacing each
integral by its \(x^k\) tilt. The denominator's explicit atom
contribution becomes \(p^{-4}x_p^k\), and (7) acquires the factor
\(x_p^{-k}\).

\hypertarget{resolution-nonconvex-finite-tests-and-a-different-additive-family}{%
\subsection{Resolution, nonconvex finite tests, and a different additive
family}\label{resolution-nonconvex-finite-tests-and-a-different-additive-family}}

For \(F=\{2,3\}\), the order-one Hankel matrix is the scalar

\[
M_0(\lambda_2,\lambda_3)=105\pi^{-4}
\left(1+\frac{\lambda_2}{17}\right)
\left(1+\frac{\lambda_3}{82}\right).
\]

Thus \((0,0)\) and \((-34,-164)\) have the same positive order-one
Hankel matrix. Their midpoint \((-17,-82)\) has zero total mass. The
unrestricted real parameter set passing the order-one
positive-definiteness test is therefore nonconvex. This is a concrete
warning about a finite test, while the true positivity domain in Theorem
1 is a convex closed orthant.

The positive-semidefinite feasible set is nonconvex as well: the points
\((0,0)\) and \((-34,-246)\) have positive mass factors 1 and 2, while
their convex combination with weights \(3/5\) and \(2/5\), namely
\((-68/5,-492/5)\), has mass factor \(-1/25\).

An additive prime family is a different deformation:

\[
\widetilde\nu_{F,\boldsymbol\gamma}
=\sum_{n\ge1}\mu(n)^2n^{-4}
\left(1+\sum_{p\in F}\gamma_p\mathbf1_{p\mid n}\right)\delta_{x_n}.
\]

Its exact positivity domain is the polyhedron

\[
1+\sum_{p\in S}\gamma_p\ge0\quad(S\subseteq F),
\quad\text{equivalently}\quad
1+\sum_{\gamma_p<0}\gamma_p\ge0.
\]

Necessity uses the squarefree integer \(n=\prod_{p\in S}p\); sufficiency
checks every atom's prime subset. Its finite Hankel matrices are affine
functions of the strengths, so their positive-semidefinite feasible sets
are convex spectrahedra. In contrast, the multiplicative family includes
the arithmetic prime-overlap terms and has polynomial parameter
dependence. For example, additive strengths \(\gamma_2=\gamma_3=-3/4\)
leave both individual prime weights positive but make the \(n=6\) weight
negative; multiplicative strengths \(\lambda_2=\lambda_3=-3/4\) remain
safely positive.

\hypertarget{a-limited-nonreduction-statement}{%
\subsection{A limited nonreduction
statement}\label{a-limited-nonreduction-statement}}

If at least two strengths in a positive multiplicative family are
nonzero, its entire depth-zero monic polynomial sequence cannot equal
any one-prime sequence on the same variable and node system. A direct
proof avoids importing a determinacy theorem. Write
\(P_m=x^m+\sum_{j<m}c_{mj}x^j\). The equations \(\int P_m\,d\nu=0\)
recursively determine every normalized moment from the monic sequence.
Equal sequences therefore give equal normalized Stieltjes series for
\(|z|>1\). Analytic continuation on the connected complement of the
common closed node set makes their Stieltjes functions equal, and
residues at each isolated node give equal normalized atom weights. The
mass 1 at zero fixes the normalization constant. Finally, the atom at
each prime \(p\) reads off \(1+\lambda_p\), whereas a one-prime
deformation changes at most one such atom. This is a precise separator
from the existing one-prime family, not a global novelty certificate. It
does not address arbitrary affine variable changes.

The multi-prime Stieltjes kernel also retains infinitely many isolated
poles at primes outside \(F\), so the preceding one-prime proof of
non-D-finiteness applies unchanged. Orthogonal-polynomial existence,
definite-pencil spectral theory, and compact moment determinacy are
classical mechanisms. The specific arithmetic coordinates, parameter
geometry, and explicit finite-resolution detector are the objects
developed here.

\hypertarget{worked-arithmetic-examples}{%
\subsection{Worked arithmetic
examples}\label{worked-arithmetic-examples}}

For \(F=\{2,3\}\), the following coefficients are rounded high-precision
evaluations of (1) and the Hankel formulas. They are not exact decimal
constants.

\begin{longtable}[]{@{}
  >{\raggedright\arraybackslash}p{(\columnwidth - 4\tabcolsep) * \real{0.3333}}
  >{\raggedright\arraybackslash}p{(\columnwidth - 4\tabcolsep) * \real{0.3333}}
  >{\raggedright\arraybackslash}p{(\columnwidth - 4\tabcolsep) * \real{0.3333}}@{}}
\toprule\noalign{}
\begin{minipage}[b]{\linewidth}\raggedright
\((\lambda_2,\lambda_3)\)
\end{minipage} & \begin{minipage}[b]{\linewidth}\raggedright
\(P_1(x)\)
\end{minipage} & \begin{minipage}[b]{\linewidth}\raggedright
\(P_2(x)\)
\end{minipage} \\
\midrule\noalign{}
\endhead
\bottomrule\noalign{}
\endlastfoot
\((-1,-1)\) & \(x-0.00200014863265043361\) &
\(x^2-0.933902261393426053x+0.000000921271332319859170\) \\
\((-1/2,1)\) & \(x-0.0377167661732468061\) &
\(x^2-0.725293001442481628x+0.000932648146085806112\) \\
\((1,1)\) & \(x-0.0822907496962840844\) &
\(x^2-0.655873200757636623x+0.00211153000026086129\) \\
\end{longtable}

Simultaneously deleting \(2,3,5\) gives
\(P_1(x)=x-0.000528721422047782375\) and
\(P_2(x)=x^2-0.966491641778821472x+0.0000000795531413796910514\), again
rounded. This removes the low prime shelves while retaining the fixed
mass at zero and an infinite tail.

For the variable strength at \(p=2\), the nearest conditional margins
are:

\begin{longtable}[]{@{}
  >{\raggedright\arraybackslash}p{(\columnwidth - 4\tabcolsep) * \real{0.2727}}
  >{\raggedleft\arraybackslash}p{(\columnwidth - 4\tabcolsep) * \real{0.3636}}
  >{\raggedleft\arraybackslash}p{(\columnwidth - 4\tabcolsep) * \real{0.3636}}@{}}
\toprule\noalign{}
\begin{minipage}[b]{\linewidth}\raggedright
Hankel order
\end{minipage} & \begin{minipage}[b]{\linewidth}\raggedleft
Other strength \(\lambda_3=-1\)
\end{minipage} & \begin{minipage}[b]{\linewidth}\raggedleft
Other strength \(\lambda_3=3\)
\end{minipage} \\
\midrule\noalign{}
\endhead
\bottomrule\noalign{}
\endlastfoot
1 & \(16\), exact & \(16\), exact \\
2 & \(0.0939109919936561\) & \(1.43829172807931\) \\
3 & \(0.000302450514004102\) & \(0.0346655628067259\) \\
4 & \(0.000000528145808944195\) & \(0.000223749375519984\) \\
\end{longtable}

The noninteger entries are approximate symmetric-pencil calculations.
They show an actual prime-overlap effect beyond the mass test. In these
calculations the defect \(\lambda_2=-1.001\) is visible at order 3 when
the 3-shelf is deleted, but only at order 4 when that shelf is
strengthened to 3. This comparison is a numerical example, not a
universal monotonicity theorem.

The separate rational certificates (5) do not depend on those
approximate roots. The script constructs negative detectors for primes
\(2,3,5,7\), at Hankel orders \(8,11,19,40\), respectively. These
sufficient orders are not asserted to be minimal. In the \(p=7\)
example, the other strengths are
\(\lambda_2=1/2,\lambda_3=-1,\lambda_5=2\), and \(\lambda_7=-1.0001\);
the exact rational upper bound on the detector norm is approximately
\(-9.74271\times10^{-9}\).

\hypertarget{computation-contract}{%
\subsection{Computation contract}\label{computation-contract}}

\path{checks/multi_prime_sieve.py} uses exact rational arithmetic
for Euler-factor versus subset-expansion comparisons, Bernoulli-based
even-zeta coordinates, isolating-polynomial values, and the rational
detection criterion (5). It separately evaluates the infinite moments at
high precision, tests positive orthogonality grids, compares them with a
direct finite squarefree sum plus a proved tail envelope, calculates
conditional definite-pencil roots, and evaluates signed detector norms.
These decimal checks are not interval certificates. Universal positivity
and negativity conclusions follow from the displayed proofs and exact
rational bounds.

The passing run records 1,019 exact rational assertions, 102 positive
parameter/depth cases, 1,395 lower-monomial orthogonality conditions,
381 recurrence identities, 40 conditional-pencil cases, and four exact
negative-atom certificates. Working precision is 320 decimal digits.
Maximum absolute orthogonality residual was approximately
\(6.33\times10^{-322}\), maximum recurrence residual
\(2.23\times10^{-303}\), and maximum scale-normalized pencil determinant
residual \(5.06\times10^{-321}\). These are finite numerical consistency
results, not interval certificates. No execution failure occurred in
this route's run.

The script records actual versions, source hash, parameter bounds,
finite counts, precision, residuals, examples, and any failed criteria
in \path{evidence/multi_prime_sieve_checks.json}. It accepts an
optional output path so a parent bounded runner can place the result
inside a fresh evidence directory. No broad priority search was
performed for this specific extension; historical novelty remains
unresolved.

\clearpage\hypertarget{orbit-spectral-zeta-continuation-trace-zeros-and-regularized-product}{%
\section{Orbit spectral zeta: continuation, trace zeros, and regularized
product}\label{orbit-spectral-zeta-continuation-trace-zeros-and-regularized-product}}

This chapter extends the previously defined orbit family with a self-contained continuation proof.
The Mittag-Leffler seed and Mellin continuation mechanism are classical.
Historical priority for this exact specialization is not established.

\hypertarget{the-completed-result}{%
\subsection{The completed result}\label{the-completed-result}}

Let \(N\ge2\) be an integer, and put \(\alpha=1/N\). Define

\[
d_N(w)=\sum_{n=0}^{\infty}\frac{(-1)^n w^n}{(Nn)!}
      =E_{N,1}(-w).
\]

List its nonzero zeros \(w_j\) with multiplicity, and use the principal
logarithm, with \(-\pi<\operatorname{Arg}w_j<\pi\). No zero lies on the
omitted negative ray. The orbit zeta function

\[
\mathcal Z_N(s)=\sum_j e^{-s\operatorname{Log}w_j}
\]

converges absolutely exactly when \(\Re s>1/N\). It has a meromorphic
continuation to the whole plane, with exactly one pole: a simple pole at
\(s=1/N\), of residue

\[
\boxed{\operatorname{Res}_{s=1/N}\mathcal Z_N(s)
      =\frac{\sin(\pi/N)}{N\pi}.}
\]

At every nonpositive integer,

\[
\boxed{\mathcal Z_N(-m)=0\quad(m=0,1,2,\ldots).}
\]

At the origin,

\[
\boxed{\mathcal Z_N'(0)=-\log N,\qquad
       \exp[-\mathcal Z_N'(0)]=N.}
\]

The last equality is the zeta regularized product of the orbit zeros
under this branch convention. It is a regularized product, not
convergence of the ordinary product of the zeros. Positive integer values retain the exact rational coefficient law established in the preceding construction.

\hypertarget{the-exponential-packet-and-its-exact-asymptotic}{%
\subsection{The exponential packet and its exact
asymptotic}\label{the-exponential-packet-and-its-exact-asymptotic}}

Let \(\omega=e^{2\pi i/N}\). The finite root-of-unity filter gives, for
every \(x\in\mathbb C\),

\[
F_N(x):=d_N(-x^N)
=\frac1N\sum_{j=0}^{N-1}e^{\omega^j x}
=\sum_{n=0}^{\infty}\frac{x^{Nn}}{(Nn)!}.
\tag{S1}
\]

This equality follows by expanding the exponentials and using
\(N^{-1}\sum_j\omega^{jm}=1\) when \(N\mid m\), and zero otherwise. For
real \(x\ge0\), the power series makes \(F_N(x)>0\). Put

\[
h_N(x)=\sum_{j=1}^{N-1}e^{(\omega^j-1)x},
\qquad c_N=1-\cos(2\pi/N)>0.
\]

Then

\[
F_N(x)=\frac{e^x}{N}(1+h_N(x)),\qquad
|h_N(x)|\le(N-1)e^{-c_Nx},\qquad
|h_N'(x)|\le2(N-1)e^{-c_Nx}.
\tag{S2}
\]

Since \(1+h_N(x)=Ne^{-x}F_N(x)>0\), the real logarithm is well defined
for all \(x\ge0\). At infinity,

\[
\log F_N(x)=x-\log N+O(e^{-c_Nx}),
\qquad g_N(x):=\frac{F_N'(x)}{F_N(x)}
=1+O(e^{-c_Nx}).
\tag{S3}
\]

At the origin \(g_N(x)=O(x^{N-1})\). Its only Taylor powers are
\(x^{Nm-1}\), with \(m\ge1\).

\hypertarget{the-branch-and-absolute-interchange-issue}{%
\subsection{The branch and absolute-interchange
issue}\label{the-branch-and-absolute-interchange-issue}}

The pre-existing genus-zero product is

\[
d_N(w)=\prod_j(1-w/w_j).
\tag{S4}
\]

For completeness, the series bound \(|d_N(w)|\le\exp(|w|^{1/N})\) gives
order at most \(1/N<1\). Hadamard factorization therefore has genus zero
and a constant exponential factor, which equals one because
\(d_N(0)=1\). The usual zero-count consequence gives
\(\sum_j|w_j|^{-\sigma}<\infty\) for every \(\sigma>1/N\). Along the
negative real axis (S3) also shows the order is exactly \(1/N\).

It is not enough simply to say that no zero lies on the integration
contour. We need a uniform angular gap from the branch cut to
interchange the Mellin integral with the zero sum. Here is that gap.

For \(|\arg(-w)|\le\pi/2\), choose \(x=(-w)^{1/N}\) with
\(|\arg x|\le\pi/(2N)\). For every \(j\ne0\),

\[
\Re x-\Re(\omega^j x)
\ge \left[\cos\frac{\pi}{2N}-\cos\frac{3\pi}{2N}\right]|x|
=:\gamma_N|x|>0.
\tag{S5}
\]

Indeed, the nearest competing root is separated from the positive
direction by at least \(3\pi/(2N)\), while the distinguished root is
within \(\pi/(2N)\). Cosine decreases on \([0,\pi]\). Thus, once
\((N-1)e^{-\gamma_N|x|}<1\), the distinguished term \(e^x\) in (S1) is
larger in modulus than the sum of all competitors, and \(d_N(w)\ne0\).
Consequently, sufficiently large zeros avoid the entire closed left
half-plane. There are only finitely many remaining zeros in a bounded
disk. None lies on the negative ray, since \(d_N(-t)>0\) for \(t\ge0\).
It follows that some \(\eta_N>0\) satisfies

\[
|\operatorname{Arg}w_j|\le\pi-\eta_N\qquad\text{for every }j.
\tag{S6}
\]

With \(c=\sin(\eta_N/2)>0\), this implies

\[
|t+w_j|\ge c(t+|w_j|),\qquad t\ge0.
\tag{S7}
\]

To check the elementary estimate, write \(w=re^{i\theta}\):
\(|t+w|^2=(t+r)^2-4tr\sin^2(\theta/2)\ge(t+r)^2\cos^2(\theta/2)\).

\hypertarget{mellin-representation-in-the-original-strip}{%
\subsection{Mellin representation in the original
strip}\label{mellin-representation-in-the-original-strip}}

Let \(H_N(t)=d_N(-t)=E_{N,1}(t)\), and \(G_N(t)=H_N'(t)/H_N(t)\).
Logarithmic differentiation of (S4), locally uniformly off the zeros,
gives

\[
G_N(t)=\sum_j\frac1{t+w_j},\qquad t\ge0.
\tag{S8}
\]

For \(1/N<\sigma<1\), (S7) yields

\[
\sum_j\int_0^\infty\frac{t^{-\sigma}}{|t+w_j|}\,dt
\le\frac{\pi}{c\sin(\pi\sigma)}\sum_j|w_j|^{-\sigma}<\infty.
\tag{S9}
\]

This is the absolute interchange needed for complex \(s\) with
\(\Re s=\sigma\). Euler's beta integral and the reflection identity give

\[
\int_0^\infty\frac{t^{-s}}{t+w}\,dt
=\frac{\pi}{\sin(\pi s)}e^{-s\operatorname{Log}w},
\quad 0<\Re s<1,\quad |\arg w|<\pi.
\tag{S10}
\]

First prove this for positive \(w\) by substitution and the beta
integral, then extend to the slit \(w\)-plane by holomorphic dependence.
On compact subsets of that slit plane the denominator estimate supplies
the needed local domination.

Combining (S8)--(S10), and substituting \(t=x^N\), proves

\[
\boxed{\mathcal Z_N(s)=\frac{\sin(\pi s)}{\pi}
\int_0^\infty x^{-Ns}g_N(x)\,dx,
\qquad \frac1N<\Re s<1.}
\tag{S11}
\]

For \(x>0\), \(x^{-Ns}=e^{-Ns\log x}\) uses the real logarithm.
Multiplicities are already incorporated in (S4), so no simple-zero
assumption entered the proof.

\hypertarget{global-continuation-and-the-only-pole}{%
\subsection{Global continuation and the only
pole}\label{global-continuation-and-the-only-pole}}

One can see directly that \(d_N\) has no zero on \(|w|\le1\):

\[
|d_N(w)-1|\le\sum_{n\ge1}\frac1{(Nn)!}
\le\cosh(1)-1<1.
\]

The strict inequality also gives a slightly larger zero-free disk by
continuity. Thus \(G_N(t)\) has a Taylor series converging uniformly on
\([0,1]\),

\[
G_N(t)=\sum_{k\ge0}a_{N,k}t^k,
\qquad g_N(x)=\sum_{m\ge1}N a_{N,m-1}x^{Nm-1}.
\tag{S12}
\]

For any integer \(K\ge0\), define

\[
\begin{aligned}
B_{N,K}(s)={}&
\int_0^1x^{-Ns}\left[g_N(x)-\sum_{m=1}^{K}N a_{N,m-1}x^{Nm-1}\right]dx\\
&+\sum_{m=1}^{K}\frac{a_{N,m-1}}{m-s}
 +\int_1^\infty x^{-Ns}\bigl[g_N(x)-1\bigr]dx
 +\frac1{Ns-1}.
\end{aligned}
\tag{S13}
\]

The first integral is holomorphic for \(\Re s<K+1\). This includes local
uniform domination, since the Taylor remainder is \(O(x^{N(K+1)-1})\).
The last integral is entire in \(s\), because (S3) gives exponential
decay against every power of \(x\), uniformly on compact \(s\)-sets,
including after any number of \(s\)-derivatives. The rational terms
exhibit the only possible poles: positive integers \(1,\ldots,K\) and
\(1/N\).

On their overlapping domains the functions (S13) agree: adding the next
Taylor term subtracts its integral and adds precisely the same rational
function. They also agree with (S11) in its original strip. Therefore

\[
\mathcal Z_N(s)=\frac{\sin(\pi s)}{\pi}B_{N,K}(s)
\tag{S14}
\]

provides compatible meromorphic continuations covering the whole plane
as \(K\to\infty\). At each positive integer \(m\), the apparent simple
pole \(a_{N,m-1}/(m-s)\) cancels with the simple zero of
\(\sin(\pi s)\). The term \(1/(Ns-1)\) has residue \(1/N\), and its
multiplier at \(s=1/N\) is nonzero. This proves the stated unique simple
pole and its residue.

It also proves that the original series cannot converge absolutely at
\(s=1/N\). If it did, its absolute terms would dominate the series for
real \(s\downarrow1/N\), except for finitely many zeros of modulus below
one. Dominated convergence would then keep that real limit bounded,
contradicting the nonzero pole. Absolute convergence at any smaller real
part would imply convergence at \(1/N\). The abscissa is consequently
exactly \(1/N\).

\hypertarget{trace-zeros-rational-values-and-the-origin}{%
\subsection{Trace zeros, rational values, and the
origin}\label{trace-zeros-rational-values-and-the-origin}}

For \(\Re s<1\) we may use \(K=0\):

\[
B_{N,0}(s)=\int_0^1x^{-Ns}g_N(x)\,dx
+\int_1^\infty x^{-Ns}(g_N(x)-1)\,dx+\frac1{Ns-1}.
\tag{S15}
\]

This expression is holomorphic at every nonpositive integer. Hence (S14)
proves \(\mathcal Z_N(-m)=0\). These statements do not assert that every
such zero is simple.

At a positive integer \(m\), taking the removable limit in (S14) gives

\[
\mathcal Z_N(m)=(-1)^{m+1}a_{N,m-1}.
\tag{S16}
\]

If \(c_n=(-1)^n/(Nn)!\), the previous logarithmic-product law therefore
remains

\[
b_m:=\mathcal Z_N(m)
=-m c_m-\sum_{j=1}^{m-1}c_jb_{m-j}.
\tag{S17}
\]

In particular all these values are rational. This is the finite
coefficient dictionary, independent of numerical zero location.

At \(s=0\), the elementary antiderivatives in (S15) give

\[
\begin{aligned}
B_{N,0}(0)
&=\log F_N(1)
 +\lim_{R\to\infty}\bigl[\log F_N(R)-R-\log F_N(1)+1\bigr]-1\\
&=-\log N.
\end{aligned}
\]

Since \(\sin(\pi s)/\pi=s+O(s^3)\), (S14) implies \(\mathcal Z_N(0)=0\)
and \(\mathcal Z_N'(0)=-\log N\).

\hypertarget{classical-benchmarks-and-numerical-examples}{%
\subsection{Classical benchmarks and numerical
examples}\label{classical-benchmarks-and-numerical-examples}}

The inherited identities

\[
\mathcal Z_2(s)=\pi^{-2s}\zeta(2s,1/2),\qquad
\mathcal Z_4(s)=(\sqrt2\pi)^{-4s}\zeta(4s,1/2)
\]

now hold as meromorphic identities, by equality in the original
half-plane and uniqueness of continuation. Their pole residues are
respectively \(1/(2\pi)\) and \(1/(4\sqrt2\pi)\), agreeing with the
general theorem. Their zero values at nonpositive integers follow from
the Hurwitz formula, and their regularized products are respectively 2
and 4.

The companion script performs 411 exact rational assertions: independent
formal-logarithm and quotient calculations for \(2\le N\le10\) through
degree 40; three explicit integer values for each \(N\); and 12
Bernoulli/Hurwitz anchors each for \(N=2,4\). It then evaluates (S13)
using a truncated local Taylor sum and high-precision exponential-tail
quadrature. Sixteen \(N=2,4\) benchmarks sample both sides of the
original convergence boundary, negative real values, and complex values.
Values for \(N=3,5\) and checks of \(\mathcal Z_N'(0)\) are recorded in
the JSON evidence.

Those quadrature results are approximations. Agreement at two precisions
and with known functions tests normalization and implementation; it does
not replace the proof, certify intervals, or settle historical novelty.

\hypertarget{a-worked-hurwitz-deformation}{%
\subsection{A worked Hurwitz
deformation}\label{a-worked-hurwitz-deformation}}

There is a concrete shifted family. For a real shift \(a\ge0\), define

\[
\mathcal Z_N(s;a)=\sum_j(w_j+a)^{-s},\qquad \Re s>1/N,
\tag{T1}
\]

again with the principal logarithm and multiplicities. A rightward shift
preserves the angular gap from the negative ray: for each nonreal \(w\),
\(|\arg(w+a)|\) is nonincreasing in \(a\), and positive real zeros
remain positive. The zero sequence is unchanged at infinity up to a
bounded translation, so the previous absolute convergence argument
applies. The normalized determinant is exactly

\[
\frac{H_N(a+t)}{H_N(a)}
=\prod_j\left(1+\frac{t}{w_j+a}\right),\qquad
H_N(t)=E_{N,1}(t)>0\quad(t\ge0).
\tag{T2}
\]

This follows by taking the ratio of the two convergent products in (S4),
not by introducing an unverified exponential factor. Its logarithmic
derivative is \(G_{N,a}(t)=H_N'(a+t)/H_N(a+t)\). The same dominated
beta-integral argument gives

\[
\mathcal Z_N(s;a)=\frac{\sin(\pi s)}{\pi}
\int_0^\infty t^{-s}G_{N,a}(t)\,dt,
\qquad 1/N<\Re s<1.
\tag{T3}
\]

At infinity, (S2)--(S3), with \(x=(a+t)^{1/N}\), imply

\[
G_{N,a}(t)=\alpha(t+a)^{\alpha-1}
+O\!\left((t+a)^{\alpha-1}e^{-c_N(t+a)^\alpha}\right).
\tag{T4}
\]

For fixed \(a\) and any integer \(K\ge0\), the binomial expansion
therefore gives

\[
G_{N,a}(t)=\sum_{k=0}^{K}\alpha\binom{\alpha-1}{k}a^k t^{\alpha-k-1}
+O(t^{\alpha-K-2}).
\tag{T5}
\]

To be explicit about continuation, choose a finite splitting point
\(T>2a+1\). Subtract the first \(K+1\) terms of (T5) inside the integral
from \(T\) to infinity, and add their integrals back as the rational
terms

\[
\frac{\alpha\binom{\alpha-1}{k}a^k
 T^{\alpha-k-s}}{s-\alpha+k}.
\tag{T6}
\]

The remainder integral is holomorphic for \(\Re s>\alpha-K-1\). The
integral on any compact interval away from zero is entire. At zero,
\(G_{N,a}\) is analytic since \(H_N(a)>0\); subtracting any finite
Taylor polynomial there extends the other side, with apparent poles only
at positive integers. Such Taylor poles cancel against \(\sin(\pi s)\),
just as in Section 4. Choosing arbitrarily large subtractions at both
ends covers every \(s\in\mathbb C\). The resulting continuations agree
on overlaps because all added terms are exactly the integrals of the
terms removed.

For \(a>0\), the complete pole set is consequently

\[
\boxed{s=\frac1N-k,\quad k=0,1,2,\ldots,}
\]

and every pole is simple. Its residue is

\[
\begin{aligned}
\operatorname{Res}_{s=\alpha-k}\mathcal Z_N(s;a)
&=\frac{\alpha\sin(\pi(\alpha-k))}{\pi}
  \binom{\alpha-1}{k}a^k\\
&=\boxed{\frac{\sin(\pi/N)}{N\pi}
          \frac{(1-1/N)_k}{k!}a^k.}
\end{aligned}
\tag{T7}
\]

Here \((b)_0=1\) and \((b)_k=b(b+1)\cdots(b+k-1)\). The sign calculation
uses both \(\sin(\pi(\alpha-k))=(-1)^k\sin(\pi\alpha)\) and
\(\binom{\alpha-1}{k}=(-1)^k(1-\alpha)_k/k!\); their signs cancel. All
residues are positive when \(a>0\). At \(a=0\), the secondary
coefficients vanish, so the only remaining pole is the unshifted pole at
\(1/N\). As before, its nonzero residue makes the absolute convergence
abscissa exactly \(1/N\).

None of the pole locations \(\alpha-k\) is an integer. Thus the sine
factor vanishes at every nonpositive integer without an opposing pole:

\[
\boxed{\mathcal Z_N(-m;a)=0\quad(m\ge0).}
\tag{T8}
\]

Finally, the one-term subtraction is already convergent at \(s=0\),
since \(G_{N,a}(t)-\alpha t^{\alpha-1}=O(t^{\alpha-2})\) and
\(\alpha<1\). Let

\[
B_{N,a}(0)=\int_0^1G_{N,a}(t)dt
 +\int_1^\infty[G_{N,a}(t)-\alpha t^{\alpha-1}]dt-1.
\]

Integrating exactly and using \(\log H_N(a+t)-t^\alpha\to-\log N\) gives

\[
B_{N,a}(0)=-\log N-\log H_N(a).
\]

The multiplier \(\sin(\pi s)/\pi\) then gives

\[
\boxed{\mathcal Z_N'(0;a)=-\log[N E_{N,1}(a)],\qquad
       \exp[-\mathcal Z_N'(0;a)]=N E_{N,1}(a).}
\tag{T9}
\]

For \(N=2\), the product reduces to \(2\cosh\sqrt a\). Four direct
quadrature checks at \(a=0,1/4,2,9\) agree with \(-\log(2\cosh\sqrt a)\)
at 80-digit working precision. They use a local integral, an exactly
integrated algebraic leading tail, and an exponentially decaying
remainder; the quadrature does not sample numerical zero locations.
These are four additional numerical benchmarks, with the original 411
exact assertions unchanged. Results are in
\path{evidence/spectral_continuation_checks.json} under
\path{shifted_N2_product_benchmarks}.

This is a classical Hurwitz-type spectral deformation, with a
packet-specific explicit regularized determinant. It is useful because a
single shift reveals an entire pole ladder that collapses at the
unshifted boundary. No historical-priority claim is made for the
deformation or determinant identity.

\hypertarget{source-comparison-and-bounded-search}{%
\subsection{Source comparison and bounded
search}\label{source-comparison-and-bounded-search}}

No external global-continuation theorem is required by the local proof above beyond the
finite-order Hadamard product and the elementary beta/reflection
identities. The Mellin subtraction mechanism itself is established.

\begin{itemize}
\tightlist
\item
  NIST DLMF 10.46.3: the definition of \(E_{a,b}\) identifies this seed
  with the classical Mittag-Leffler function. Source:
  \url{https://dlmf.nist.gov/10.46}.
\item
  NIST DLMF 5.12.1 and 5.5.3: Euler's beta integral and the gamma
  reflection identity supply (S10), with the project's exponents
  obtained by \(a=1-s,b=s\). Sources: \url{https://dlmf.nist.gov/5.12} and
  \url{https://dlmf.nist.gov/5.5}. The extension to complex \(w\) and the
  uniform branch-gap argument are supplied locally above.
\item
  Noam Elkies, finite-order-function teaching notes, first Hadamard
  theorem: supplies the genus-zero product for an entire function of
  order below one. Source:
  \url{https://people.math.harvard.edu/~elkies/M259.02/prod.pdf}.
  The constant factor and zero-count consequences are established in the preceding sections.
\item
  Andre Voros, \textbf{Zeta functions for the Riemann zeros},
  \emph{Annales de l'Institut Fourier} 53(3) (2003), 665--699. The
  authenticated journal PDF, rather than an evolving preprint, was
  checked. Section 2.1 specifies general admissible zero sequences of
  order below one and logarithmic asymptotics; Section 2.2 derives
  global meromorphic continuation by Mellin subtraction,
  positive-integer cancellations, and origin trace constants. Source:
  \url{https://www.numdam.org/article/AIF_2003__53_3_665_0.pdf}. This
  authenticates the shared method. It does not state this particular
  Mittag-Leffler packet theorem verbatim. All hypotheses needed here,
  including the complex-zero branch and domination conditions, were
  established independently in Sections 1--4. Classification:
  \textbf{known mechanism after translation, with a worked seed-specific
  specialization}.
\item
  A separate Voros preprint, \url{https://arxiv.org/pdf/math/0311029} was
  compared as a secondary reference. For method attribution this monograph relies on the fixed 2003 journal article cited above.
\end{itemize}

The bounded discovery searches used both available search engines, with
the exact and adjacent terms ``Mittag Leffler zeros zeta function E\_n,1
analytic continuation spectral zeta log n'', ``hyper Bessel zeta
function analytic continuation Mittag Leffler zeros integer order sin pi
n'', ``Mittag-Leffler zeta zeros spectral meromorphic'', ``hyper-Bessel
zeta function continuation'', and ``E\_\{N,1\} zeros zeta''. The closest
authenticated result inspected was Voros's general framework. Search
results referring to the unrelated Mittag-Leffler theorem or Helson zeta
were discarded. Specialist book coverage and a complete backward/forward
citation search were not exhausted. No global-priority claim follows
from this search.

\hypertarget{scope-and-remaining-questions}{%
\subsection{Scope and remaining
questions}\label{scope-and-remaining-questions}}

The global continuation gap for this unweighted orbit family, and for
its nonnegative real shifts, is closed by the local proof. A functional
equation is not supplied or presumed. The only-pole theorem concerns the
unshifted case; positive shifts have the full pole ladder in (T7). The
pole and trace-zero theorems follow from the exact positive-axis
exponential asymptotic, and do not imply an RH-type critical-line
property. Residue-weighted sums still require separate multiplicity and
convergence analysis. A historical comparison against all Mittag-Leffler
and hyper-Bessel zero-zeta literature remains open.

No mathematical counterexample was found during this route. Any
computational failure and its repair, if encountered, must remain
recorded in the evidence history rather than silently erased.

\clearpage\begingroup\setlength{\parskip}{3pt}\hypertarget{exact-sequence-polyphase-and-tap-laboratory}{%
\section{Exact sequence, polyphase and tap
laboratory}\label{exact-sequence-polyphase-and-tap-laboratory}}



This chapter connects the nonlinear family to exact sequences and finite
filter algebra. It supplies reusable code without replacing the analytic
proofs in \path{routes/asymmetric_dilation.md}. All mathematical
sequence indices begin at zero. All finite arrays are zero extended
beyond their displayed support. The catalog contains explicit candidates
for OEIS lookup. This chapter does not assert an OEIS accession or
historical priority for an unsearched row.

\hypertarget{the-integer-polynomial-object}{%
\subsection{The integer polynomial
object}\label{the-integer-polynomial-object}}

Write the unique odd formal solution as

\[
T_{q,r}(z)=\sum_{n\ge0}c_n(q,r)z^{2n+1}
          =\sum_{n\ge0}A_n(q,r)\frac{z^{2n+1}}{(2n+1)!}.
\]

The two coefficient systems are related by \(A_n=(2n+1)!c_n\). In
particular, \(A_n\) is the derivative of order \(2n+1\) at zero. The
coefficient \(c_n\), the derivative \(A_n\), and the full EGF sequence
\(a_m\) with \(a_{2n+1}=A_n, a_{2n}=0\) are different objects. None of
the code silently substitutes one for another.

The nonlinear differential equation gives the recurrence

\[
\boxed{
A_0=1,\qquad
A_n=\sum_{i=0}^{n-1}\binom{2n}{2i+1}
 q^{2i+1}r^{2(n-1-i)+1}A_i A_{n-1-i},\quad n\ge1.}
\]

\textbf{Proof.} Multiply the ordinary Taylor recurrence by \((2n+1)!\).
The factor attached to \(A_iA_j\), where \(i+j=n-1\), is
\((2n)!/[(2i+1)!(2j+1)!]=\binom{2n}{2i+1}\). This also proves the
recurrence directly in the EGF convention. \(\square\)

\textbf{Polynomial theorem.} For every \(n\), \(A_n\in\mathbb Z[q,r]\),
all its nonzero coefficients are positive integers, and
\(A_n(q,r)=A_n(r,q)\). Every monomial \(q^u r^v\) satisfies

\[
u,v\ge n,\qquad u\equiv v\equiv n\pmod2,\qquad u+v\le n(n+1).
\]

\textbf{Proof.} The recurrence is an induction using positive integer
binomial coefficients. Swapping \(q,r\) and replacing \(i\) by \(n-1-i\)
leaves the sum unchanged. If child indices are \(i,j\), their exponent
pairs receive the shifts \((2i+1,2j+1)\). The lower bounds follow from
\(i+j+1=n\). The two parities are the parity of \(i+j+1\), also \(n\).
For the degree bound, a recurrence term has total degree at most

\[
2n+i(i+1)+j(j+1)\le n(n+1),\qquad i+j=n-1,
\]

since \(i(i+1)+j(j+1)\) is largest at \((i,j)=(0,n-1)\) or its swap. The
base \(A_0=1\) satisfies every assertion. \(\square\)

The first four polynomials are

\[
\begin{aligned}
A_0={}&1,\\
A_1={}&2qr,\\
A_2={}&8q^2r^2(q^2+r^2),\\
A_3={}&q^3r^3\bigl[48(q^2+r^2)(q^4+r^4)+80q^2r^2\bigr].
\end{aligned}
\]

Consequently each integer pair \((q,r)\) gives an integer sequence for
all \(n\), independently of whether its formal Taylor series converges.
For example, the first values of \(A_n(1,2)\) are
\(1,4,160,35200,38041600\).

\hypertarget{rational-feedback-banks-also-admit-integer-sequences}{%
\subsubsection{Rational feedback banks also admit integer
sequences}\label{rational-feedback-banks-also-admit-integer-sequences}}

For rational channel tuples \((h_t,q_t,r_t)\), use

\[
F'(z)=1+\sum_t h_tF(q_tz)F(r_tz),\qquad F(0)=0.
\]

The preceding recurrence is summed over channels with the factor
\(h_t\). In the universal polynomial over the weights and contractions,
every term of \(A_n\) has weight degree exactly \(n\) and total
contraction degree at most \(n(n+1)\). The weight statement follows from
one new weight at each recurrence step and \(1+i+j=n\); the contraction
bound uses the same induction.

Let \(H\) be a common denominator for all channel weights, and \(D\) a
common denominator for all channel contractions. Then

\[
\boxed{B_n=H^nD^{n(n+1)}A_n\in\mathbb Z\quad\hbox{for all }n.}
\]

This is a proved normalization, even with signed rational weights. It is
usually not the smallest possible denominator clearing. The catalog
records both \(A_n\) and \(B_n\), and the exact multiplier, so that a
proposed sequence lookup can state its normalization honestly.

\hypertarget{a-sharp-distinction-between-sequences-and-analytic-germs}{%
\subsection{A sharp distinction between sequences and analytic
germs}\label{a-sharp-distinction-between-sequences-and-analytic-germs}}

The entire-function theorem uses \(0<q,r<1\). Specializations with large
positive integers produce valid formal and integer sequences but do not
necessarily produce analytic functions near zero.

\textbf{Zero-radius theorem.} If \(q>0\) and \(r>1\), the formal series
for \(T_{q,r}\) has radius of convergence zero. The same holds after
swapping the parameters.

\textbf{Proof.} Every coefficient is positive. In the ordinary
recurrence keep only the term \(i=0\). This gives

\[
c_n\ge\frac{q r^{2n-1}}{2n+1}c_{n-1}
\quad\Longrightarrow\quad
c_n\ge\frac{q^n r^{n^2}}{3\cdot5\cdots(2n+1)}.
\]

Since the denominator is at most \((2n+1)^n\),

\[
c_n^{1/(2n+1)}\ge
q^{n/(2n+1)}r^{n^2/(2n+1)}(2n+1)^{-n/(2n+1)}\longrightarrow\infty.
\]

The Cauchy--Hadamard criterion therefore gives radius zero. \(\square\)

More generally, a positive feedback bank with a channel satisfying
\(h>0,q>0,r>1\) has the same obstruction, replacing \(q^n\) by
\((hq)^n\). Positivity prevents cancellation. No conclusion of this form
is asserted for a bank with signed weights.

\begin{longtable}[]{@{}
  >{\raggedright\arraybackslash}p{(\columnwidth - 4\tabcolsep) * \real{0.3333}}
  >{\raggedright\arraybackslash}p{(\columnwidth - 4\tabcolsep) * \real{0.3333}}
  >{\raggedright\arraybackslash}p{(\columnwidth - 4\tabcolsep) * \real{0.3333}}@{}}
\toprule\noalign{}
\begin{minipage}[b]{\linewidth}\raggedright
Parameters or bank
\end{minipage} & \begin{minipage}[b]{\linewidth}\raggedright
First five \(A_n\)
\end{minipage} & \begin{minipage}[b]{\linewidth}\raggedright
Mathematical status
\end{minipage} \\
\midrule\noalign{}
\endhead
\bottomrule\noalign{}
\endlastfoot
\((1,1)\) & \(1\), \(2\), \(16\), \(272\), \(7936\) & Known tangent derivatives \\
\((0,1)\) & \(1\), \(0\), \(0\), \(0\), \(0\) & Known linear boundary \\
\((1,-1)\) & \(1\), \(-2\), \(16\), \(-272\), \(7936\) & Known hyperbolic tangent
derivatives \\
\((1,2)\) & \(1\), \(4\), \(160\), \(35200\), \(38041600\) & Integer sequence, radius
zero \\
\((1/2,1/2)\) & \(1\), \(1/2\), \(1/4\), \(1/8\), \(1/16\) & Known scaled hyperbolic sine \\
\((1/4,3/4)\) & \(1\), \(3/8\), \(45/256\), \(21465/262144\), \(39833775/1073741824\) &
Entire asymmetric family, order one \\
\end{longtable}

The complete catalog contains \textbf{18 rows}, each through \(n=12\):
three classical or degenerate integer anchors, eight positive integer
formal specializations, four rational contractions, one entire edge case
of infinite order, and two normalized feedback banks. Integer
integrality and rational clearing are established by the preceding
proofs, not merely inferred from the finite sample.

\hypertarget{the-uncontracted-boundary-is-entire-of-infinite-order}{%
\subsubsection{The uncontracted boundary is entire of infinite
order}\label{the-uncontracted-boundary-is-entire-of-infinite-order}}

\textbf{Boundary theorem.} For \(0<q<1\), the normalized formal solution
of

\[
T'(z)=1+T(qz)T(z),\qquad T(0)=0
\]

extends to an entire function of infinite order. The swapped boundary
has the same property.

\textbf{Proof of entire extension.} The tangent-coefficient majorant
gives an initial holomorphic germ. Suppose \(T\) is known on \(|z|<R\).
The function \(T(qz)\) is known on \(|z|<R/q\). On this larger disk
solve the linear ODE \(U'=1+T(qz)U\), \(U(0)=0\). Explicitly,

\[
U(z)=e^{B(z)}\int_0^z e^{-B(w)}\,dw,
\qquad B(z)=\int_0^zT(qw)\,dw.
\]

The integrals are path independent on the disk. Linear-ODE uniqueness
shows \(U=T\) on the previous disk, and there \(T(qz)\) equals the newly
extended function at \(qz\). Thus the same nonlinear equation holds on
the enlarged disk. Iterating \(R\mapsto R/q\) exhausts the plane.

\textbf{Proof of infinite order.} Positive coefficients imply
\(T(x)\ge x\) for \(x>0\) and \(M(x)=T(x)\). The logarithmic derivative
obeys

\[
\frac{d}{dx}\log T(x)=\frac1{T(x)}+T(qx)\ge qx.
\]

Integrating from one, with \(T(1)\ge1\), gives
\(T(x)\ge\exp(q(x^2-1)/2)\) for \(x\ge1\). For \(R\ge2/q\), integrate
the same logarithmic derivative over \([R/2,R]\). Since \(T(R/2)\ge1\),

\[
\boxed{
\log M(R)\ge\frac R2\exp\left(\frac{q^3R^2}{8}-\frac q2\right).}
\]

Therefore \(\log\log M(R)/\log R\to\infty\), which proves infinite
order. \(\square\)

For positive parameters this gives a complete analytic/formal phase
chart:

\begin{longtable}[]{@{}
  >{\raggedright\arraybackslash}p{(\columnwidth - 2\tabcolsep) * \real{0.5000}}
  >{\raggedright\arraybackslash}p{(\columnwidth - 2\tabcolsep) * \real{0.5000}}@{}}
\toprule\noalign{}
\begin{minipage}[b]{\linewidth}\raggedright
Parameter region
\end{minipage} & \begin{minipage}[b]{\linewidth}\raggedright
Normalized solution
\end{minipage} \\
\midrule\noalign{}
\endhead
\bottomrule\noalign{}
\endlastfoot
\(0<q,r<1\) & Entire, finite order defined by \(q^\rho+r^\rho=1\) \\
One parameter equals one, the other lies in \((0,1)\) & Entire, infinite
order \\
\(q=r=1\) & Tangent, meromorphic with its usual poles \\
\(q>0,r>0,\max(q,r)>1\) & Formal solution with radius zero \\
One parameter equals zero & Linear solution \(z\), regardless of the
other scale \\
\end{longtable}

The finite-order formula applies in the open contraction square. It does
not assign a finite order to the edge, where its defining scale equation
has no positive root.

\hypertarget{exact-root-filter-sections}{%
\subsection{Exact root-filter
sections}\label{exact-root-filter-sections}}

For an EGF \(G(z)=\sum_{m\ge0}a_m z^m/m!\), an integer base \(b\ge1\), a
residue \(0\le j<b\), and \(\omega=e^{2\pi i/b}\), define

\[
\boxed{
\mathcal P_{b,j}G(z)=\frac1b\sum_{k=0}^{b-1}\omega^{-jk}G(\omega^kz)
=\sum_{m\equiv j\,(b)}a_m\frac{z^m}{m!}.}
\]

This holds as a formal series and wherever the original series
converges. Indeed the coefficient multiplier is the geometric sum
\(b^{-1}\sum_{k=0}^{b-1}\omega^{k(m-j)}\), equal to one or zero
according to the residue. It proves all the projector identities

\[
\sum_{j=0}^{b-1}\mathcal P_{b,j}=I,\qquad
\mathcal P_{b,j}\mathcal P_{b,l}=\delta_{jl}\mathcal P_{b,j}.
\]

The full-position section and the decimated lane are distinct arrays. A
section keeps the original positions, inserting zeros outside the
residue. A lane \(x_j[n]=a_{bn+j}\) changes the index to \(n\).

For an \textbf{ordinary} generating function \(X(z)=\sum_ma_m z^m\),
define \(X_j(w)=\sum_na_{bn+j}w^n\). Then

\[
X(z)=\sum_{j=0}^{b-1}z^jX_j(z^b).
\]

For an EGF the corresponding exact reconstruction is

\[
G(z)=\sum_{j=0}^{b-1}z^j
 \sum_{n\ge0}\frac{a_{bn+j}}{(bn+j)!}(z^b)^n.
\]

One must retain \((bn+j)!\). Replacing it by \(n!\) gives a different
function. The root-filter section identity has no such factorial
ambiguity.

\hypertarget{ordinary-taps-egf-multiplication-and-differentiation}{%
\subsection{Ordinary taps, EGF multiplication and
differentiation}\label{ordinary-taps-egf-multiplication-and-differentiation}}

Let \(h_0,\ldots,h_{L-1}\) be a finite tap vector and let \(a_m=0\) for
\(m<0\). Ordinary taps form

\[
b_m=\sum_{t=0}^{L-1}h_ta_{m-t}.
\]

For OGFs this is multiplication by \(H(z)=\sum_th_tz^t\):
\(B(z)=H(z)X(z)\). For EGFs it is a sum of repeated zero-constant
antiderivatives:

\[
\boxed{B_{\rm EGF}(z)=\sum_{t=0}^{L-1}h_t I^tG(z),\qquad
IG(z)=\int_0^zG(w)\,dw.}
\]

\textbf{Proof.} The EGF of the delayed sequence is
\(\sum_{m\ge t}a_{m-t}z^m/m!=I^tG(z)\). Summing taps gives the formula.
Conversely, \(G^{(t)}\) has coefficient sequence \(a_{m+t}\), so
derivatives sample future coefficients and do not implement these past
delays. \(\square\)

For two EGFs \(G_a,G_h\), ordinary function multiplication instead gives
the \textbf{Hurwitz product}

\[
G_a(z)G_h(z)=\sum_m\left[
\sum_{t=0}^m\binom mt h_ta_{m-t}\right]\frac{z^m}{m!}.
\]

This is a binomial convolution. For example, with \(a_m=1\) and taps
\(h=(1,1)\), ordinary taps give \(b_2=2\), while the Hurwitz product
gives \(b_2=3\). This explicit negative control catches a common
normalization error.

\textbf{Intrinsic nonlinear feedback taps} are the parameters
\((h_t,q_t,r_t)\) inside \(F'=1+\sum_th_tF(q_tz)F(r_tz)\). They change
the function's defining equation and its derivative recurrence. Ordinary
FIR taps are a linear transformation of an already defined coefficient
sequence. Both are useful, but they operate on different objects. The
intrinsic-bank growth theorem belongs to
\path{routes/asymmetric_dilation.md}.

\hypertarget{the-exact-finite-polyphase-tap-matrix}{%
\subsection{The exact finite polyphase tap
matrix}\label{the-exact-finite-polyphase-tap-matrix}}

For an OGF sequence \(a_m\), define lanes \(x_l[n]=a_{bn+l}\) and
\(X_l(w)=\sum_nx_l[n]w^n\). Apply finite ordinary taps and obtain
\(y_j[n]=b_{bn+j}\). Then the exact matrix identity is

\[
\boxed{Y_j(w)=\sum_{l=0}^{b-1}H_{jl}(w)X_l(w),\qquad
H_{jl}(w)=\sum_{\substack{0\le t<L\\t+l\equiv j\,(b)}}
h_t w^{(t+l-j)/b}.}
\]

Every exponent is a nonnegative integer. It records the carry across
phase boundaries, rather than erasing it.

\textbf{Proof.} A term \(h_t a_{bn+j-t}\) is in phase \(l\) exactly when
\(t+l\equiv j\pmod b\). Write \(d=(t+l-j)/b\). That term is
\(h_t x_l[n-d]\), whose OGF is \(h_tw^dX_l(w)\) by zero extension.
Summing over taps and phases gives the matrix. \(\square\)

For two phases and two taps,

\[
\begin{pmatrix}Y_0\\Y_1\end{pmatrix}
=\begin{pmatrix}h_0&h_1w\\h_1&h_0\end{pmatrix}
\begin{pmatrix}X_0\\X_1\end{pmatrix}.
\]

The upper-right factor \(w\) is necessary: a delay from odd phase to
even phase crosses one block. A diagonal-only phase operation cannot
implement this filter. The code constructs the sparse matrix for any
finite base and tap vector and checks that recombination gives the full
direct convolution.

\hypertarget{reuse-and-evidence}{%
\subsection{Reuse and evidence}\label{reuse-and-evidence}}

From the project directory, standard Python is enough:

\begin{verbatim}
python family_lab.py coefficients --q 1/4 --r 3/4 --n 8
python family_lab.py catalog --n 12 --output-dir examples
python checks/family_lab_checks.py
\end{verbatim}

The reusable functions include \path{asymmetric_odd_derivatives},
\path{feedback_odd_derivatives}, \path{asymmetric_polynomials},
\texttt{section}, \path{phase_lanes}, \path{merge_lanes},
\path{ordinary_taps}, \path{hurwitz_product},
\path{derivative_shift}, \path{integral_shift},
\path{polyphase_tap_matrix}, and \path{apply_polyphase_matrix}.
All core arithmetic is \texttt{fractions.Fraction} or integer
arithmetic. The phase and tap API applies to any finite rational
sequence, not just the special-function catalog.

The final completed run passed \textbf{8,038 exact assertions} and is
recorded in \path{evidence/family_lab_checks.json}. Its scope
includes:

\begin{itemize}
\tightlist
\item
  18 catalog rows through \(A_{12}\), independently compared with dense
  formal Picard iteration through degree 25.
\item
  Bivariate polynomials through \(A_8\), with 540 nonzero monomials
  checked for integer positivity, symmetry, parity and the degree bound.
\item
  Tangent coefficients independently obtained by division of sine and
  cosine, plus hyperbolic tangent, scaled sine and linear anchors.
\item
  Root-filter multipliers reduced exactly in cyclotomic fields for bases
  one through seven and indices zero through 31.
\item
  Direct FIR convolution compared with phase-matrix convolution for 26
  input lengths, seven bases and four tap vectors, including signed
  rational taps.
\item
  The antiderivative interpretation of EGF taps and the factorial
  conversion of Hurwitz products, independently compared with ordinary
  Taylor products.
\item
  Negative controls for omitted phase carries, aliasing lanes by
  concatenation, silent truncation, derivative/delay confusion and
  Hurwitz/FIR confusion.
\end{itemize}

The bounded checks validate the implemented paths in the recorded
ranges. The universal polynomial, degree, clearing, root-filter and
phase-matrix statements follow from the proofs above. The finite checks
do not establish entire continuation, growth order, or historical
novelty. Those claims are separately proved or explicitly left
unresolved in the project.

\path{evidence/family_lab_checks.json} records software version,
code hashes, exact arithmetic, bounds, category counts and run time. It
reports no failed assertion in this run. Independent audit identified
two API input-validation defects: a fractional residue was accepted by
the section operator, and an empty phase matrix accepted base zero. Both
are now rejected and have passing regression checks. These defects
affected invalid inputs, not the displayed identities. A draft prose
example initially transcribed \(A_3(1,2)\) and \(A_4(1,2)\) incorrectly;
the exact catalog comparison exposed the mismatch and corrected them to
\(35200\) and \(38041600\). The code result did not change. Larger
numerical evidence is unnecessary for these algebraic identities; an
independent mathematical audit should examine the reductions.
\endgroup

\clearpage\hypertarget{v11-coupled_growth_full-coupled-contraction-functions-and-their-perron-growth-law}{%
\section{Coupled contraction functions and their Perron growth
law}\label{v11-coupled_growth_full-coupled-contraction-functions-and-their-perron-growth-law}}

\hypertarget{v11-coupled_growth_full-definition-and-exact-coefficient-construction}{%
\subsection{Definition and exact coefficient
construction}\label{v11-coupled_growth_full-definition-and-exact-coefficient-construction}}

Fix real parameters

\[
0<q,r,a,b<1.
\]

Consider the normalized pair

\[
\boxed{\begin{aligned}
U'(z)&=1+U(qz)V(rz),\\
V'(z)&=1+U(az)V(bz),\\
U(0)&=V(0)=0.
\end{aligned}}
\]

\textbf{Theorem 1.} There is a unique pair of entire solutions. Both
functions are odd and have strictly positive coefficients at every odd
degree. Write

\[
U(z)=\sum_{n\ge0}u_nz^{2n+1},\qquad
V(z)=\sum_{n\ge0}v_nz^{2n+1}.
\]

Then \(u_0=v_0=1\), and, for \(n\ge1\),

\[
\begin{aligned}
u_n&=\frac1{2n+1}\sum_{i=0}^{n-1}
 q^{2i+1}r^{2(n-1-i)+1}u_i v_{n-1-i},\\
v_n&=\frac1{2n+1}\sum_{i=0}^{n-1}
 a^{2i+1}b^{2(n-1-i)+1}u_i v_{n-1-i}.
\end{aligned}
\]

\textbf{Proof.} Comparing successive coefficients in the differential
equations determines the coefficient of degree \(k+1\) from coefficients
of degrees at most \(k\). Thus the formal pair is unique. The displayed
recurrence constructs an odd formal pair satisfying both equations.
Formal uniqueness excludes any additional even coefficients. Strict
positivity follows by induction.

The elementary bounds

\[
0<u_n\le2^{-n},\qquad 0<v_n\le2^{-n}
\]

hold by simultaneous induction. If they hold below \(n\), each product
in either recurrence is at most \(2^{-(n-1)}\), and there are \(n\)
products. The resulting bound is \(n2^{-(n-1)}/(2n+1)<2^{-n}\). Hence
the two series converge on \(|z|<\sqrt2\), and define a holomorphic
solution there.

Put \(Q=\max(q,r,a,b)<1\). If the pair is already defined on \(|z|<R\),
then

\[
\int_0^z[1+U(qw)V(rw)]\,dw,\qquad
\int_0^z[1+U(aw)V(bw)]\,dw
\]

define holomorphic extensions on \(|z|<R/Q\). The integrands are
holomorphic on that disk. On the old disk the new functions agree with
the old pair because their derivatives and values at zero agree. On the
larger disk, every contracted argument still lies in the old disk, so
the extensions satisfy the same equations. Iterating exhausts the
complex plane. Any other entire solution agrees with the unique germ and
therefore everywhere. \(\square\)

In particular,

\[
\begin{aligned}
U(z)&=z+\frac{qr}{3}z^3
 +\frac{qr(qr\,q^2+ab\,r^2)}{15}z^5+O(z^7),\\
V(z)&=z+\frac{ab}{3}z^3
 +\frac{ab(qr\,a^2+ab\,b^2)}{15}z^5+O(z^7).
\end{aligned}
\]

The factorial coefficients \(A_n=(2n+1)!u_n\) and \(B_n=(2n+1)!v_n\)
satisfy

\[
\begin{aligned}
A_n&=\sum_{i=0}^{n-1}\binom{2n}{2i+1}
 q^{2i+1}r^{2(n-1-i)+1}A_iB_{n-1-i},\\
B_n&=\sum_{i=0}^{n-1}\binom{2n}{2i+1}
 a^{2i+1}b^{2(n-1-i)+1}A_iB_{n-1-i}.
\end{aligned}
\]

Thus \(A_n,B_n\) are polynomials with nonnegative integer coefficients
in the four scales. These formulas are exact algebraic constructions.
Finite checks of them do not replace the analytic proof.

\hypertarget{v11-coupled_growth_full-the-characteristic-exponent-and-its-positive-vector}{%
\subsection{The characteristic exponent and its positive
vector}\label{v11-coupled_growth_full-the-characteristic-exponent-and-its-positive-vector}}

For \(s>0\), let

\[
P(s)=\begin{pmatrix}q^s&r^s\\a^s&b^s\end{pmatrix}.
\]

There is a unique \(\rho>0\) with

\[
\boxed{(1-q^\rho)(1-b^\rho)=(ar)^\rho.}
\]

Indeed, the left side strictly increases continuously from zero to one,
while the right side strictly decreases continuously from one to zero.
This equation is precisely \(\det(I-P(\rho))=0\). Define

\[
\widetilde w=(r^\rho,1-q^\rho)^T,
\qquad
w=\frac{\widetilde w}{\max(r^\rho,1-q^\rho)}.
\]

Both coordinates of \(w\) are positive, \(\max(w_1,w_2)=1\), and direct
multiplication gives

\[
P(\rho)w=w.
\]

The other eigenvalue is \(q^\rho+b^\rho-1\), whose absolute value is
strictly less than one. Consequently

\[
\boxed{\operatorname{spr}P(\rho)=1.}
\]

Conversely, a positive matrix of this form having spectral radius one
has eigenvalue one and satisfies the determinant equation. This proves
the proposed spectral characterization without requiring a numerical
eigensolver. A positive left eigenvector is

\[
\ell=(a^\rho,1-q^\rho)^T,\qquad \ell^TP(\rho)=\ell^T.
\]

The notation \(w\) for the growth vector is deliberately different from
\(v_n\), the coefficients of \(V\).

\hypertarget{v11-coupled_growth_full-exact-order-and-explicit-positive-type-bounds}{%
\subsection{Exact order and explicit positive type
bounds}\label{v11-coupled_growth_full-exact-order-and-explicit-positive-type-bounds}}

For \(R>0\), positivity gives

\[
M_U(R)=\max_{|z|=R}|U(z)|=U(R),\qquad
M_V(R)=V(R).
\]

Set \(F_1=U\), \(F_2=V\), and \(H_i(R)=\log F_i(R)\). Since
\(F_i(R)\ge R\), these logarithms are positive for \(R>1\). The type
convention used here is

\[
\underline\tau_i=\liminf_{R\to\infty}\frac{\log M_{F_i}(R)}{R^\rho},
\qquad
\overline\tau_i=\limsup_{R\to\infty}\frac{\log M_{F_i}(R)}{R^\rho}.
\]

\textbf{Theorem 2.} Both functions have order \(\rho\), positive lower
type and finite upper type. More precisely, the explicit constants below
satisfy

\[
0<c_-w_i\le\underline\tau_i\le\overline\tau_i\le c_+w_i<\infty,
\qquad i=1,2.
\]

Moreover,

\[
\lim_{R\to\infty}\frac{\log\log M_{F_i}(R)}{\log R}=\rho.
\]

This radial exponent limit is stronger than the limsup used to define
order. It does not assert that either type has a limit or that growth is
uniform in the angular variable.

\hypertarget{v11-coupled_growth_full-explicit-lower-comparison}{%
\subsubsection{Explicit lower
comparison}\label{v11-coupled_growth_full-explicit-lower-comparison}}

Let

\[
m=\min(q,r,a,b),\qquad Q=\max(q,r,a,b),
\]

\[
C=\frac{Q}{1-Q},\qquad B=1+\frac2m,
\qquad c_-=\frac{\log2}{(B+C)^\rho}.
\]

We prove the stronger statement

\[
\boxed{H_i(R)\ge c_-w_i(R+C)^\rho\quad(R\ge2,\ i=1,2).}
\]

On \(2\le R\le B\), its right side is at most \(\log2\), because
\(w_i\le1\), whereas \(H_i(R)\ge\log R\ge\log2\).

For the first equation, integrate only over the final unit interval to
obtain, when \(R\ge B\),

\[
U(R)\ge\int_{R-1}^R U(qt)V(rt)\,dt
\ge U(q(R-1))V(r(R-1)).
\]

The second equation gives the same inequality with \((q,r)\) replaced by
\((a,b)\) and \(U(R)\) replaced by \(V(R)\). All four child radii are at
least two. Every scale \(p\in\{q,r,a,b\}\) satisfies

\[
p(R-1)+C\ge p(R+C),
\]

because \((1-p)C\ge p\). If the comparison already holds at the child
radii, then

\[
\begin{aligned}
H_1(R)&\ge c_-\{w_1[q(R-1)+C]^\rho
                     +w_2[r(R-1)+C]^\rho\}\\
&\ge c_-(q^\rho w_1+r^\rho w_2)(R+C)^\rho
=c_-w_1(R+C)^\rho.
\end{aligned}
\]

The second row of \(P(\rho)w=w\) proves the second component in exactly
the same way.

For a formal propagation, put

\[
d=(1-Q)B+Q>0.
\]

Every child of \(R\ge B\) is at most \(Q(R-1)\le R-d\). Simultaneous
induction over the intervals \([B+kd,B+(k+1)d]\), starting with the
established seed interval, therefore proves both comparisons for every
\(R\ge2\). Continuity handles shared endpoints. The argument uses only
monotonicity of positive powers and works unchanged when \(0<\rho<1\).

\hypertarget{v11-coupled_growth_full-explicit-upper-comparison}{%
\subsubsection{Explicit upper
comparison}\label{v11-coupled_growth_full-explicit-upper-comparison}}

Set

\[
\Delta=\log(1/m),\qquad
D=\max\left(\log\frac1{qr},\log\frac1{ab}\right)+\log2.
\]

Choose any integer \(N\ge0\) satisfying \(Q^{-N}\ge3/m\), and define

\[
W_N=(2^{N+1}-1)\log2
 +(2^{N+1}-N-2)\log(1/Q),
\]

\[
S=W_N+1+\Delta+D,\qquad
c_+=\frac{e^{-\rho}S}{\min(w_1,w_2)}.
\]

These constants require only the parameters and the uniquely defined
root \(\rho\). They require no evaluation of an unknown solution.

To bound a seed interval, the coefficient majorant from Theorem 1 gives
\(F_i(1)\le2\). Put \(R_k=Q^{-k}\), \(Z_0=2\), and

\[
Z_k=Q^{-k}(1+Z_{k-1}^2).
\]

Monotonicity of the positive functions and the integral equations give
\(F_i(R_k)\le Z_k\) simultaneously. As \(Z_k\ge1\),

\[
\log Z_k\le k\log(1/Q)+\log2+2\log Z_{k-1}.
\]

Summing this elementary recurrence gives \(\log Z_N\le W_N\). Therefore
\(H_i(e^t)\le W_N\) throughout \(1\le t\le1+\Delta\), since
\(e^t\le e/m<3/m\le R_N\).

Now let

\[
G_i(t)=H_i(e^t)+t+D.
\]

For \(t\ge1+\Delta\), all child radii are at least \(e\), so their
function values are at least one. The full integral gives

\[
H_1(e^t)\le t+\log2
 +H_1(e^{t-\log(1/q)})+H_2(e^{t-\log(1/r)}).
\]

The definition of \(D\) converts this into

\[
G_1(t)\le G_1(t-\log(1/q))+G_2(t-\log(1/r)).
\]

Similarly,

\[
G_2(t)\le G_1(t-\log(1/a))+G_2(t-\log(1/b)).
\]

On the seed interval, \(G_i(t)\le S\le c_+w_i e^{\rho t}\). The Perron
equation propagates this comparison: substituting the bounds at the
children in the first inequality gives
\(c_+e^{\rho t}(q^\rho w_1+r^\rho w_2)=c_+w_1e^{\rho t}\), and the
second row works likewise. All logarithmic delays belong to
\([\log(1/Q),\Delta]\). Induction in intervals of length
\(\log(1/Q)>0\), after the seed interval, therefore proves

\[
\boxed{H_i(R)\le c_+w_iR^\rho-\log R-D\quad(R\ge e).}
\]

The lower and upper comparisons imply the asserted positive, finite
types. Taking logarithms of these bounds for sufficiently large \(R\)
gives the claimed radial exponent limits and hence exact orders.
\(\square\)

\hypertarget{v11-coupled_growth_full-the-two-type-vectors-have-exactly-the-perron-proportions}{%
\subsection{The two type vectors have exactly the Perron
proportions}\label{v11-coupled_growth_full-the-two-type-vectors-have-exactly-the-perron-proportions}}

The explicit constants above are conservative. There is nevertheless an
exact statement about the relative types of the two components.

\textbf{Theorem 3.} There exist positive finite constants
\(\underline\tau\le\overline\tau\) such that

\[
\boxed{\begin{aligned}
(\underline\tau_1,\underline\tau_2)^T&=\underline\tau\,w,\\
(\overline\tau_1,\overline\tau_2)^T&=\overline\tau\,w.
\end{aligned}}
\]

In particular,

\[
\frac{\underline\tau_1}{\underline\tau_2}
=\frac{\overline\tau_1}{\overline\tau_2}
=\frac{r^\rho}{1-q^\rho}
=\frac{1-b^\rho}{a^\rho}.
\]

Either both ordinary type limits exist or neither does. When they exist,
their vector is a positive multiple of \(w\).

\textbf{Proof.} Use the final-unit integral inequality from the lower
proof and divide by \(R^\rho\). Taking the liminf gives

\[
\underline t\ge P(\rho)\underline t,
\qquad
\underline t=(\underline\tau_1,\underline\tau_2)^T.
\]

Here \(p(R-1)\) tends to infinity continuously and
\([p(R-1)]^\rho/R^\rho\to p^\rho\), which justifies each scaled liminf.
The liminf of a sum is at least the sum of the liminfs. All quantities
are positive and finite by Theorem 2.

The full-integral inequality

\[
H_1(R)\le\log R+\log2+H_1(qR)+H_2(rR)
\]

and its second-row counterpart imply, on taking the limsup,

\[
\overline t\le P(\rho)\overline t.
\]

The terms \((\log R+\log2)/R^\rho\) vanish. Dotting either componentwise
inequality with the strictly positive left eigenvector \(\ell\) from
Section~\ref{v11-coupled_growth_full-the-characteristic-exponent-and-its-positive-vector} shows that its nonnegative difference vector has zero positive
weighted sum. Every coordinate of that difference is therefore zero.
Thus both type vectors are fixed by \(P(\rho)\). Its eigenvalue-one
space is one dimensional, so each is a positive multiple of \(w\).

If one component has an ordinary type limit, its lower and upper types
are equal. Because that coordinate of \(w\) is positive, this forces
\(\underline\tau=\overline\tau\), and the other component also has a
type limit. \(\square\)

This theorem does not assert a limit for \(\log U(R)/\log V(R)\) when
type oscillations persist. It also does not identify the scalar values
of either type or prove their equality.

\hypertarget{v11-coupled_growth_full-reductions-a-worked-unequal-pair-and-parameter-structure}{%
\subsection{Reductions, a worked unequal pair, and parameter
structure}\label{v11-coupled_growth_full-reductions-a-worked-unequal-pair-and-parameter-structure}}

\textbf{Exact scalar reduction.} The components coincide if and only if
the unordered pairs of scales coincide:

\[
U=V\quad\Longleftrightarrow\quad\{q,r\}=\{a,b\}.
\]

For sufficiency, the scalar solution \(T_{q,r}\) from
\path{routes/asymmetric_dilation.md} makes \((T_{q,r},T_{q,r})\)
solve the coupled equations. Theorem 1 gives uniqueness. For necessity,
equality at degree three gives \(qr=ab=p>0\). Equality at degree five
then gives \(q^2+r^2=a^2+b^2\). The two positive unordered pairs have
the same product and the same sum \(\sqrt{q^2+r^2+2p}\), so they
coincide. This characterizes equality of the normalized components. A
stronger assertion about all possible changes of variables requires its
own nonreduction theorem.

On either scalar reduction locus, the characteristic equation is
equivalent to \(q^\rho+r^\rho=1\), recovering the scalar order theorem.
At \(q=r=a=b=1/2\),

\[
U(z)=V(z)=\sqrt2\sinh(z/\sqrt2),\qquad \rho=1.
\]

\textbf{A genuinely unequal order-one pair.} Set

\[
(q,r,a,b)=\left(\frac15,\frac25,\frac35,\frac7{10}\right).
\]

The determinant identity at \(\rho=1\) is
\((1-1/5)(1-7/10)=(3/5)(2/5)=6/25\). Thus the exact order is one, and

\[
w=(1/2,1)^T.
\]

The first terms are

\[
\begin{aligned}
U(z)&=z+\frac2{75}z^3+\frac{88}{234375}z^5
 +\frac{18814}{8544921875}z^7+O(z^9),\\
V(z)&=z+\frac7{50}z^3+\frac{8211}{1250000}z^5
 +\frac{85664601}{625000000000}z^7+O(z^9).
\end{aligned}
\]

Both lower and upper exponential types of \(V\) are exactly twice the
corresponding types of \(U\). The row sums of \(P(1)\) are \(3/5\) and
\(13/10\). Consequently neither an unweighted average row condition nor
the maximum of the two scalar row orders yields this example's coupled
order. The positive eigenvector accounts for the exchange of growth
between the components.

For this example the elementary lower constants are \(m=1/5\),
\(Q=7/10\), \(C=7/3\), \(B=11\), and \(c_-=3\log2/40\). Hence, for every
\(R\ge2\),

\[
\log U(R)\ge\frac{3\log2}{80}\left(R+\frac73\right),\qquad
\log V(R)\ge\frac{3\log2}{40}\left(R+\frac73\right).
\]

\textbf{Dependence of the order on the scales.} The determinant equation
shows that the order depends on the self scales \(q,b\) and the product
\(ar\) of the two cross scales. Changing \(a\) and \(r\) while keeping
their product fixed leaves the order fixed, provided both remain in
\((0,1)\). The Perron vector and the normalized functions generally
change.

For any prescribed \(s>0\) and any \(q,r\in(0,1)\), choices satisfying

\[
a^s r^s=(1-q^s)(1-b^s),\qquad 0<a,b<1,
\]

lie on the order-\(s\) surface. For instance, one may select \(q,r,b\)
with \((1-q^s)(1-b^s)<r^s\), then define

\[
a=\frac{[(1-q^s)(1-b^s)]^{1/s}}{r}.
\]

This constructs exact order surfaces without fitting growth from
truncated series.

\hypertarget{v11-coupled_growth_full-proof-provenance-and-remaining-questions}{%
\subsection{Proof provenance and remaining
questions}\label{v11-coupled_growth_full-proof-provenance-and-remaining-questions}}

The decisive extension is from a scalar positive comparison to a
comparison weighted by a positive eigenvector. The fixed-width integral
controls the lower type at the exact exponent. A fixed proportional
interval \([\alpha R,R]\), with \(\alpha<1\), would only give a smaller
exponent through \(\alpha^sP(s)\). Letting \(\alpha\) tend to one
recovers exact order but does not, by itself, yield positive lower type.
This is a limitation of that alternative route, not a failure of the
theorem.

The complete proofs above close the target on the strictly positive
contraction domain. No computed example is used as a universal analytic
premise. The displayed rational Taylor coefficients were evaluated
directly from the exact recurrence with Python
\texttt{fractions.Fraction} as a transcription check. Their independent
formal residual checks belong to the release's exact computation record.

Remaining questions include the existence and value of ordinary type
limits, possible oscillations, zeros, angular growth, connection
formulas, and regimes with zero, negative, complex, or noncontracting
scales. Positivity and strict contraction are hypotheses of the results
here. This route makes no historical priority claim and no claim of a
Lean formalization.

Self-review coverage: this entire file. Routine checks: consistent
coefficient and eigenvector notation, domains for every logarithm and
power, seed and child interval endpoints, and both vector orientations.
No unresolved mathematical gap was identified in the author's
self-review. Independent review remains a separate record.

\clearpage\hypertarget{v11-coupled_inverse_full-when-the-pair-reduces-and-when-its-parameters-are-intrinsic}{%
\section{When the pair reduces, and when its parameters are
intrinsic}\label{v11-coupled_inverse_full-when-the-pair-reduces-and-when-its-parameters-are-intrinsic}}

Throughout, \(q,r,a,b\in(0,1)\). Let the normalized coupled solution be

\[
U'=1+U(qz)V(rz),\qquad V'=1+U(az)V(bz),\qquad U(0)=V(0)=0.
\]

Write \(U(z)=z+u_1z^3+u_2z^5+u_3z^7+\cdots\) and
\(V(z)=z+v_1z^3+v_2z^5+v_3z^7+\cdots\). Put \(p=qr\), \(k=ab\).
Coefficient matching gives

\[
u_1=\frac p3,\qquad v_1=\frac k3,
\] \[
u_2=\frac p{15}(p q^2+k r^2),\qquad
v_2=\frac k{15}(p a^2+k b^2),
\] \[
u_3=\frac p7(q^4u_2+p^2u_1v_1+r^4v_2),\qquad
v_3=\frac k7(a^4u_2+k^2u_1v_1+b^4v_2).
\]

Define the fifth-order separation quantity

\[
\boxed{D=k^2u_2-p^2v_2
=\frac{pk ar}{15}(q^2-b^2)(qb-ar).}
\]

This factorization will distinguish equality of normalized functions up
to amplitude and argument rescaling from a change of coordinates in
their defining functional differential operator. These are different
equivalence relations.

\hypertarget{v11-coupled_inverse_full-exact-classification-of-component-rescaling}{%
\subsection{Exact classification of component
rescaling}\label{v11-coupled_inverse_full-exact-classification-of-component-rescaling}}

\textbf{Theorem.} The following statements are equivalent:

\begin{enumerate}
\def\labelenumi{\arabic{enumi}.}
\tightlist
\item
  There are nonzero complex constants \(A,B\) such that \(V(z)=A U(Bz)\)
  for every complex \(z\).
\item
  \(D=0\).
\item
  \(q=b\) or \(qb=ar\).
\end{enumerate}

If these conditions hold, put \(s=\sqrt{k/p}>0\), and let \(T_{x,y}\)
denote the normalized scalar solution \(T'=1+T(xz)T(yz)\), \(T(0)=0\).
Then

\[
\boxed{
U(z)=\sqrt{s}\,T_{q,sr}(z/\sqrt{s}),\qquad
V(z)=s^{-1/2}T_{q,sr}(\sqrt{s}\,z).
}
\]

The effective scales belong to \((0,1)\): if \(q=b\), then
\(sr=\sqrt{ar}\), and if \(qb=ar\), then \(s=b/r\) and \(sr=b\).

\textbf{Proof.} Normalization gives \(AB=1\). Comparing cubic
coefficients gives \(B^2=k/p\), and comparing fifth coefficients gives
\(v_2=(k/p)^2u_2\). Thus component rescaling implies \(D=0\), even when
complex rescalings are allowed. The displayed factorization, strict
positivity, and \(q,b>0\) show that \(D=0\) is equivalent to \(q=b\) or
\(qb=ar\).

For the converse define the displayed pair using \(T=T_{q,sr}\). Direct
differentiation gives

\[
U'(z)=1+T(qz/\sqrt{s})T(r\sqrt{s}\,z)
=1+U(qz)V(rz).
\]

The second equation reduces to equality of the unordered pairs of
arguments

\[
\{q\sqrt{s},\;s^{3/2}r\}
=\{a/\sqrt{s},\;b\sqrt{s}\}.
\]

When \(qb=ar\), one has \(a=sq\) and \(b=sr\), giving equality in the
displayed order. When \(q=b\), one has \(a=s^2r\), giving equality after
interchanging the factors. The pair also has the required initial values
and derivatives. Uniqueness of the formal coefficient recursion
identifies it with the coupled solution. Finally, \(V(z)=s^{-1}U(sz)\),
proving rescaling.

The positive parameter domain therefore contains two explicit reduction
surfaces. Outside them the two component functions cannot be obtained
from one another by amplitude and argument rescaling. This statement
does not exclude every representation using previously known functions.

\hypertarget{v11-coupled_inverse_full-seven-derivatives-recover-all-four-ordered-scales-away-from-reduction}{%
\subsection{Seven derivatives recover all four ordered scales away
from
reduction}\label{v11-coupled_inverse_full-seven-derivatives-recover-all-four-ordered-scales-away-from-reduction}}

\textbf{Theorem.} If \(D\ne0\), the six coefficients
\((u_1,v_1,u_2,v_2,u_3,v_3)\) uniquely determine the ordered tuple
\((q,r,a,b)\) within \((0,1)^4\). Consequently, equality of the two
labeled normalized germs implies equality of these parameters in this
regime.

\textbf{Proof.} First recover \(p=3u_1\), \(k=3v_1\). If \(x=q^2\), then
\(r^2=p^2/x\), and the equation for \(u_2\) gives

\[
p x^2-\frac{15u_2}{p}x+kp^2=0.
\]

This leaves at most two positive candidates. If there are two distinct
positive roots \(x,x'\), their product is \(xx'=kp\). The corresponding
seventh-coefficient values differ by

\[
u_3(x)-u_3(x')
=\frac p7(x^2-x'^2)
\left(u_2-\frac{p^2}{k^2}v_2\right)
=\frac{pD}{7k^2}(x^2-x'^2).
\]

The difference is nonzero because \(D\ne0\) and the distinct positive
roots have distinct squares. Thus \(u_3\) selects the unique root
\(q^2\), giving \(q=\sqrt{x}>0\) and \(r=p/q\). A repeated quadratic
root already gives a unique value and requires no selection.

Likewise \(y=a^2\) obeys

\[
p y^2-\frac{15v_2}{k}y+k^3=0,
\]

with root product \(yy'=k^3/p\). The two candidate seventh coefficients
have difference

\[
v_3(y)-v_3(y')
=\frac k7(y^2-y'^2)
\left(u_2-\frac{p^2}{k^2}v_2\right)
=\frac{D}{7k}(y^2-y'^2),
\]

which is again nonzero for distinct positive candidates. Recover
\(a=\sqrt{y}\) and \(b=k/a\). All selections must also respect the
contraction domain. Since the supplied coefficients arise from an
admissible pair, at least one admissible selection exists. The preceding
differences give uniqueness.

This is an exact-data inverse theorem. The displayed differences become
small as \(D\to0\), so the argument does not supply a uniform stable
inverse near the reduction surfaces. It also does not identify unlabeled
components without the additional component-permutation equivalence.

\hypertarget{v11-coupled_inverse_full-a-worked-pair-with-exact-order-one}{%
\subsection{A worked pair with exact order
one}\label{v11-coupled_inverse_full-a-worked-pair-with-exact-order-one}}

Take

\[
(q,r,a,b)=\left(\frac15,\frac25,\frac35,\frac7{10}\right),\qquad
p=\frac2{25},\quad k=\frac{21}{50}.
\]

At \(\rho=1\), the characteristic matrix has positive eigenvector
\((1/2,1)^T\) with eigenvalue one. Its other eigenvalue is \(-1/10\).
The growth theorem therefore gives exact order one for both components.
Here \(q\ne b\) and \(qb\ne ar\), so \(D\ne0\). The pair is outside both
reduction surfaces and its four labeled parameters are recoverable from
the coefficients through degree seven.

\hypertarget{v11-coupled_inverse_full-verification-scope}{%
\subsection{Verification scope}\label{v11-coupled_inverse_full-verification-scope}}

The identities above are algebraic and the proofs use exact positive
parameters. The companion check script compares the first coefficients
against the full independent recursion, verifies the factorization over
a rational grid, checks the candidate-selection formulas, and tests both
reduction surfaces by comparing many exact Taylor coefficients. Finite
checks test the implementation and the worked examples. The equivalence
and recovery theorems follow from the proofs.

Historical priority of these specific statements remains a separate
literature question.

\clearpage\hypertarget{v11-coupled_structure_full-constant-linear-coordinates-and-coupled-dilation}{%
\section{Constant linear coordinates and coupled
dilation}\label{v11-coupled_structure_full-constant-linear-coordinates-and-coupled-dilation}}

\hypertarget{v11-coupled_structure_full-exact-meaning-of-decoupling}{%
\subsection{Exact meaning of
decoupling}\label{v11-coupled_structure_full-exact-meaning-of-decoupling}}

Let \(q,r,a,b>0\), with contraction \(<1\) imposed when discussing the
normalized entire solution. For an arbitrary analytic vector input
\(X=(u,v)\) with \(u(0)=v(0)=0\), define the functional vector field

\[
\mathcal F[X](z)=
\begin{pmatrix}1+u(qz)v(rz)\\1+u(az)v(bz)\end{pmatrix}.
\tag{S1}
\]

A nonzero constant row \(c=(c_1,c_2)\in\mathbb C^2\) is an
\textbf{autonomous scalar projection} if some scalar functional
\(\mathcal H\) satisfies

\[
c\mathcal F[X]=\mathcal H[cX]
\quad\text{for every such analytic input }X.
\tag{S2}
\]

The functional \(\mathcal H\) may use the whole scalar input. Thus this
definition allows any scalar dilation equation, not only a prescribed
quadratic ansatz. Equality (S2) means that two vector inputs with the
same scalar projection have the same projected vector field. It is an
identity of functional differential operators. It does not assert an
identity only along the distinguished normalized solution.

A constant invertible change \(W=PX\) completely decouples the operator
if both rows of \(P\) are autonomous scalar projections. No assumption
about the form of their scalar functionals is needed.

\hypertarget{v11-coupled_structure_full-classification-of-autonomous-linear-projections}{%
\subsection{Classification of autonomous linear
projections}\label{v11-coupled_structure_full-classification-of-autonomous-linear-projections}}

\textbf{Theorem S1.} For positive nonzero scales, a nonzero row \(c\)
satisfies (S2) if and only if

\[
q=a,\qquad r=b,\qquad c=\gamma(1,-1),\qquad\gamma\ne0.
\tag{S3}
\]

In that exceptional case the projected vector field is identically zero.
Consequently, no invertible constant linear change completely decouples
(S1), including when some or all scales coincide. If \((q,r)\ne(a,b)\),
there is not even one nonzero autonomous scalar projection. Four
pairwise distinct contraction scales are a special case of this stronger
result.

\textbf{Proof.} Subtract the constant forcing. Define

\[
Q_c[u,v](z)=c_1u(qz)v(rz)+c_2u(az)v(bz).
\]

If (S2) holds, \(Q_c\) is constant on every fiber of
\((u,v)\mapsto c_1u+c_2v\). Fix \(z_0\ne0\). Polynomial interpolation
gives a scalar polynomial \(f\) with \(f(0)=0\) and \(f(sz_0)=1\) at
every distinct scale \(s\in\{q,r,a,b\}\). For inputs \(u=xf,v=yf\),
where \(x,y\in\mathbb C\),

\[
Q_c[xf,yf](z_0)=(c_1+c_2)xy.
\tag{S4}
\]

It must depend only on \(c_1x+c_2y\). If \(c_1+c_2\ne0\),
differentiating the polynomial \(xy\) along the fiber direction
\((c_2,-c_1)\) would give

\[
c_2y-c_1x=0\quad\text{for every }x,y,
\]

which forces \(c=0\), a contradiction. Therefore \(c_2=-c_1\), and
scaling the row reduces the question to \(c=(1,-1)\).

For this row, put \(u=w+g,v=g\). This input and \((w,0)\) have the same
projection \(w\). The second input has \(Q_c[w,0]=0\), so autonomy
forces

\[
w(qz)g(rz)-w(az)g(bz)
+g(qz)g(rz)-g(az)g(bz)=0
\tag{S5}
\]

for all zero-constant analytic \(w,g\). Setting \(w=0\) and subtracting
that instance of (S5) leaves

\[
w(qz)g(rz)=w(az)g(bz).
\tag{S6}
\]

Taking successively \((w,g)=(z,z),(z^2,z),(z,z^2)\) yields \(qr=ab\),
\(q^2r=a^2b\), and \(qr^2=ab^2\). Nonzero scales then give \(q=a\) and
\(r=b\). Conversely these equalities make the two components of (S1)
identical on every input, so \((1,-1)\mathcal F[X]=0\). The only
possible autonomous rows therefore span a one-dimensional space, which
cannot be the two rows of an invertible matrix. \(\square\)

The same proof works on odd analytic inputs for positive scales: choose
\(f(z)=z p(z^2)\) by interpolation at the distinct positive scale
squares, and use \((z,z),(z^3,z),(z,z^3)\) in (S6). Positivity then
resolves the equalities of squares. The theorem does not depend on
allowing even inputs.

\hypertarget{v11-coupled_structure_full-tensor-explanation-for-four-distinct-scales}{%
\subsection{Tensor explanation for four distinct
scales}\label{v11-coupled_structure_full-tensor-explanation-for-four-distinct-scales}}

The preceding proof also makes a coordinate invariant obstruction
explicit. Under an invertible state change, write the original
coordinates as \(u=\ell_1(W),v=\ell_2(W)\), where \(\ell_1,\ell_2\) are
independent linear forms. At the pair of distinct nodes \((qz,rz)\), the
first quadratic term has bilinear coefficient

\[
B=\ell_1\otimes\ell_2.
\]

It has rank one, but is not symmetric:
\(B-B^{\mathsf T}=\ell_1\wedge\ell_2\ne0\). At four distinct nonzero
nodes, polynomial interpolation lets their vector values vary
independently, so terms supported at \((az,bz)\) cannot cancel this
coefficient. A scalar autonomous channel at the same node pair would
require \(B=\kappa\ell\otimes\ell\), a symmetric rank-one tensor, after
projecting the output onto that channel. For any output row with a
nonzero coefficient on this term, the alternating part remains nonzero.
The other node pair gives the same obstruction if the first output
coefficient vanishes.

This is not a rank-only argument: the tensor \(\ell_1\otimes\ell_2\)
already has rank one. Its nonzero alternating part is the decisive
invariant at distinct nodes. When the nodes of a product coincide, the
symmetric quadratic form \(uv\) has matrix

\[
\begin{pmatrix}0&1/2\\1/2&0\end{pmatrix},
\]

of rank two, whereas a square of one linear coordinate has rank one.
Rank is then the appropriate obstruction. The fiber proof in Theorem S1
handles node coincidences and cancellations without mixing these two
cases.

\hypertarget{v11-coupled_structure_full-particular-solutions-and-degenerate-parameters}{%
\subsection{Particular solutions and degenerate
parameters}\label{v11-coupled_structure_full-particular-solutions-and-degenerate-parameters}}

Theorem S1 is an operator statement. It does not exclude coincident
components for particular initial data, algebraic relations along a
solution, elimination into a scalar equation of higher order, nonlinear
coordinate changes, or identities with other named special functions.

For contraction scales, the normalized analytic germ is unique. If
\((a,b)=(q,r)\), its components coincide because their derivatives are
identical and their initial values agree. The common component is the
scalar family \(T'=1+T(qz)T(rz),T(0)=0\). The same coincidence holds if
\((a,b)=(r,q)\): \((T,T)\) solves both equations because multiplication
is commutative, and uniqueness identifies the normalized solution. For
\(q\ne r\), this reversed-pair case has no autonomous scalar operator
projection by Theorem S1, despite having \(U=V\) along its normalized
solution. This is a concrete reason the two notions must remain
distinct.

In fact, equality of normalized components occurs exactly when the
unordered scale pairs agree:

\[
U=V\quad\Longleftrightarrow\quad\{q,r\}=\{a,b\}.
\tag{S7}
\]

To prove necessity, write \(U=V=T=z+c_1z^3+\cdots\). Then
\(3c_1=qr=ab>0\). Comparing the \(z^4\) coefficients of
\(T(qz)T(rz)=T(az)T(bz)\) gives \(c_1qr(q^2+r^2)=c_1ab(a^2+b^2)\), so
both products and sums of squares agree. Positive scales then have equal
sums and are the roots of the same quadratic. Sufficiency is the
uniqueness argument above. For example, \((q,r,a,b)=(1/2,1/3,1/3,1/2)\)
has coincident normalized components and no nonzero autonomous operator
projection.

At \(q=r=a=b=s>0\), the variables \(w=u+v,d=u-v\) give

\[
w'=2+\tfrac12\bigl(w(sz)^2-d(sz)^2\bigr),\qquad d'=0.
\]

This triangular system is not a completely decoupled operator. The
normalized solution lies in the invariant subspace \(d=0\).

Zero scales lie outside Theorem S1. On zero-constant inputs, a product
vanishes identically if either of its scales is zero. If both products
vanish, both equations are already \(u'=v'=1\). If only the first
product vanishes, the projection \(u\) is autonomous, regardless of
whether the second pair agrees with the first. Thus extending Theorem S1
to zero scales would be false. Scale one and positive scales greater
than one do not affect the algebraic operator theorem, although the
contraction-domain entire growth theorems do not automatically extend to
them. Negative or complex scales and odd-input coincidences caused by
signs are not classified here.

\hypertarget{v11-coupled_structure_full-claim-status}{%
\subsection{Claim status}\label{v11-coupled_structure_full-claim-status}}

Theorem S1 and (S7) are proved above. The proof is elementary and
independent of the coupled growth theorem. Historical priority for these
precise operator statements is not asserted. The theorem provides a
specified nondecoupling result, not a universal assertion that the
normalized functions cannot reduce to any previously studied function
class.

\clearpage\hypertarget{v11-resolution_full-a-visual-atlas-of-finite-resolution}{%
\section{A visual atlas of finite
resolution}\label{v11-resolution_full-a-visual-atlas-of-finite-resolution}}

A structural change need not produce a large change in the measurements
we can make. A tiny feedback tap can change the order of an entire
function. An excluded prime can change the exact arithmetic inverse. A
negative atom can destroy positivity of a measure. Each change becomes
visible at a different resolution.

This atlas makes those statements quantitative. Its three controls are
observation radius, absolute precision and polynomial degree. Each
example includes a bound that can be evaluated without trusting the
appearance of a curve. The large radii below are sufficient guarantees
from conservative estimates. The degree bounds bracket a detection
problem. Neither is presented as an optimal threshold.

\begin{longtable}[]{@{}
  >{\raggedright\arraybackslash}p{(\columnwidth - 6\tabcolsep) * \real{0.2500}}
  >{\raggedright\arraybackslash}p{(\columnwidth - 6\tabcolsep) * \real{0.2500}}
  >{\raggedright\arraybackslash}p{(\columnwidth - 6\tabcolsep) * \real{0.2500}}
  >{\raggedright\arraybackslash}p{(\columnwidth - 6\tabcolsep) * \real{0.2500}}@{}}
\toprule\noalign{}
\begin{minipage}[b]{\linewidth}\raggedright
Family
\end{minipage} & \begin{minipage}[b]{\linewidth}\raggedright
Structural difference
\end{minipage} & \begin{minipage}[b]{\linewidth}\raggedright
Finite observation
\end{minipage} & \begin{minipage}[b]{\linewidth}\raggedright
Certificate that resolves it
\end{minipage} \\
\midrule\noalign{}
\endhead
\bottomrule\noalign{}
\endlastfoot
Positive feedback & A nonzero fast tap changes entire order & Values on
a fixed complex disk & A sufficient radius for a factor of two
separation \\
Arithmetic depth & The excluded prime set changes & Values or finitely
many derivatives on a disk & A bracket for the negative depth needed to
reach an absolute tolerance \\
Signed sieve & One prime shelf has negative mass & All squared
polynomial norms up to a given degree & An explicit rational polynomial
with negative norm \\
\end{longtable}

\hypertarget{v11-resolution_full-a-weak-tap-that-changes-global-growth}{%
\subsection{A weak tap that changes global
growth}\label{v11-resolution_full-a-weak-tap-that-changes-global-growth}}

Let the normalized positive feedback function satisfy

\[
F_\varepsilon'(z)=1+(1-\varepsilon)F_\varepsilon(z/4)^2
 +\varepsilon F_\varepsilon(3z/4)^2,
\qquad F_\varepsilon(0)=0,
\qquad 0\le\varepsilon\le1.
\]

The base theorem gives a unique entire function with nonnegative odd
coefficients. At zero weight, the order is \(1/2\). At every positive
weight it is

\[
\rho=\frac{\log2}{\log(4/3)}\approx2.40942084.
\]

Thus the order jumps at a parameter value where the entire functions
themselves converge uniformly on every compact set.

\hypertarget{v11-resolution_full-the-disk-has-a-numerical-error-budget}{%
\subsubsection{The disk has a numerical error
budget}\label{v11-resolution_full-the-disk-has-a-numerical-error-budget}}

Write \(f=f_{3/4}\) for the normalized diagonal function, and
\(f(z)=\sum_{n\ge0}d_nz^{2n+1}\). The inherited tap continuity theorem
gives

\[
\sup_{|z|\le R}|F_\varepsilon(z)-F_0(z)|
\le \varepsilon L(R),\qquad
L(R)=\sum_{n\ge1}n d_nR^{2n+1}
=\frac{Rf'(R)-f(R)}2.
\tag{RA-1}
\]

The following elementary certificate makes the unit disk example fully
computable. Positive coefficient domination by tangent, followed by
\(\sin x\le x\) and \(\cos x\ge1-x^2/2\), gives

\[
f(4/3)\le\tan(4/3)\le12.
\]

For \(N\ge1\) and \(0\le R\le1\), put \(t=3R/4\). If
\((N+2)t^2\le N+1\), then \(n t^{2n+1}\) is nonincreasing for
\(n\ge N+1\), and hence

\[
0\le L(R)-\sum_{n=1}^{N}n d_nR^{2n+1}
\le12(N+1)t^{2N+3}.
\tag{RA-2}
\]

Indeed the left side is a sum of \(d_n(4/3)^{2n+1}\) multiplied by
\(n t^{2n+1}\). Bound the multiplier by its first value, then use the
positive total mass \(f(4/3)\le12\). The coefficients \(d_n\) are
rational. With \(N=40\), the exact rational upper bound is less than

\[
L(1)<0.241890267.
\]

The decimal shown here is rounded upward from the rational certificate.
Consequently

\[
\boxed{\sup_{|z|\le1}|F_{10^{-12}}(z)-F_0(z)|<2.41890267\times10^{-13},}
\]

although their entire orders are different. If observations have an
absolute error allowance at least half this bound, midpoint data can be
consistent with both functions throughout the disk. This is an ambiguity
under a stated noise allowance, not equality of their exact values.

\emph{The corresponding figure appears in the resolution visual
chapter.}

\textbf{How to read the figure.} The left curves divide the difference
by the tap weight, so the relevant small quantity is the plotted height
multiplied by \(\varepsilon\). The right panel records the proved order,
which is determined by arbitrarily large radii. The curve computation
uses powers through \(z^{201}\), with a normalized analytic omitted-tail
bound below \(5.26\times10^{-23}\) on the unit interval. Ordinary
floating rounding is separate from that analytic tail bound.

\hypertarget{v11-resolution_full-a-radius-at-which-a-difference-is-guaranteed}{%
\subsubsection{A radius at which a difference is
guaranteed}\label{v11-resolution_full-a-radius-at-which-a-difference-is-guaranteed}}

The base lower-type theorem supplies, for every \(\varepsilon>0\),

\[
\log F_\varepsilon(R)\ge c_\varepsilon(R+3/\varepsilon)^\rho
\ge c_\varepsilon R^\rho\qquad(R\ge2),
\] \[
c_\varepsilon=\frac{\log2}{(4/\varepsilon+8/3)^\rho}.
\]

For the slow function, the base upper-type construction with
\(q=r=1/4\), \(N=2\), yields

\[
\log F_0(R)\le C_0\sqrt R\qquad(R\ge e),
\qquad C_0=e^{-1/2}(1+22\log2)\approx9.85566103.
\]

The value of \(C_0\) comes directly from the earlier explicit constants:
\(V_2=15\log2\), \(\Delta=2\log2\), and \(B_+=5\log2\). Thus
\(V_2+1+\Delta+B_+=1+22\log2\).

Define

\[
R_2(\varepsilon)=\max\left\{
 e,\left(\frac{2C_0}{c_\varepsilon}\right)^{1/(\rho-1/2)},
 \left(\frac{2\log2}{c_\varepsilon}\right)^{1/\rho}\right\}.
\tag{RA-3}
\]

Then, for every \(R\ge R_2(\varepsilon)\),

\[
\boxed{F_\varepsilon(R)\ge2F_0(R).}
\]

To prove this, the last two constraints imply that
\(c_\varepsilon R^\rho\) is at least both \(2C_0\sqrt R\) and
\(2\log2\). It is therefore at least their half-sum
\(C_0\sqrt R+\log2\). Exponentiating proves the claim.

\begin{longtable}[]{@{}
  >{\raggedright\arraybackslash}p{(\columnwidth - 4\tabcolsep) * \real{0.2727}}
  >{\raggedleft\arraybackslash}p{(\columnwidth - 4\tabcolsep) * \real{0.3636}}
  >{\raggedleft\arraybackslash}p{(\columnwidth - 4\tabcolsep) * \real{0.3636}}@{}}
\toprule\noalign{}
\begin{minipage}[b]{\linewidth}\raggedright
Tap weight
\end{minipage} & \begin{minipage}[b]{\linewidth}\raggedleft
Difference on the unit disk is at most
\end{minipage} & \begin{minipage}[b]{\linewidth}\raggedleft
Sufficient radius \(R_2\), rounded
\end{minipage} \\
\midrule\noalign{}
\endhead
\bottomrule\noalign{}
\endlastfoot
\(10^{-3}\) & \(2.419\times10^{-4}\) & \(2.029\times10^5\) \\
\(10^{-6}\) & \(2.419\times10^{-7}\) & \(1.237\times10^9\) \\
\(10^{-9}\) & \(2.419\times10^{-10}\) & \(7.549\times10^{12}\) \\
\(10^{-12}\) & \(2.419\times10^{-13}\) & \(4.608\times10^{16}\) \\
\end{longtable}

The radius column displays the analytic formula approximately. For a
certificate use the formula itself or a rigorously rounded upper value.
These large numbers reflect the conservatism of the global comparison.
They do not mean that observations at all smaller radii fail. The
formula grows proportionally to \(\varepsilon^{-\rho/(\rho-1/2)}\) as
the tap vanishes.

\emph{The corresponding figure appears in the resolution visual
chapter.}

\hypertarget{v11-resolution_full-an-excluded-prime-and-the-precision-needed-to-see-it}{%
\subsection{An excluded prime and the precision needed to see
it}\label{v11-resolution_full-an-excluded-prime-and-the-precision-needed-to-see-it}}

Fix \(a=\beta=2\), and use the arithmetic depth family

\[
A_Q(s)=\sum_{\substack{n\ge2\\\mu(n)^2=1\\(n,Q)=1}}
 n^{-4}(1-n^{-2})^{2s}.
\]

For a prime \(p\), let \(D_p=A_1-A_p\). This positive difference
consists exactly of the squarefree atoms divisible by \(p\). Its largest
frequency is

\[
\lambda_p=-2\log(1-p^{-2}),
\]

and its total mass is \(D(4)/(p^4+1)\), where
\(D(4)=\zeta(4)/\zeta(8)=105/\pi^4\).

\hypertarget{v11-resolution_full-exact-disk-bounds-include-derivative-measurements}{%
\subsubsection{Exact disk bounds include derivative
measurements}\label{v11-resolution_full-exact-disk-bounds-include-derivative-measurements}}

For every integer \(k\ge0\) and \(R\ge0\),

\[
\sup_{|s|\le R}|D_p^{(k)}(s)|
\le \lambda_p^k e^{R\lambda_p}\frac{D(4)}{p^4+1}.
\tag{RA-4}
\]

This is the inherited arithmetic inverse theorem, now read as a
precision requirement. For fixed \(R\), the envelope is of order
\(p^{-4-2k}\). Each additional derivative suppresses the large-prime
difference even further in absolute units. This does not rule out
methods that rescale derivatives and account for their corresponding
noise amplification.

At integer radius \(R=5\), an especially simple rational certificate
follows from \(D(4)<\zeta(4)<4/3\):

\[
\boxed{\sup_{|s|\le5}|D_p(s)|
 <\frac{4}{3(p^4+1)}(1-p^{-2})^{-10}.}
\tag{RA-5}
\]

For derivatives, multiply this rational bound by \([2/(p^2-1)]^k\),
using \(-\log(1-u)\le u/(1-u)\). Thus this part of the certificate needs
no transcendental evaluation at all.

\begin{longtable}[]{@{}lr@{}}
\toprule\noalign{}
Excluded prime & Rational upper bound on \(|s|\le5\), rounded upward \\
\midrule\noalign{}
\endhead
\bottomrule\noalign{}
\endlastfoot
101 & \(1.283\times10^{-8}\) \\
1009 & \(1.287\times10^{-12}\) \\
10007 & \(1.330\times10^{-16}\) \\
\end{longtable}

\emph{The corresponding figure appears in the resolution visual
chapter.}

\hypertarget{v11-resolution_full-increasing-negative-depth-gives-a-two-sided-threshold-bracket}{%
\subsubsection{Increasing negative depth gives a two-sided threshold
bracket}\label{v11-resolution_full-increasing-negative-depth-gives-a-two-sided-threshold-bracket}}

There is also an elementary lower bound, because the removed prime atom
is present:

\[
\boxed{p^{-4}e^{\lambda_pT}\le D_p(-T)
\le\frac{D(4)}{p^4+1}e^{\lambda_pT}\qquad(T\ge0).}
\tag{RA-6}
\]

Choose an absolute tolerance \(\eta>D(4)/(p^4+1)\). The strict
inequality makes the initial difference smaller than the tolerance.
Define

\[
T_-=\frac{\log\bigl(\eta(p^4+1)/D(4)\bigr)}{\lambda_p},
\qquad T_+=\frac{\log(\eta p^4)}{\lambda_p}.
\tag{RA-7}
\]

Then \(D_p(-T)\le\eta\) for \(0\le T\le T_-\), while \(D_p(-T)\ge\eta\)
for \(T\ge T_+\). The true crossing exists and is unique because
\(D_p(-T)\) is continuous and strictly increasing. Therefore it belongs
to \([T_-,T_+]\). If the tolerance hypothesis fails, the interval
starting at zero is not certified by this formula. The executable record
represents an absent lower interval by \texttt{null}.

The bracket width and its large-prime scale are explicit:

\[
T_+-T_-=\frac{\log\bigl(D(4)p^4/(p^4+1)\bigr)}{\lambda_p},
\qquad
T_+=\frac{p^2}{2}\bigl(4\log p+\log\eta\bigr)(1+O(p^{-2})).
\]

For \(\eta=10^{-6}\), rounded values are:

\begin{longtable}[]{@{}
  >{\raggedright\arraybackslash}p{(\columnwidth - 4\tabcolsep) * \real{0.2727}}
  >{\raggedleft\arraybackslash}p{(\columnwidth - 4\tabcolsep) * \real{0.3636}}
  >{\raggedleft\arraybackslash}p{(\columnwidth - 4\tabcolsep) * \real{0.3636}}@{}}
\toprule\noalign{}
\begin{minipage}[b]{\linewidth}\raggedright
Prime
\end{minipage} & \begin{minipage}[b]{\linewidth}\raggedleft
No larger than tolerance through \(T_-\)
\end{minipage} & \begin{minipage}[b]{\linewidth}\raggedleft
At least tolerance by \(T_+\)
\end{minipage} \\
\midrule\noalign{}
\endhead
\bottomrule\noalign{}
\endlastfoot
101 & 23,307.789 & 23,690.516 \\
1009 & 7,012,695.633 & 7,050,894.424 \\
10007 & 1,149,287,766.447 & 1,153,045,061.461 \\
\end{longtable}

These are approximate evaluations of exact analytic endpoints, not
directed-rounding intervals.

\hypertarget{v11-resolution_full-the-direct-difference-and-subtraction-of-full-functions-have-different-conditioning}{%
\subsubsection{The direct difference and subtraction of full functions
have different
conditioning}\label{v11-resolution_full-the-direct-difference-and-subtraction-of-full-functions-have-different-conditioning}}

The positive series for \(D_p\) can be evaluated directly if the
excluded prime is already specified. Recovering it by subtracting
independently evaluated full functions can be much harder. The atom
\(n=2\) in \(A_1\) gives

\[
\boxed{\frac{A_1(-T)}{D_p(-T)}
\ge\frac{p^4+1}{16D(4)}e^{(\lambda_2-\lambda_p)T}.}
\tag{RA-8}
\]

At the sufficient endpoint \(T_+\), the base-ten logarithm of this lower
bound is about 5,924 for \(p=101\), 1,761,864 for \(p=1009\), and
288,119,991 for \(p=10007\). These are conditioning scales for
subtracting full signals. They are not running times, required precision
for every possible algorithm, or claimed numerical evaluations of the
full functions at those enormous depths. The logarithmic bound is
computed directly without forming those functions.

\emph{The corresponding figure appears in the resolution visual
chapter.}

\hypertarget{v11-resolution_full-a-negative-atom-that-every-low-degree-test-misses}{%
\subsection{A negative atom that every low-degree test
misses}\label{v11-resolution_full-a-negative-atom-that-every-low-degree-test-misses}}

Consider the one-prime signed sieve at \(p=2\), with
\(\lambda_2=-1-\delta\), \(\delta>0\). Its nodes are
\(x_n=(1-n^{-2})^2\), and its quadratic form is

\[
Q_\delta(f)=\sum_{\substack{n\ge1\\\mu(n)^2=1\\n\text{ odd}}}n^{-4}f(x_n)^2
-\delta\sum_{\substack{n\ge1\\\mu(n)^2=1\\n\text{ even}}}n^{-4}f(x_n)^2.
\]

Every even squarefree atom has negative mass. We will certify that all
polynomials through degree four nevertheless have positive norm for
\(\delta=10^{-12}\), and then construct a degree-seven polynomial of
negative norm. This is an assertion about every polynomial in a
finite-dimensional space, not a plot of a few test vectors.

\hypertarget{v11-resolution_full-a-rational-certificate-for-every-polynomial-through-a-chosen-degree}{%
\subsubsection{A rational certificate for every polynomial through a
chosen
degree}\label{v11-resolution_full-a-rational-certificate-for-every-polynomial-through-a-chosen-degree}}

Choose \(d+1\) distinct odd squarefree labels \(n_0,\ldots,n_d\),
including 1, and set \(y_i=x_{n_i}\). Their Lagrange polynomials are

\[
\ell_i(x)=\prod_{j\ne i}\frac{x-y_j}{y_i-y_j},
\quad
b_i=\prod_{j\ne i}\frac{\max(y_j,1-y_j)}{|y_i-y_j|},
\quad K_d=\sum_{i=0}^{d}n_i^4b_i^2.
\]

For \(x\in[0,1]\), we have \(|\ell_i(x)|\le b_i\). Interpolation and
weighted Cauchy--Schwarz give, for every real polynomial of degree at
most \(d\),

\[
|f(x)|^2\le
\left(\sum_{i=0}^{d}n_i^{-4}f(y_i)^2\right)
\left(\sum_{i=0}^{d}n_i^4\ell_i(x)^2\right)
\le K_d\sum_{i=0}^{d}n_i^{-4}f(y_i)^2.
\]

The total even mass satisfies
\(\sum_{n\text{ even}}\mu(n)^2n^{-4}<\frac1{16}(1+\frac13)=1/12\).
Keeping the selected nodes in the positive part yields the explicit
bound

\[
\boxed{Q_\delta(f)\ge
\left(1-\frac{\delta K_d}{12}\right)
\sum_{i=0}^{d}n_i^{-4}f(y_i)^2.}
\tag{RA-9}
\]

A nonzero polynomial of degree at most \(d\) cannot vanish at all
\(d+1\) nodes. Thus \(\delta<12/K_d\) proves strict positive
definiteness of the Hankel matrix of order \(d+1\).

For \(d=4\), use labels \(1,3,5,7,11\). Exact rational arithmetic gives

\[
K_4\approx1.50393441996\times10^{11},\qquad
12/K_4\approx7.9790713217\times10^{-11}>10^{-12}.
\]

The inequality is checked as an exact rational comparison. Therefore
every nonzero polynomial of degree at most four has positive norm in the
chosen signed measure.

\hypertarget{v11-resolution_full-a-shorter-detector-by-killing-the-first-positive-shelves}{%
\subsubsection{A shorter detector by killing the first positive
shelves}\label{v11-resolution_full-a-shorter-detector-by-killing-the-first-positive-shelves}}

Let \(S\) be exactly the set of all odd squarefree labels \(n\le K\),
where \(K\) is an odd squarefree integer, and let \(n_+\) be the next
such label. Assume \(n_+>2\), and define

\[
P_{S,m}(x)=\prod_{n\in S}\frac{x-x_n}{x_2-x_n}
\left(\frac{1-x}{1-x_2}\right)^m,
\] \[
C_S=\prod_{n\in S}\frac{1-x_n}{|x_2-x_n|},
\qquad r_S=\frac{1-x_{n_+}}{1-x_2}<1.
\]

This polynomial has degree \(|S|+m\), equals one at \(x_2\), and kills
every positive atom through the cutoff. At each remaining positive node,
the product of the absolute root factors is at most \(C_S\), and the
last factor is at most \(r_S^m\). Dropping all other negative atoms and
bounding the positive tail by the integral
\(\sum_{n>K}n^{-4}\le1/(3K^3)\) gives

\[
\boxed{Q_\delta(P_{S,m})\le-\frac\delta{16}
 +\frac{C_S^2}{3K^3}r_S^{2m}.}
\tag{RA-10}
\]

For any chosen \(S\), a computable sufficient exponent is the first
nonnegative integer \(m\) for which

\[
 r_S^{2m}<\frac{3K^3\delta}{16C_S^2}.
\]

It can be found by multiplying rational numbers, without logarithms or
eigenvalues. Searching finitely many choices of \(S\) improves the
sufficient degree but does not prove its optimality.

For the worked defect \(\delta=10^{-12}\), take

\[
S=\{1,3,5,7\},\quad K=7,\quad n_+=11,\quad m=3.
\]

Then

\[
C_S=\frac{756482048}{20777538825},\qquad
r_S=\frac{3856}{102487},
\]

and \(P_{S,3}\) has degree seven. The exact rational right side of
(RA-10) is strictly negative and is approximately

\[
\boxed{Q_{10^{-12}}(P_{S,3})\le-5.8845708899\times10^{-14}.}
\]

The negative atom contributes \(-6.25\times10^{-14}\). The full positive
tail contributes at most \(3.6542911008\times10^{-15}\), a bound that
covers every omitted atom. The exact certificate is stored as numerator
and denominator. The original generic isolating detector needed degree
21 for this defect. Using the known sign pattern and killing later
positive shelves reduces the exhibited degree to seven.

\emph{The corresponding figure appears in the resolution visual
chapter.}

The fifth-order Hankel matrix is therefore positive definite, while the
eighth-order matrix is not positive semidefinite. The two intervening
polynomial degrees, five and six, are not decided by these certificates.
A numerical eigenvalue could suggest a sharper boundary, but it would be
a different calculation requiring its own conditioning and error
analysis.

\hypertarget{v11-resolution_full-what-the-three-examples-establish-together}{%
\subsection{What the three examples establish
together}\label{v11-resolution_full-what-the-three-examples-establish-together}}

The three resolution controls act on different mathematical objects.
Radius controls what part of an entire function is sampled. Precision
controls whether two different arithmetic parameters are distinguishable
within an observation model. Degree controls the richness of a
positivity test. The common lesson has a precise form: a property
determined by an unbounded collection of tests need not have a uniform
margin on each bounded subcollection.

Here the positive tap changes a proved global order. The excluded prime
changes an exactly recoverable radical. The signed sieve lies outside
the exact positivity domain even while a whole low-degree polynomial
space passes. Each conclusion follows from explicit formulas and local
proofs. There is no claim that the three mechanisms are identical, and
no historical novelty claim for interpolation, spectral peeling, or the
use of isolating polynomials.

\hypertarget{v11-resolution_full-reproduction-and-evidence}{%
\subsection{Reproduction and
evidence}\label{v11-resolution_full-reproduction-and-evidence}}

Run from the release root:

\begin{verbatim}
python checks/resolution_v11.py
python build_resolution_figures.py
\end{verbatim}

The first command uses only the Python standard library. It records
\textbf{985 exact rational assertions} and \textbf{16 floating algebra
checks}, with source hashes and exact certificate numerators and
denominators. It tests interpolation identities in the monomial basis,
all selected root values, sample envelope inequalities, exact negative
norms for finite sums, coefficient continuity bounds and each displayed
certificate. The universal analytic conclusions rely on the proofs
above, rather than on the finite count.

The figure command uses NumPy, Matplotlib and mpmath. Six figures are
saved in PDF and PNG form. The weak-tap curves use 100 decimal digits
and an explicit positive-coefficient tail bound. Every other curve
displays a closed analytic bound or rational certificate, so it has no
omitted arithmetic series. Displayed floating values and polynomial
curves are not interval arithmetic. The script never evaluates the huge
full arithmetic functions in the cancellation example.

One preliminary presentation issue was corrected during development: a
negative formula for \(T_-\) was initially clamped to zero. That would
misleadingly label an already visible difference as invisible at zero.
The finished formula requires its stated tolerance hypothesis, and the
executable record marks an absent interval by \texttt{null}. The exact
assertion run passed. This correction did not affect any of the three
listed prime examples.

\clearpage\hypertarget{v11-coupled_lab_full-exact-laboratory-and-certified-evaluation}{%
\section{Exact laboratory and certified
evaluation}\label{v11-coupled_lab_full-exact-laboratory-and-certified-evaluation}}

The reusable module \path{coupled_lab.py} uses Python's standard
library and imports the existing exact coefficient, phase and tap
operations from \path{family_lab.py}. The figure builder additionally
uses NumPy and Matplotlib. Neither script contacts an external service.
The generic integer and signed catalog rows are formal power series,
with no claim of analytic convergence.

\hypertarget{v11-coupled_lab_full-normalization-and-recurrence}{%
\subsection{Normalization and
recurrence}\label{v11-coupled_lab_full-normalization-and-recurrence}}

Write

\[
U(z)=\sum_{n\ge0}A_n\frac{z^{2n+1}}{(2n+1)!},\qquad
V(z)=\sum_{n\ge0}B_n\frac{z^{2n+1}}{(2n+1)!}.
\]

The exact recurrence is

\[
A_n=\sum_{i+j=n-1}\binom{2n}{2i+1}q^{2i+1}r^{2j+1}A_iB_j,
\qquad
B_n=\sum_{i+j=n-1}\binom{2n}{2i+1}a^{2i+1}b^{2j+1}A_iB_j,
\]

with \(A_0=B_0=1\). Coefficients with even derivative index vanish.
Swapping the components sends the parameter tuple \((q,r,a,b)\) to
\((b,a,r,q)\). Identical rows \((a,b)=(q,r)\), or exchanged rows
\((a,b)=(r,q)\), yield the scalar reduction \(U=V=T_{q,r}\).

Both derivative arrays are polynomials in four scale variables with
nonnegative integer coefficients. Their total degree is at most
\(n(n+1)\). Indeed, a recurrence term with \(i+j=n-1\) has degree at
most \(2n+i(i+1)+j(j+1)\le n(n+1)\). If all four denominators divide
\(D\), then \(D^{n(n+1)}A_n\) and \(D^{n(n+1)}B_n\) are integers. The
clearing factor is a convenient uniform choice and need not be minimal.

\hypertarget{v11-coupled_lab_full-an-elementary-exact-evaluation-certificate}{%
\subsection{An elementary exact evaluation
certificate}\label{v11-coupled_lab_full-an-elementary-exact-evaluation-certificate}}

This argument uses \(0<q,r,a,b<1\) and does not require the order
theorem. Let \(Q=\max(q,r,a,b)\). All ordinary odd Taylor coefficients
are positive and at most one: induction in the ordinary recurrence
bounds the coefficient of \(z^{2n+1}\) by \(n/(2n+1)\) for \(n\ge1\).
Thus both series converge on the unit disk and, for \(0\le x<1\),

\[
\max(U(x),V(x))\le \frac{x}{1-x^2}.
\]

Analytic continuation by the integral equations expands each known disk
by the factor \(1/Q\), so both solutions are entire. Positive
coefficients give \(\max_{|z|\le x}|F(z)|=F(x)\) for either component.
The integral equations therefore imply

\[
M(x):=\max(U(x),V(x))\le x\bigl(1+M(Qx)^2\bigr).
\]

Start at the first member \(Q^kx<1\) of this finite contraction ladder,
use \(x/(1-x^2)\) there, and repeatedly apply \(B(x)=x(1+B(Qx)^2)\)
backward. This defines the rational \path{disk_majorant} bound for
rational parameters and radius. The implementation has explicit
recursion and integer size guards, which raise errors rather than
returning a weakened result.

If \(F_N\) retains odd indices \(0,\ldots,N\), then for \(S>R\ge0\),

\[
\max_{|z|\le R}|F(z)-F_N(z)|
\le B(S)\left(\frac RS\right)^{2N+3}.
\]

For every omitted exponent \(m\ge2N+3\), use \(R^m\le (R/S)^{2N+3}S^m\),
sum positive terms, and bound their sum by \(F(S)\le B(S)\). No extra
geometric denominator is needed. For rational \(x\ge0\),
\path{certified_evaluation} returns the exact interval
\([F_N(x),F_N(x)+B(S)(x/S)^{2N+3}]\) for each component.

\hypertarget{v11-coupled_lab_full-worked-calls}{%
\subsection{Worked calls}\label{v11-coupled_lab_full-worked-calls}}

\begin{verbatim}
from fractions import Fraction as Q
import coupled_lab as lab

p = (Q(1,5), Q(2,5), Q(3,5), Q(7,10))
A, B = lab.coupled_odd_derivatives(*p, last=12)
certificate = lab.certified_evaluation(*p, Q(1), last=30, outer_radius=2)
root = lab.pf_order(*p)  # floating approximation, not interval arithmetic

derivatives = lab.full_egf_coefficients(B)
phases = lab.phase_lanes(derivatives, 3)
matrix = lab.polyphase_tap_matrix([1, -2, 1], 3)
length = len(derivatives) + 2
output_phases = lab.apply_polyphase_matrix(phases, matrix, length)
output = lab.merge_lanes(output_phases, length)
\end{verbatim}

At \(x=1\), with \(N=30\) and \(S=2\), the common tail bound is below
\(1.600\times10^{-13}\). The partial values are approximately
\(U(1)=1.027044341827337\) and \(V(1)=1.1467074957737624\). These
displayed decimals are rounded illustrations. The returned fractions are
the actual certificate.

Run the independent checks and regenerate figures from the project root:

\begin{verbatim}
python3 checks/coupled_v11.py
python3 build_coupled_figures.py
\end{verbatim}

The first command passes 1,559 exact assertions, 10 numerical
comparisons, and 13 separate input/resource guards. Independent
constructions include dense ordinary Taylor Picard iteration,
sine/cosine division for tangent, the exact sinh formula, component
symmetry, and direct convolution versus phase carry. The test counts
describe bounded implementation checks, not universal proofs.


\MonographSubpart{Evidence and limitations}
\pageheading{EVIDENCE / COUPLED FAMILY}{What the coupled extension establishes}
\renewcommand{\arraystretch}{1.28}
\begin{tabular}{@{}p{.28\linewidth}p{.65\linewidth}@{}}
\toprule
Result & Scope\\\midrule
Coupled entire construction & Unique normalized pair for four positive strict contractions, with positive odd coefficients\\
Common growth law & Exact order from the characteristic matrix, positive lower type, finite upper type, and exact Perron proportions for both type vectors\\
Component reduction & Rescaling occurs exactly on $q=b$ or $qb=ar$\\
Generic exact inverse & Coefficients through degree seven identify all four ordered scales when $D\ne0$\\
Operator obstruction & No invertible constant linear change decouples the defining operator on all analytic inputs\\
Resolution certificates & Explicit compact, radius, prime-depth and polynomial-degree bounds for the three worked examples\\
\bottomrule
\end{tabular}

\claim{Verification and reproducibility}
The coupled-system finite suite passes \textbf{5,272 exact assertions and 26 numerical
assertions}. Thirteen input and resource guards are recorded separately.
The exact inverse is checked on 256 rational parameter tuples. The coupled
API is compared with an independent dense Picard construction and classical
anchors. The resolution checks verify the exact rational certificates.
These counts are separate from the earlier four-family evidence of 19,871 exact
and 5,576 numerical assertions.

Independent mathematical review reconstructs the proof implications and the
preserved source files. It corrected an initial-segment quantifier in the
detector theorem. The worked certificate already used that initial segment.
The source package records the correction, the repair of a figure label, and the successful checks.

\claim{Questions that remain open}
The coupled theorem does not determine an ordinary type limit, uniform angular
growth, or the complex zero geometry. Stable finite-precision recovery near the
reduction surfaces remains open. The resolution bounds are constructive and
can be conservative. The sieve example leaves degrees five and six undecided.
The literature comparison does not settle historical priority.

The exact API and its checks run with the Python standard library. NumPy, Matplotlib, and
mpmath generate the figures. The full manuscript supplies the proof
appendices and the source archive contains the code and inputs needed for the computations.

\pageheading{EVIDENCE / FOUNDATIONS}{Established results and open questions}
The foundational families close selected structural questions and admit a reproducible visual treatment. Historical priority is a distinct question that the bounded literature review does not settle.

\begin{tabular}{@{}>{\raggedright\arraybackslash}p{0.46\linewidth}>{\raggedright\arraybackslash}p{0.47\linewidth}@{}}
\toprule
Established results & Remaining questions\\\midrule
Positive lower and finite upper dilation type & Existence of an optimal type limit and logarithmic modulation\\
Joint analytic parameter dependence and compact tap bounds & General minimal-bank identifiability\\
Exact arithmetic parameter recovery up to modulus radical & Stable recovery within bounded prime classes\\
All-prime signed-atom detectors and positivity domain & Sharp sieve-root asymptotics and multiplicities\\
Global orbit-zeta continuation and shifted determinants & Functional equations and derivative-weighted zero sums\\
Integer coefficient and exact polyphase/tap machinery & Historical priority of the candidate extensions\\
\bottomrule\end{tabular}

\claim{The strongest direct function candidate}
The asymmetric dilation extension remains the strongest direct discovery
candidate: its parameters are intrinsic, its off-diagonal members cannot be
rescaled from the known diagonal, and its nonnegative real scale regime has now been
mapped. The arithmetic inverse and signed-resolution results add structural
depth. The orbit continuation is a worked application of established spectral
methods, with its exact seed and branch conventions verified locally.

\claim{A common finite-resolution lesson}
A weak tap can change asymptotic order while preserving compact convergence.
A large excluded prime can change an exact arithmetic parameter while hiding
within a small bounded-window error. A negative atom can evade low-degree
positivity and become detectable at a calculable larger degree.
These are three different proved mechanisms. Their shared resolution theme
does not make them identical constructions.

\pageheading{EVIDENCE / EXACT CHECKS}{Evidence with declared ranges}
The independent audit reconstructed five proof chapters and examined their mathematical statements. The release runner executed ten scripts once, with
25 pinned inputs and byte-identical isolated snapshots. All inputs remained
unchanged, every declared result was present, and the schema validator passed.

\[
\boxed{19{,}871\text{ exact assertion statements}\quad+
5{,}576\text{ numerical assertion statements passed}.}
\]
The exact total comprises 19,331 principal checks and 540 separately recorded
helper checks in the inverse reconstruction. Operational path and process
guards are excluded. Counts are assertion statements, not counts of theorems
or every scalar operation inside a grouped assertion.

The API contributes 8,038 exact assertions and 18 catalog rows through index
12. Dilation type and continuity add 6,091 exact checks. Arithmetic inverse
examples include rank 10,000, negative depth 500, and large excluded primes.
The multiple-prime run tests 102 positive parameter/depth cases, 1,395
orthogonality conditions, 381 recurrence cases, and 40 conditional pencils.
Spectral continuation tests classical unshifted and shifted anchors.

\claim{Recovered failures are retained}
The preserved audit documents early environment and plotting issues, an unsuitable absolute determinant residual, and the previously known one-dilation case.
Two API validation defects were reproduced and fixed. A draft
sequence example had two miscopied integers and was corrected against the
exact catalog. Initial figure footer and label overlaps were repaired.
A direct publisher-source retrieval failed; an authenticated publisher
preview supplied the classical Laplace-uniqueness comparison instead.

Exact analytic tail envelopes are distinguished from floating-point rounding.
All decimal arithmetic and root approximations are ordinary high precision,
not interval certificates. No analytic Lean certificate is claimed.
The universal conclusions rely on the proofs, not on the number of checks.

\pageheading{EVIDENCE / REPRODUCIBILITY}{A reproducible research package}
The preserved source archive includes its manuscript, all 18 vector figures and PNG
previews, figure generators, exact API, JSON/CSV catalog, inherited baseline
routes, the five new proof chapters, independent audit, and bounded release
evidence with input and output hashes.

The source files preserve the earlier claim ledger and open questions as provenance. Later proof chapters supersede the earlier statements where a new argument supplies a closed result. The distinction is recorded in the verification documents rather than folded into the mathematical theorem count.

\claim{Reproduce the work}
Use the versions listed in \texttt{requirements.txt}; execute
\texttt{checks/run\_release.py} in a fresh output directory. The checked
instructions give the bounded-run invocation and validator. Figure scripts
then rebuild the atlas from formulas and recorded numerical examples.
The manuscript is typeset with LaTeX. The build script resolves the contents and cross references through successive passes.

The standard-library API is portable. Numerical checks additionally require
mpmath. No network access, external submission, or dependency installation is
part of the verification run. PDF byte identity can change with timestamps
and font environments; mathematical outputs retain their own provenance.

The remaining questions admit concrete starting points: type-limit analysis
for dilation, precision-controlled arithmetic inversion within bounded
parameter classes, sharper sieve-root asymptotics, and spectral weighted sums.
They remain research questions rather than hidden dependencies of the
completed release claims.

\pageheading{NOTATION / GLOSSARY}{Symbols and structural vocabulary}
The notation is introduced where each family is defined. The entries here collect the conventions used across the five parts.

\begin{description}
\setlength{\itemsep}{7pt}
\item[Contracted feedback $T_{q,r}$.] The normalized odd entire function determined by $T'_{q,r}(z)=1+T_{q,r}(qz)T_{q,r}(rz)$ and $T_{q,r}(0)=0$, for $0<q,r<1$.
\item[Feedback bank $F_{\Theta}$.] A weighted finite family of contraction channels whose evaluations enter the defining nonlinear differential equation.
\item[Coupled pair $(U,V)$.] Two normalized entire functions satisfying $U'=1+U(qz)V(rz)$ and $V'=1+U(az)V(bz)$, with strict positive contractions $q,r,a,b$.
\item[Entire order $\rho(f)$.] For an entire nonpolynomial function, the limsup of $\log\log M_f(R)/\log R$ as $R\to\infty$, with $M_f(R)=\max_{|z|=R}|f(z)|$.
\item[Arithmetic depth $A_{a,\beta,Q}$.] A complex parameter family formed from squarefree coefficients, positive logarithmic frequencies, and exclusion of integers sharing prime factors with the modulus $Q$.
\item[M\"obius function $\mu$.] The arithmetic function for which $\mu(n)^2$ selects squarefree integers.
\item[Prime-sieve measure $\nu_{F,\boldsymbol\lambda}$.] An atomic arithmetic measure whose masses are modified at primes in a finite set $F$. Signed parameter choices may invalidate positive definiteness.
\item[Hankel matrix.] A moment matrix $H_d=(M_{i+j})_{0\le i,j\le d}$. The positivity of every such matrix differs from positivity at any fixed degree.
\item[Orbit spectral zeta $\mathcal Z_N$.] The sum over nonzero zero orbits of the seed $d_N(w)=\sum_{n\ge0}(-1)^nw^n/(Nn)!$, with the stated branch and multiplicity conventions.
\item[Polyphase section.] A decomposition of a coefficient sequence by residue class, with carry terms required when finite taps cross class boundaries.
\item[Resolution certificate.] An explicit bound or rational witness quantifying which structural distinction is detectable at a finite radius, precision, or polynomial degree.
\end{description}

The same word \emph{packet} can describe different analytic, arithmetic, and spectral representations. Its meaning is fixed by the ambient objects and maps in each construction, not by a claim that these families are automatically equivalent.

\MonographSubpart{Bibliography and prior work}
\pageheading{SOURCE MAP}{Primary references and bounded comparisons}
\small
\textbf{M\'endez, Miguel A.} \textit{Combinatorial differential operators in: Fa\`a di Bruno formula, enumeration of ballot paths, enriched rooted trees and increasing rooted trees},\;
\textit{Australasian Journal of Combinatorics} 89 (2024), 350--384.
Example 2.4, printed p.356, contains $f'=1+f(qz)^2$.
\url{https://ajc.maths.uq.edu.au/pdf/89/ajc_v89_p350.pdf}

\textbf{Dascaliuc, Pham, Thomann, Waymire.} Primary treatment of the
$\alpha$-Riccati equation $u'=-u+u(\alpha t)^2$, an adjacent feedback model.
Its equation is not asserted identical to the asymmetric model here.
\url{https://arxiv.org/html/2411.00629v2}

\textbf{Coffey, Mark W.} \textit{On the coefficients of the B\'aez-Duarte
criterion for the Riemann hypothesis and their extensions} (2006).
Signed M\"obius kernels and Newton methods are close antecedents. Their
signed weight differs from the positive squarefree weight used here.
\url{https://arxiv.org/pdf/math-ph/0608050}

\textbf{NIST DLMF.} \S16.8, equation 16.8.3, and \S16.21, equation 16.21.1,
give the differential equations used for the limited D-finite comparison.
\url{https://dlmf.nist.gov/16.8}\quad
\url{https://dlmf.nist.gov/16.21}

\textbf{NIST DLMF.} \S18.2 gives general orthogonal-polynomial theory;
\S10.46, equation 10.46.3, defines the Mittag--Leffler seed.
\url{https://dlmf.nist.gov/18.2}\quad
\url{https://dlmf.nist.gov/10.46}

\textbf{Uvarov, V. B.} \textit{The connection between systems of polynomials
that are orthogonal with respect to different distribution functions} (1969).
Original metadata was verified; the mass-modification formula is also proved
locally in the route file.
\url{https://doi.org/10.1016/0041-5553(69)90124-4}

\textbf{Elkies, Noam D.} Harvard Math 259 notes on finite-order entire
functions and Hadamard products. The genus-zero theorem supports the orbit
factorization, with multiplicities retained.
\url{https://people.math.harvard.edu/~elkies/M259.02/prod.pdf}

\textbf{Byrnes, Jiu, Moll, Vignat.} \textit{Recursion Rules for the
Hypergeometric Zeta Function} (2013). Hypergeometric zero spectra and
coefficient recursions provide adjacent prior machinery.
\url{https://arxiv.org/pdf/1305.1892}

\normalsize
\textbf{Voros, Andr\'e.} \textit{Zeta functions for the Riemann zeros},
Annales de l'Institut Fourier 53(3) (2003), 665--699. Sections 2.1--2.2
describe canonical products, Mellin continuation, and regularization for
general sequences. The present orbit proof verifies its own hypotheses.
\url{https://www.numdam.org/article/AIF_2003__53_3_665_0.pdf}

\textbf{OEIS.} A000182 records the tangent-number anchor. Candidate sequence
exports retain parameter, normalization, and regime labels; no OEIS submission
or unmatched-prefix novelty claim is made.
\url{https://oeis.org/A000182}

This literature comparison is bounded rather than exhaustive.
Full query lists and source classifications are in the route files. An
attempted EMIS retrieval of the original B\'aez-Duarte paper failed; no theorem
here depends on that failed retrieval. The absence of a located antecedent is not evidence of priority.

The integer arithmetic seed and its real-depth Gamma--Laguerre comparisons
are inherited project material. Their reuse is identified rather than counted
again as a discovery.

\pageheading{REFERENCES / RELATED SYSTEMS}{Where the coupled construction sits}
The entire-order and type results are proved directly in this edition.
The surrounding mechanisms have classical antecedents, whose precise
relationship is recorded in the included literature comparison.

\textbf{Miguel A. M\'endez, 2024.} \textit{Combinatorial differential operators
in: Fa\`a di Bruno formula, enumeration of ballot paths, enriched rooted trees
and increasing rooted trees}, Australasian Journal of Combinatorics 89(3),
350--384. Theorem 2.1 and Example 2.4 give the scalar diagonal equation.
\url{https://ajc.maths.uq.edu.au/pdf/89/ajc_v89_p350.pdf}

\textbf{H. Alrabaiah, I. Ahmad, K. Shah and G. Ur Rahman, 2020.}
\textit{Qualitative analysis of nonlinear coupled pantograph differential
equations of fractional order with integral boundary conditions}, Boundary
Value Problems, article 138. Its coupled fractional boundary-value system is
an adjacent construction with different hypotheses and conclusions.
\url{https://doi.org/10.1186/s13661-020-01432-2}

\textbf{Pawel Blasiak and Philippe Flajolet, 2011.}
\textit{Combinatorial Models of Creation--Annihilation}, S\'eminaire
Lotharingien de Combinatoire 65, B65c. Proposition 24 and Note 25 describe
multitype increasing trees and differential systems.
\url{https://arxiv.org/abs/1010.0354v3}

\textbf{R. Daniel Mauldin and S. C. Williams, 1988.}
\textit{Hausdorff Dimension in Graph Directed Constructions}, Transactions
of the American Mathematical Society 309(2), 811--829. Theorems 2 and 3 use
a spectral-radius equation for a geometric dimension. Here the same form
of spectral balance determines an analytic growth order through a separate
proof.
\url{https://sites.math.unt.edu/~mauldin/papers/no67.pdf}

\textbf{Radu Dascaliuc, Tuan N. Pham, Enrique Thomann and Edward C. Waymire.}
\textit{Transformation of Stochastic Recursions and Critical Phenomena in the
Analysis of the Aldous-Shields-Athreya Cascade and Related Mean Flow Equations},
arXiv:2411.00629v2. Equation (1.8) is an adjacent nonlinear Riccati delay model.
\url{https://arxiv.org/abs/2411.00629v2}

\textbf{Charles Knessl and Wojciech Szpankowski, 2006.}
\textit{Binary Trees, Left and Right Paths, WKB Expansions, and Painlev\'e
Transcendents}, ANALCO, 198--204. Its ordinary generating function satisfies
$C(z)=1+zC(wz)C(vz)$, a related dilation-product identity with different
normalization and coefficient rules.
\url{https://www.cs.purdue.edu/homes/spa/papers/analco06.pdf}

This is a bounded comparison rather than an exhaustive historical survey. The complete record
contains exact locators, inspected statements, terminology variants and search
limits. It leaves priority unresolved for the particular results of this edition.

\end{document}
